%% GENERATED by generate_papers.py -- edit the bank_*.py files, not this file.
\documentclass[11pt,a4paper]{article}
\usepackage[a4paper, top=20mm, bottom=20mm, left=18mm, right=18mm, headheight=14pt]{geometry}
\usepackage[T1]{fontenc}
\usepackage{lmodern}
\usepackage{amsmath, amssymb, mathtools}
\usepackage{braket}
\usepackage{array, tabularx, booktabs, multicol}
\usepackage{enumitem}
\usepackage[dvipsnames]{xcolor}
\usepackage[most]{tcolorbox}
\usepackage{needspace}
\usepackage{fancyhdr, lastpage}
\usepackage{pifont}
\usepackage{tikz}
\usepackage[hidelinks]{hyperref}
\newcolumntype{L}{>{\raggedright\arraybackslash}X}
\definecolor{brand}{HTML}{1F4E79}
\definecolor{okgreen}{HTML}{1B7F3B}
\setlength{\parindent}{0pt}
\setlength{\parskip}{3pt}
\newcommand{\optlabel}[1]{\textbf{(#1)}}
\newcommand{\correctlabel}[1]{\textcolor{okgreen}{\textbf{(#1)}\,\ding{51}}}
\newenvironment{qitem}[1]{\begin{list}{}{\setlength{\leftmargin}{3.3em}\setlength{\labelwidth}{3.1em}%
  \setlength{\labelsep}{0.2em}\setlength{\topsep}{3pt}\setlength{\parsep}{3pt}}\item[\textbf{Q.#1}\hfill]}{\end{list}}
\newcommand{\qsep}{\par\vspace{2pt}{\color{black!15}\hrule}\vspace{4pt}}
\newcommand{\qmeta}[2]{{\footnotesize\color{black!60}\textsc{#1} \;$\cdot$\; #2}\par}
\newtcolorbox{passagebox}[1]{enhanced, breakable, colback=brand!4, colframe=brand!70, fonttitle=\bfseries,
  title={#1}, left=6pt, right=6pt, top=4pt, bottom=4pt}
\newtcolorbox{sectionbanner}{enhanced, colback=brand, colframe=brand, coltext=white, fontupper=\large\bfseries,
  left=8pt, right=8pt, top=6pt, bottom=6pt, sharp corners}
\newtcolorbox{infobox}[1]{enhanced, breakable, colback=black!3, colframe=black!45, fonttitle=\bfseries,
  title={#1}, left=6pt, right=6pt, top=4pt, bottom=4pt}
\newtcolorbox{solbox}[1]{enhanced, breakable, colback=okgreen!3, colframe=okgreen!60!black, fonttitle=\bfseries,
  title={Answer: (#1) \quad Detailed solution}, left=6pt, right=6pt, top=3pt, bottom=4pt, before skip=4pt, after skip=6pt}
\newcommand{\HHLop}{\mathrm{HHL}}
\newcommand{\bub}[1]{\tikz[baseline=-0.65ex]\node[draw, circle, inner sep=0pt, minimum size=4.6mm, font=\scriptsize]{#1};}
\allowdisplaybreaks

\begin{document}
\pagestyle{fancy}\fancyhf{}
\fancyhead[L]{\small\color{black!70}The HHL Algorithm --- Assessment}
\fancyhead[R]{\small\color{black!70}\textbf{Answer Book}}
\fancyfoot[C]{\small Page \thepage\ of \pageref*{LastPage}}
\renewcommand{\headrulewidth}{0.4pt}

\thispagestyle{empty}
\begin{center}
{\color{brand}\rule{\linewidth}{1.2pt}}\\[6pt]
{\Huge\bfseries\color{brand} The HHL Algorithm}\\[4pt]
{\Large Assessment of the Harrow--Hassidim--Lloyd algorithm for $A\vec x = \vec b$}\\[4pt]
{\large\bfseries ANSWER BOOK}\\[2pt]
{\color{brand}\rule{\linewidth}{1.2pt}}
\end{center}
This book accompanies the separate \emph{Question Paper} (80 questions, 3 hours, 240 marks). It contains
\begin{enumerate}[leftmargin=1.4em, itemsep=0pt, topsep=2pt]
\item the \textbf{answer key}, section by section, with the topic tested by every question (page~\pageref{key}),
\item a \textbf{topic-wise map} grouping the questions by theme and a \textbf{quick-marking grid} (page~\pageref{topicmap}), and
\item the \textbf{Answer Bank}: every question restated with the correct option marked \textcolor{okgreen}{\ding{51}} and a
      complete solution with derivations (from page~\pageref{bank}).
\end{enumerate}
\begin{infobox}{Marking scheme (every question)}
Each question has \textbf{four} options (A), (B), (C) and (D). \textbf{ONLY ONE} of these four options is correct.\\[2pt]
\begin{tabular}{@{}ll@{}}
\emph{Full marks}: & $+3$ if ONLY the correct option is chosen;\\
\emph{Zero marks}: & $\phantom{+}0$ if none of the options is chosen (question unanswered);\\
\emph{Negative marks}: & $-1$ in all other cases (wrong option, or more than one option chosen).
\end{tabular}
\end{infobox}

\begin{infobox}{How to evaluate}
For each section count $N_{\mathrm c}$ = number of correct answers and $N_{\mathrm w}$ = number of questions answered wrongly (or with more than
one bubble darkened). Section marks $= 3N_{\mathrm c} - N_{\mathrm w}$; unanswered questions score $0$. The total is out of $240$.
Every numerical answer was re-computed independently (NumPy linear algebra, NumPy state vectors of the HHL circuits in the notebook's convention, and CUDA-Q for the notebook's HHL kernel and the phase-estimation circuits) by
\texttt{verify\_answers.py} in the source folder.
\end{infobox}
\begin{infobox}{Conventions (identical to the Question Paper)}

\begin{itemize}[leftmargin=1.2em, itemsep=1pt, topsep=2pt]
\item Linear system $A\vec x = \vec b$, $A$ a $d\times d$ matrix, $d = 2^{n_b}$. $A$ Hermitian with eigenpairs $A\ket{u_j} = \lambda_j\ket{u_j}$;
      $\ket b = \vec b/\lVert\vec b\rVert = \sum_j\beta_j\ket{u_j}$; $\ket x = A^{-1}\ket b/\lVert A^{-1}\ket b\rVert$.
      Condition number $\kappa = |\lambda|_{\max}/|\lambda|_{\min}$ (Hermitian $A$), in general $\kappa = \sigma_{\max}/\sigma_{\min}$.
      Sparsity $s$ = largest number of non-zero entries in a row.
\item Basis states are written with \textbf{qubit 0 first = most significant bit}: $\ket{q_0q_1\dots}$, e.g.\ $\ket{10} = \ket{2}$.
      $H = \frac{1}{\sqrt2}\begin{psmallmatrix}1&1\\1&-1\end{psmallmatrix}$; $X, Y, Z$ are the Pauli matrices;
      $R_Y(\theta) = \begin{psmallmatrix}\cos\frac\theta2&-\sin\frac\theta2\\ \sin\frac\theta2&\cos\frac\theta2\end{psmallmatrix}$;
      $\ket\pm = \frac{1}{\sqrt2}(\ket0\pm\ket1)$. $\mathrm{QFT}\ket{x} = \frac{1}{\sqrt N}\sum_{y}e^{2\pi i\,xy/N}\ket y$, $N = 2^n$.
\item \textbf{Phase estimation (QPE):} $n_c$ clock qubits in $\ket+$ control $U^{2^j}$, then the inverse QFT; an eigenphase $U\ket u = e^{2\pi i\varphi}\ket u$
      gives the outcome $k \approx \varphi\,2^{n_c}$.
\item \textbf{HHL (as in \texttt{HHL.ipynb}):} registers clock ($n_c$), system ($n_b$), ancilla (1), total $n_c + n_b + 1$ qubits.
      $U = e^{iAt}$, so eigenvalue $\lambda$ gives clock value $k = \lambda t\,2^{n_c}/(2\pi)$. $C = 2\pi/(2^{n_c}t)$ is the smallest representable
      eigenvalue; for clock value $k$ ($\tilde\lambda_k = kC$) the ancilla is rotated by $R_Y(\theta_k)$, $\theta_k = 2\arcsin(C/\tilde\lambda_k)$.
      Clock values $k \ge 2^{n_c-1}$ mean the negative value $k - 2^{n_c}$ (two's complement) unless a question says otherwise.
      The run succeeds when the ancilla is measured as $1$ (and the clock as $0\dots0$); $P_{\text{success}} = \sum_j|\beta_j|^2C^2/\lambda_j^2$ when every $\lambda_j$ is exactly representable.
\item \textbf{Notebook default:} rows of $A$: $(2.5, -0.5, -1, 0)$, $(-0.5, 2.5, 0, -1)$, $(-1, 0, 2.5, -0.5)$, $(0, -1, -0.5, 2.5)$ (eigenvalues $1, 2, 3, 4$),
      $\vec b = (1, 0, 0, 0)$, $n_c = 4$, $t = \pi/8$, $C = 1$, 7 qubits, 20\,000 shots. A \emph{previously recorded run} means the outputs stored in the
      notebook (CUDA-Q 0.16.0, CPU target \texttt{qpp-cpu}).
\item $\log$ means $\log_2$; $\tilde O$ hides polylogarithmic factors. Numerical answers are rounded.
\end{itemize}
\end{infobox}
\newpage
\section*{Answer key}\label{key}
{\small Answer distribution: (A) 20, (B) 20, (C) 20, (D) 20.}

\needspace{14\baselineskip}
\subsection*{Section A --- Linear Algebra and Mathematical Foundations}
{\small\renewcommand{\arraystretch}{1.15}
\begin{tabularx}{\linewidth}{@{}r c L @{\hspace{1em}} r c L@{}}
\toprule
\textbf{Q.} & \textbf{Ans.} & \textbf{Topic tested} & \textbf{Q.} & \textbf{Ans.} & \textbf{Topic tested}\\ \midrule
1 & \textbf{D} & Solving a $2\times2$ system & 11 & \textbf{B} & Normalisation and amplitude encoding \\
2 & \textbf{C} & Eigenvalues and $\kappa$ of a symmetric matrix & 12 & \textbf{C} & Solving $A\vec x=\vec b$ in the eigenbasis \\
3 & \textbf{B} & Matrix properties (Hermitian, unitary, singular) & 13 & \textbf{D} & Norms, $\kappa$ and determinant \\
4 & \textbf{D} & Condition number of a diagonal matrix & 14 & \textbf{D} & Hermitian embedding of a non-Hermitian matrix \\
5 & \textbf{C} & Eigenvalues of a complex Hermitian matrix & 15 & \textbf{C} & When is $e^{iAt}$ unitary? \\
6 & \textbf{A} & Inverse via spectral decomposition & 16 & \textbf{A} & Error amplification by $\kappa$ \\
7 & \textbf{B} & Spectral theorem vs.\ real eigenvalues & 17 & \textbf{D} & $2\times2$ example: eigenvalues \\
8 & \textbf{A} & Condition number of $U$ and of $A^{-1}$ & 18 & \textbf{C} & $2\times2$ example: normalised solution \\
9 & \textbf{B} & Sparsity of the 1-D Laplacian & 19 & \textbf{B} & Embedding + padding: qubit count \\
10 & \textbf{A} & Spectrum and $\kappa$ of the Laplacian & 20 & \textbf{A} & $\kappa$ from singular values \\
\bottomrule
\end{tabularx}}

\needspace{14\baselineskip}
\subsection*{Section B --- Quantum Computing Foundations}
{\small\renewcommand{\arraystretch}{1.15}
\begin{tabularx}{\linewidth}{@{}r c L @{\hspace{1em}} r c L@{}}
\toprule
\textbf{Q.} & \textbf{Ans.} & \textbf{Topic tested} & \textbf{Q.} & \textbf{Ans.} & \textbf{Topic tested}\\ \midrule
21 & \textbf{C} & Born rule & 31 & \textbf{B} & Global vs.\ relative phase \\
22 & \textbf{A} & State-vector size and memory & 32 & \textbf{A} & Controlled-$R_Y$ and ancilla probability \\
23 & \textbf{C} & Gate identity $HZH = X$ & 33 & \textbf{C} & Selecting a clock value with $X$ gates \\
24 & \textbf{B} & $R_Y$ rotations of $\ket0$ & 34 & \textbf{A} & Post-selection on an ancilla \\
25 & \textbf{C} & Entanglement test & 35 & \textbf{D} & Inverse of a circuit \\
26 & \textbf{D} & Phase kickback & 36 & \textbf{D} & Gate counts of a 4-qubit QFT \\
27 & \textbf{D} & Inverse QFT of a phase state & 37 & \textbf{B} & Uniform superposition \\
28 & \textbf{B} & QPE with an exact 3-bit phase & 38 & \textbf{D} & Product states, reduced states \\
29 & \textbf{A} & QPE with a non-representable phase & 39 & \textbf{B} & QPE trace: phase state \\
30 & \textbf{C} & Unitarity vs.\ measurement & 40 & \textbf{A} & QPE trace: outcome \\
\bottomrule
\end{tabularx}}

\needspace{14\baselineskip}
\subsection*{Section C --- HHL Algorithm Fundamentals}
{\small\renewcommand{\arraystretch}{1.15}
\begin{tabularx}{\linewidth}{@{}r c L @{\hspace{1em}} r c L@{}}
\toprule
\textbf{Q.} & \textbf{Ans.} & \textbf{Topic tested} & \textbf{Q.} & \textbf{Ans.} & \textbf{Topic tested}\\ \midrule
41 & \textbf{C} & What HHL outputs & 51 & \textbf{C} & Negative eigenvalues (two's complement) \\
42 & \textbf{D} & State after phase estimation & 52 & \textbf{A} & Code: $X$ gates of the inversion \\
43 & \textbf{C} & Eigenvalue $\to$ clock bit string & 53 & \textbf{B} & Norm recovery $\vec x = \lVert\vec b\rVert\,\vec a/C$ \\
44 & \textbf{B} & Constant $C$ and rotation angle & 54 & \textbf{A} & Default $4\times4$ solution \\
45 & \textbf{D} & Ancilla amplitude $C/\lambda$ & 55 & \textbf{D} & Algorithm step $\leftrightarrow$ CUDA-Q code \\
46 & \textbf{A} & Roles of the registers & 56 & \textbf{B} & Preparing $\ket b$ \\
47 & \textbf{D} & Success probability & 57 & \textbf{B} & Role of the inverse QFT \\
48 & \textbf{B} & Qubit count $n_b + n_c + 1$ & 58 & \textbf{A} & $\ket b$ an eigenvector \\
49 & \textbf{C} & Why uncompute? & 59 & \textbf{D} & $2\times2$ HHL trace: after QPE \\
50 & \textbf{A} & Choice of $C$ & 60 & \textbf{C} & $2\times2$ HHL trace: output \\
\bottomrule
\end{tabularx}}

\needspace{14\baselineskip}
\subsection*{Section D --- Advanced HHL and Implementation}
{\small\renewcommand{\arraystretch}{1.15}
\begin{tabularx}{\linewidth}{@{}r c L @{\hspace{1em}} r c L@{}}
\toprule
\textbf{Q.} & \textbf{Ans.} & \textbf{Topic tested} & \textbf{Q.} & \textbf{Ans.} & \textbf{Topic tested}\\ \midrule
61 & \textbf{C} & Complexity of HHL & 71 & \textbf{C} & Recorded $P_{\text{success}}$: shot statistics \\
62 & \textbf{B} & Repetitions and amplitude amplification & 72 & \textbf{A} & Recorded read-out of $|x_0|^2$ \\
63 & \textbf{C} & Quantum linear-system algorithms & 73 & \textbf{B} & Circuit resources vs.\ $n_c$ \\
64 & \textbf{B} & Read-out cost of all $x_i$ & 74 & \textbf{D} & Error sources and their effects \\
65 & \textbf{D} & ``Exponential speed-up'' claims & 75 & \textbf{B} & Useful outputs of HHL \\
66 & \textbf{A} & Too few clock qubits: rescaled $t$ & 76 & \textbf{B} & Dequantisation \\
67 & \textbf{D} & Clock qubits for precision $\epsilon$ & 77 & \textbf{A} & $\kappa^2$ run-time dependence \\
68 & \textbf{D} & Memory of a large HHL simulation & 78 & \textbf{D} & Solution inside the embedded system \\
69 & \textbf{A} & Trotter steps for $e^{iAt}$ & 79 & \textbf{A} & Recorded run: $P_{\text{success}}$ exactly \\
70 & \textbf{C} & What the notebook implements & 80 & \textbf{C} & Recorded run: $\lVert\vec x\rVert$ from $P_{\text{success}}$ \\
\bottomrule
\end{tabularx}}

\needspace{18\baselineskip}
\subsection*{Topic-wise map of the questions}\label{topicmap}
{\small Every question appears under at least one theme; several appear under two.\par}
{\small\renewcommand{\arraystretch}{1.2}
\begin{tabularx}{\linewidth}{@{}>{\raggedright\arraybackslash}p{0.42\linewidth}L@{}}
\toprule
\textbf{Theme} & \textbf{Question numbers}\\ \midrule
Linear systems and their classical solution & 1, 6, 12, 16, 18, 54 \\
Eigenvalues, Hermitian matrices and spectral decomposition & 2, 3, 5, 6, 7, 10, 12, 17 \\
Condition number, norms and sparsity & 4, 8, 9, 10, 13, 16, 20, 77 \\
Hermitian embedding, padding and unitarity of $e^{iAt}$ & 14, 15, 19, 78 \\
Normalisation, amplitude encoding and state preparation & 11, 22, 37, 56 \\
Qubits, measurement, phases and entanglement & 21, 23, 25, 30, 31, 38 \\
Gates, controlled operations and circuits & 24, 26, 32, 33, 35, 36, 52 \\
QFT and quantum phase estimation & 27, 28, 29, 36, 39, 40, 42, 43, 57, 66, 67 \\
Eigenvalue inversion, the constant $C$ and the ancilla & 44, 45, 50, 51, 52, 59, 60 \\
Post-selection, success probability and uncomputation & 34, 47, 49, 53, 58, 60, 62 \\
HHL registers, circuit structure and resources & 46, 48, 55, 68, 73 \\
Complexity, read-out and advanced algorithms & 41, 61, 63, 64, 65, 69, 75, 76, 77 \\
Noise, errors and limitations & 66, 70, 74, 76 \\
Interpreting recorded results of \texttt{HHL.ipynb} & 54, 70, 71, 72, 79, 80 \\
\bottomrule
\end{tabularx}}

\subsection*{Quick-marking grid}
\begin{center}\renewcommand{\arraystretch}{1.3}
\begin{tabular}{|c|c|c|c|c|c|c|c|c|c|}\hline
1\,\textbf{D} & 2\,\textbf{C} & 3\,\textbf{B} & 4\,\textbf{D} & 5\,\textbf{C} & 6\,\textbf{A} & 7\,\textbf{B} & 8\,\textbf{A} & 9\,\textbf{B} & 10\,\textbf{A} \\
\hline
11\,\textbf{B} & 12\,\textbf{C} & 13\,\textbf{D} & 14\,\textbf{D} & 15\,\textbf{C} & 16\,\textbf{A} & 17\,\textbf{D} & 18\,\textbf{C} & 19\,\textbf{B} & 20\,\textbf{A} \\
\hline
21\,\textbf{C} & 22\,\textbf{A} & 23\,\textbf{C} & 24\,\textbf{B} & 25\,\textbf{C} & 26\,\textbf{D} & 27\,\textbf{D} & 28\,\textbf{B} & 29\,\textbf{A} & 30\,\textbf{C} \\
\hline
31\,\textbf{B} & 32\,\textbf{A} & 33\,\textbf{C} & 34\,\textbf{A} & 35\,\textbf{D} & 36\,\textbf{D} & 37\,\textbf{B} & 38\,\textbf{D} & 39\,\textbf{B} & 40\,\textbf{A} \\
\hline
41\,\textbf{C} & 42\,\textbf{D} & 43\,\textbf{C} & 44\,\textbf{B} & 45\,\textbf{D} & 46\,\textbf{A} & 47\,\textbf{D} & 48\,\textbf{B} & 49\,\textbf{C} & 50\,\textbf{A} \\
\hline
51\,\textbf{C} & 52\,\textbf{A} & 53\,\textbf{B} & 54\,\textbf{A} & 55\,\textbf{D} & 56\,\textbf{B} & 57\,\textbf{B} & 58\,\textbf{A} & 59\,\textbf{D} & 60\,\textbf{C} \\
\hline
61\,\textbf{C} & 62\,\textbf{B} & 63\,\textbf{C} & 64\,\textbf{B} & 65\,\textbf{D} & 66\,\textbf{A} & 67\,\textbf{D} & 68\,\textbf{D} & 69\,\textbf{A} & 70\,\textbf{C} \\
\hline
71\,\textbf{C} & 72\,\textbf{A} & 73\,\textbf{B} & 74\,\textbf{D} & 75\,\textbf{B} & 76\,\textbf{B} & 77\,\textbf{A} & 78\,\textbf{D} & 79\,\textbf{A} & 80\,\textbf{C} \\
\hline\end{tabular}\end{center}
\newpage
\section*{Answer Bank --- questions with detailed solutions}\label{bank}

\begin{sectionbanner}SECTION A \enspace$\cdot$\enspace Linear Algebra and Mathematical Foundations \hfill Q.1--20\end{sectionbanner}

\needspace{13\baselineskip}
\begin{qitem}{1}
\qmeta{Single correct}{Solving a $2\times2$ system}
The solution of the linear system $\begin{pmatrix}2&1\\1&3\end{pmatrix}\begin{pmatrix}x_1\\x_2\end{pmatrix} = \begin{pmatrix}1\\2\end{pmatrix}$ is
\par\smallskip\noindent\begin{tabularx}{\linewidth}{@{}*{2}{L}@{}}
\optlabel{A}\,$(x_1, x_2) = (-0.2,\ 0.6)$ & \optlabel{B}\,$(x_1, x_2) = (0.6,\ 0.2)$ \\[3pt]
\optlabel{C}\,$(x_1, x_2) = (1,\ 0.\overline{3})$ & \correctlabel{D}\,$(x_1, x_2) = (0.2,\ 0.6)$
\end{tabularx}
\end{qitem}
\begin{solbox}{D}
\textbf{\color{brand}Correct answer.} \textcolor{okgreen}{\textbf{(D)}} $(x_1, x_2) = (0.2,\ 0.6)$

\textbf{\color{brand}Concepts required.} A square system $A\vec x = \vec b$ has a unique solution iff $\det A \neq 0$; for $2\times2$ matrices the inverse has a closed form.

\textbf{\color{brand}Relevant formulae.} $\begin{psmallmatrix}a&b\\c&d\end{psmallmatrix}^{-1} = \frac{1}{ad - bc}\begin{psmallmatrix}d&-b\\-c&a\end{psmallmatrix}$; $\ \vec x = A^{-1}\vec b$.

\textbf{\color{brand}Step-by-step explanation.} $\det A = 2\cdot3 - 1\cdot1 = 5 \neq 0$, so the system has a unique solution. For a $2\times2$ matrix
\[ A^{-1} = \frac{1}{\det A}\begin{pmatrix}3&-1\\-1&2\end{pmatrix}, \qquad
   \vec x = A^{-1}\vec b = \frac15\begin{pmatrix}3\cdot1 - 1\cdot2\\ -1\cdot1 + 2\cdot2\end{pmatrix} = \frac15\begin{pmatrix}1\\3\end{pmatrix} = \begin{pmatrix}0.2\\0.6\end{pmatrix}. \]
\emph{Check:} $2(0.2) + 0.6 = 1$ and $0.2 + 3(0.6) = 2$. (Gaussian elimination gives the same: $R_2 \leftarrow R_2 - \tfrac12R_1$ gives $2.5x_2 = 1.5$.)

\textbf{\color{brand}Why the other options are incorrect.}
\begin{itemize}[leftmargin=2.4em, itemsep=1pt, topsep=2pt]
\item[\textbf{(A)}] $(-0.2, 0.6)$ has a sign error in the first component: $2(-0.2) + 0.6 = 0.2 \neq 1$.
\item[\textbf{(B)}] $(0.6, 0.2)$ swaps the two components; substituting gives $2(0.6) + 0.2 = 1.4 \neq 1$.
\item[\textbf{(C)}] $(1, 0.\overline3)$ solves each equation separately ($2x_1 = 2$?) and ignores the coupling; $2(1) + \tfrac13 \neq 1$.
\end{itemize}
\textbf{\color{brand}Key takeaway.} Always substitute the candidate back into $A\vec x = \vec b$: it takes ten seconds and catches every slip.
\end{solbox}
\qsep

\needspace{13\baselineskip}
\begin{qitem}{2}
\qmeta{Single correct}{Eigenvalues and $\kappa$ of a symmetric matrix}
The eigenvalues and the condition number $\kappa$ of $A = \begin{pmatrix}5&2\\2&2\end{pmatrix}$ are
\par\smallskip\noindent\begin{tabularx}{\linewidth}{@{}*{2}{L}@{}}
\optlabel{A}\,$3$ and $4$; $\ \kappa = 4/3$ & \optlabel{B}\,$2$ and $5$; $\ \kappa = 2.5$ \\[3pt]
\correctlabel{C}\,$1$ and $6$; $\ \kappa = 6$ & \optlabel{D}\,$1$ and $6$; $\ \kappa = 7$
\end{tabularx}
\end{qitem}
\begin{solbox}{C}
\textbf{\color{brand}Correct answer.} \textcolor{okgreen}{\textbf{(C)}} $1$ and $6$; $\ \kappa = 6$

\textbf{\color{brand}Concepts required.} Eigenvalues are the roots of $\det(A - \lambda I) = 0$; for a symmetric matrix they are real and $\kappa = |\lambda|_{\max}/|\lambda|_{\min}$.

\textbf{\color{brand}Relevant formulae.} $2\times2$: $\lambda^2 - (\operatorname{tr}A)\lambda + \det A = 0$; $\ \kappa = |\lambda|_{\max}/|\lambda|_{\min}$.

\textbf{\color{brand}Step-by-step explanation.} Characteristic polynomial: $\lambda^2 - (\operatorname{tr}A)\lambda + \det A = \lambda^2 - 7\lambda + 6 = (\lambda - 1)(\lambda - 6)$, so
$\lambda \in \{1, 6\}$. Eigenvectors: $\lambda = 1$: $4x + 2y = 0 \Rightarrow (1, -2)$; $\lambda = 6$: $-x + 2y = 0 \Rightarrow (2, 1)$ --- orthogonal,
as they must be for a symmetric matrix. $A$ is symmetric positive definite, so
\[ \kappa = \frac{|\lambda|_{\max}}{|\lambda|_{\min}} = \frac61 = 6. \]
(The diagonal entries $5, 2$ are \emph{not} the eigenvalues, because $A$ is not diagonal.)

\textbf{\color{brand}Why the other options are incorrect.}
\begin{itemize}[leftmargin=2.4em, itemsep=1pt, topsep=2pt]
\item[\textbf{(A)}] $3$ and $4$ have the right sum $7$ but product $12 \neq 6$, so they are not the roots.
\item[\textbf{(B)}] $2$ and $5$ are the diagonal entries, which are eigenvalues only for diagonal (or triangular) matrices; their product $10 \neq \det A = 6$.
\item[\textbf{(D)}] $\kappa = 7$ adds the eigenvalues ($\operatorname{tr}A$) instead of dividing the largest by the smallest.
\end{itemize}
\textbf{\color{brand}Key takeaway.} Check eigenvalues with two invariants: their sum is the trace and their product the determinant.
\end{solbox}
\qsep

\needspace{13\baselineskip}
\begin{qitem}{3}
\qmeta{Matching lists}{Matrix properties (Hermitian, unitary, singular)}
Match each matrix (List-I) with the description that fits it (List-II).
\par\smallskip\noindent{\renewcommand{\arraystretch}{1.35}%
\begin{tabularx}{\linewidth}{@{}c>{\raggedright\arraybackslash}p{0.33\linewidth}@{\hspace{1.2em}}cL@{}}
\toprule
\multicolumn{2}{@{}l}{\textbf{List-I}} & \multicolumn{2}{l@{}}{\textbf{List-II}} \\ \midrule
\textbf{(P)} & $\begin{pmatrix}0&1\\1&0\end{pmatrix}$ & \textbf{(1)} & Hermitian and unitary \\
\textbf{(Q)} & $\begin{pmatrix}1&2\\2&4\end{pmatrix}$ & \textbf{(2)} & Hermitian, invertible, not unitary \\
\textbf{(R)} & $\begin{pmatrix}1&i\\-i&2\end{pmatrix}$ & \textbf{(3)} & Hermitian but singular \\
\textbf{(S)} & $\begin{pmatrix}0&1\\0&0\end{pmatrix}$ & \textbf{(4)} & not diagonalisable \\
\bottomrule
\end{tabularx}}
\par\smallskip\noindent\begin{tabularx}{\linewidth}{@{}*{2}{L}@{}}
\optlabel{A}\,P$\to$2, Q$\to$3, R$\to$1, S$\to$4 & \correctlabel{B}\,P$\to$1, Q$\to$3, R$\to$2, S$\to$4 \\[3pt]
\optlabel{C}\,P$\to$1, Q$\to$4, R$\to$2, S$\to$3 & \optlabel{D}\,P$\to$1, Q$\to$2, R$\to$3, S$\to$4
\end{tabularx}
\end{qitem}
\begin{solbox}{B}
\textbf{\color{brand}Correct answer.} \textcolor{okgreen}{\textbf{(B)}} P$\to$1, Q$\to$3, R$\to$2, S$\to$4

\textbf{\color{brand}Concepts required.} Hermitian: $A = A^\dagger$. Unitary: $A^\dagger A = I$. Singular: $\det A = 0$. Diagonalisable: a full set of eigenvectors.

\textbf{\color{brand}Relevant formulae.} $A^\dagger = \bar A^{T}$; $\ \det\begin{psmallmatrix}a&b\\c&d\end{psmallmatrix} = ad - bc$; Hermitian $\Rightarrow$ unitarily diagonalisable.

\textbf{\color{brand}Step-by-step explanation.} \begin{itemize}\itemsep0pt
\item P $= X$: real symmetric, so Hermitian; $X^\dagger X = X^2 = I$, so unitary (1).
\item Q: symmetric (Hermitian) but $\det = 4 - 4 = 0$ (second row $= 2\times$ first): singular, eigenvalues $0$ and $5$ (3).
\item R: $R^\dagger$ = conjugate transpose $= \begin{psmallmatrix}1&\overline{-i}\\ \overline{i}&2\end{psmallmatrix} = \begin{psmallmatrix}1&i\\-i&2\end{psmallmatrix} = R$, so Hermitian;
$\det R = 2 - (i)(-i) = 2 - 1 = 1 \neq 0$, invertible; $R^\dagger R = R^2 \neq I$ (e.g.\ $(R^2)_{11} = 1 + 1 = 2$), not unitary (2).
\item S: both eigenvalues are $0$ but $S \neq 0$; the eigenspace of $0$ is only 1-dimensional ($\{(a, 0)\}$), so $S$ is not diagonalisable (4).
\end{itemize}
Only Hermitian, invertible matrices (P and R) are directly usable in HHL.

\textbf{\color{brand}Why the other options are incorrect.}
\begin{itemize}[leftmargin=2.4em, itemsep=1pt, topsep=2pt]
\item[\textbf{(A)}] P$\to$2 calls $X$ ``not unitary'', but $X^2 = I$; R$\to$1 calls the $2\times2$ matrix with diagonal $1, 2$ unitary although its first column has norm $\sqrt2$.
\item[\textbf{(C)}] Q$\to$4: Q is symmetric, hence diagonalisable (eigenvalues $0$ and $5$); S$\to$3: S is not Hermitian ($S^\dagger \neq S$).
\item[\textbf{(D)}] Q$\to$2 calls Q invertible, but its rows are proportional ($\det = 0$); R$\to$3 calls R singular, but $\det R = 1$.
\end{itemize}
\textbf{\color{brand}Key takeaway.} HHL needs a Hermitian (or embedded) and invertible matrix; check $A = A^\dagger$ and $\det A \neq 0$ first.
\end{solbox}
\qsep

\needspace{13\baselineskip}
\begin{qitem}{4}
\qmeta{Single correct}{Condition number of a diagonal matrix}
The condition number of $A = \operatorname{diag}(1,\ 0.5,\ 0.01,\ 0.2)$ is
\par\smallskip\noindent\begin{tabularx}{\linewidth}{@{}*{4}{L}@{}}
\optlabel{A}\,$0.01$ & \optlabel{B}\,$1.71$ & \optlabel{C}\,$50$ & \correctlabel{D}\,$100$
\end{tabularx}
\end{qitem}
\begin{solbox}{D}
\textbf{\color{brand}Correct answer.} \textcolor{okgreen}{\textbf{(D)}} $100$

\textbf{\color{brand}Concepts required.} For a diagonal matrix the eigenvalues and singular values are the absolute values of the diagonal entries.

\textbf{\color{brand}Relevant formulae.} $\kappa = \lVert A\rVert_2\lVert A^{-1}\rVert_2 = \sigma_{\max}/\sigma_{\min}$; diagonal: $\sigma_i = |a_{ii}|$.

\textbf{\color{brand}Step-by-step explanation.} For a diagonal (hence Hermitian) matrix the eigenvalues are the diagonal entries: $1, 0.5, 0.01, 0.2$.
\[ \kappa = \frac{|\lambda|_{\max}}{|\lambda|_{\min}} = \frac{1}{0.01} = 100. \]
Equivalently $\kappa = \lVert A\rVert_2\,\lVert A^{-1}\rVert_2 = 1 \times 100$. A system with $\kappa = 100$ amplifies relative errors in $\vec b$
by up to a factor $100$; for HHL the run time grows with $\kappa$ (Section D).

\textbf{\color{brand}Why the other options are incorrect.}
\begin{itemize}[leftmargin=2.4em, itemsep=1pt, topsep=2pt]
\item[\textbf{(A)}] $0.01$ is $\lambda_{\min}/\lambda_{\max}$, the reciprocal; $\kappa \ge 1$ always.
\item[\textbf{(B)}] $1.71$ is the sum of the entries (the trace), which has nothing to do with conditioning.
\item[\textbf{(C)}] $50$ divides $0.5$ by $0.01$, using the second-largest instead of the largest entry.
\end{itemize}
\textbf{\color{brand}Key takeaway.} $\kappa = $ largest over smallest $|\lambda|$; $\kappa \ge 1$, and $\kappa = 1$ only for multiples of unitaries.
\end{solbox}
\qsep

\needspace{13\baselineskip}
\begin{qitem}{5}
\qmeta{Single correct}{Eigenvalues of a complex Hermitian matrix}
The Hermitian matrix $A = \begin{pmatrix}3 & 1-i\\ 1+i & 2\end{pmatrix}$ has eigenvalues
\par\smallskip\noindent\begin{tabularx}{\linewidth}{@{}*{4}{L}@{}}
\optlabel{A}\,$3$ and $2$ & \optlabel{B}\,$0$ and $5$ & \correctlabel{C}\,$1$ and $4$ & \optlabel{D}\,$2 \pm i$
\end{tabularx}
\end{qitem}
\begin{solbox}{C}
\textbf{\color{brand}Correct answer.} \textcolor{okgreen}{\textbf{(C)}} $1$ and $4$

\textbf{\color{brand}Concepts required.} Hermitian matrices may have complex entries but always have real eigenvalues; the off-diagonal pair contributes $-|a_{12}|^2$ to the determinant.

\textbf{\color{brand}Relevant formulae.} $\det\begin{psmallmatrix}a&z\\ \bar z&d\end{psmallmatrix} = ad - |z|^2$; $\ \lambda_{1,2} = \frac{\operatorname{tr}A \pm\sqrt{(\operatorname{tr}A)^2 - 4\det A}}{2}$.

\textbf{\color{brand}Step-by-step explanation.} $\operatorname{tr}A = 5$, $\det A = 3\cdot2 - (1-i)(1+i) = 6 - (1 + 1) = 4$. The eigenvalues solve $\lambda^2 - 5\lambda + 4 = 0$:
$\lambda = 1$ or $\lambda = 4$. Both are real, as for every Hermitian matrix, even though $A$ has complex entries. (Note $(1-i)(1+i) = |1+i|^2 = 2$.)
$A$ is positive definite with $\kappa = 4$.

\textbf{\color{brand}Why the other options are incorrect.}
\begin{itemize}[leftmargin=2.4em, itemsep=1pt, topsep=2pt]
\item[\textbf{(A)}] $3$ and $2$ are the diagonal entries; they ignore the coupling $1 \mp i$ (product $6 \neq \det A = 4$).
\item[\textbf{(B)}] $0$ and $5$ would need $\det A = 0$, i.e.\ $|1+i|^2 = 6$; in fact $|1+i|^2 = 2$.
\item[\textbf{(D)}] $2 \pm i$ are complex, which is impossible for a Hermitian matrix (their sum $4$ is not even $\operatorname{tr}A = 5$); the discriminant $25 - 16 = 9$ is positive.
\end{itemize}
\textbf{\color{brand}Key takeaway.} For a Hermitian $2\times2$ matrix: $\det = ad - |z|^2$, and the eigenvalues must come out real.
\end{solbox}
\qsep

\needspace{13\baselineskip}
\begin{qitem}{6}
\qmeta{Single correct}{Inverse via spectral decomposition}
$A = \begin{pmatrix}2&1\\1&2\end{pmatrix}$ has eigenpairs $\lambda_1 = 3$, $\ket{u_1} = \tfrac{1}{\sqrt2}(1, 1)$ and $\lambda_2 = 1$,
$\ket{u_2} = \tfrac{1}{\sqrt2}(1, -1)$. Using $A^{-1} = \sum_j \lambda_j^{-1}\ket{u_j}\!\bra{u_j}$, the off-diagonal entry $(A^{-1})_{12}$ is
\par\smallskip\noindent\begin{tabularx}{\linewidth}{@{}*{4}{L}@{}}
\correctlabel{A}\,$-\tfrac13$ & \optlabel{B}\,$+\tfrac13$ & \optlabel{C}\,$-\tfrac23$ & \optlabel{D}\,$\tfrac16$
\end{tabularx}
\end{qitem}
\begin{solbox}{A}
\textbf{\color{brand}Correct answer.} \textcolor{okgreen}{\textbf{(A)}} $-\tfrac13$

\textbf{\color{brand}Concepts required.} Spectral decomposition: a Hermitian matrix is a weighted sum of projectors onto its eigenvectors; functions of $A$ act on the weights.

\textbf{\color{brand}Relevant formulae.} $A = \sum_j\lambda_j\ket{u_j}\!\bra{u_j}$, $\ f(A) = \sum_j f(\lambda_j)\ket{u_j}\!\bra{u_j}$, so $A^{-1} = \sum_j\lambda_j^{-1}\ket{u_j}\!\bra{u_j}$.

\textbf{\color{brand}Step-by-step explanation.} $\ket{u_1}\!\bra{u_1} = \tfrac12\begin{psmallmatrix}1&1\\1&1\end{psmallmatrix}$ and $\ket{u_2}\!\bra{u_2} = \tfrac12\begin{psmallmatrix}1&-1\\-1&1\end{psmallmatrix}$. Hence
\[ A^{-1} = \frac13\cdot\frac12\begin{pmatrix}1&1\\1&1\end{pmatrix} + \frac11\cdot\frac12\begin{pmatrix}1&-1\\-1&1\end{pmatrix}
          = \begin{pmatrix}\tfrac16 + \tfrac12 & \tfrac16 - \tfrac12\\ \tfrac16 - \tfrac12 & \tfrac16 + \tfrac12\end{pmatrix}
          = \frac13\begin{pmatrix}2&-1\\-1&2\end{pmatrix}. \]
So $(A^{-1})_{12} = -\tfrac13$, in agreement with the $2\times2$ formula $\frac{1}{\det A}\begin{psmallmatrix}2&-1\\-1&2\end{psmallmatrix}$, $\det A = 3$.
This ``invert every eigenvalue'' formula is exactly what HHL implements on a quantum state.

\textbf{\color{brand}Why the other options are incorrect.}
\begin{itemize}[leftmargin=2.4em, itemsep=1pt, topsep=2pt]
\item[\textbf{(B)}] $+\tfrac13$ has the wrong sign: the $\lambda = 1$ projector contributes $-\frac12$ off the diagonal and dominates $+\frac16$.
\item[\textbf{(C)}] $-\tfrac23$ forgets the factor $\frac12$ in $\ket{u_j}\!\bra{u_j}$ for one of the terms ($\frac13 - 1$).
\item[\textbf{(D)}] $\tfrac16$ keeps only the $\lambda = 3$ term and forgets the $\lambda = 1$ projector.
\end{itemize}
\textbf{\color{brand}Key takeaway.} $f(A)$ = same eigenvectors, $f$ applied to the eigenvalues. HHL applies $f(\lambda) = 1/\lambda$ to a quantum state.
\end{solbox}
\qsep

\needspace{13\baselineskip}
\begin{qitem}{7}
\qmeta{Assertion--Reason}{Spectral theorem vs.\ real eigenvalues}
\textbf{Assertion (A):} Every Hermitian matrix can be diagonalised by a unitary matrix, $A = V\Lambda V^\dagger$.\\[1mm]
\textbf{Reason (R):} All eigenvalues of a Hermitian matrix are real.
\begin{itemize}[leftmargin=2.6em, labelsep=0.4em, itemsep=2pt, topsep=3pt]
\item[\optlabel{A}] Both Assertion (A) and Reason (R) are true, and R is the correct explanation of A
\item[\correctlabel{B}] Both A and R are true, but R is \emph{not} the correct explanation of A
\item[\optlabel{C}] A is true, but R is false
\item[\optlabel{D}] A is false, but R is true
\end{itemize}
\end{qitem}
\begin{solbox}{B}
\textbf{\color{brand}Correct answer.} \textcolor{okgreen}{\textbf{(B)}} Both A and R are true, but R is \emph{not} the correct explanation of A

\textbf{\color{brand}Concepts required.} Spectral theorem: normal matrices ($AA^\dagger = A^\dagger A$) are exactly the unitarily diagonalisable ones. Hermitian matrices are normal and have real eigenvalues.

\textbf{\color{brand}Relevant formulae.} $A = V\Lambda V^\dagger$, $V^\dagger V = I$; $\ \lambda\braket{u|u} = \braket{u|A|u} = \bar\lambda\braket{u|u}$ for Hermitian $A$.

\textbf{\color{brand}Step-by-step explanation.} \textbf{A is true:} this is the spectral theorem (it holds for every \emph{normal} matrix, $AA^\dagger = A^\dagger A$, and Hermitian matrices are normal).
\textbf{R is true:} $\lambda\braket{u|u} = \braket{u|A|u} = \braket{u|A^\dagger|u}^* = \bar\lambda\braket{u|u}$, so $\lambda = \bar\lambda$.
But R does not explain A: real eigenvalues do not imply diagonalisability --- $\begin{psmallmatrix}1&1\\0&1\end{psmallmatrix}$ has the real eigenvalue $1$ twice
and cannot be diagonalised, and unitary matrices are diagonalisable by a unitary although their eigenvalues are complex. The
explanation of A is normality. Answer: (B).

\textbf{\color{brand}Why the other options are incorrect.}
\begin{itemize}[leftmargin=2.4em, itemsep=1pt, topsep=2pt]
\item[\textbf{(A)}] claims R explains A, but real eigenvalues do not imply diagonalisability: $\begin{psmallmatrix}1&1\\0&1\end{psmallmatrix}$ is a counter-example.
\item[\textbf{(C)}] claims R is false, but every Hermitian matrix has real eigenvalues.
\item[\textbf{(D)}] claims A is false, but the spectral theorem guarantees a unitary diagonalisation for Hermitian matrices.
\end{itemize}
\textbf{\color{brand}Key takeaway.} Two separate facts about Hermitian matrices: unitary diagonalisability (from normality) and real eigenvalues (from $A = A^\dagger$). HHL needs both.
\end{solbox}
\qsep

\needspace{13\baselineskip}
\begin{qitem}{8}
\qmeta{Statement I/II}{Condition number of $U$ and of $A^{-1}$}
\textbf{Statement I:} Every unitary matrix has condition number $\kappa = 1$.\\[1mm]
\textbf{Statement II:} For every invertible matrix $A$, $\kappa(A^{-1}) = \kappa(A)$.
\begin{itemize}[leftmargin=2.6em, labelsep=0.4em, itemsep=2pt, topsep=3pt]
\item[\correctlabel{A}] Both Statement I and Statement II are true
\item[\optlabel{B}] Both Statement I and Statement II are false
\item[\optlabel{C}] Statement I is true, but Statement II is false
\item[\optlabel{D}] Statement I is false, but Statement II is true
\end{itemize}
\end{qitem}
\begin{solbox}{A}
\textbf{\color{brand}Correct answer.} \textcolor{okgreen}{\textbf{(A)}} Both Statement I and Statement II are true

\textbf{\color{brand}Concepts required.} The condition number is the ratio of the largest to the smallest singular value.

\textbf{\color{brand}Relevant formulae.} $\kappa(A) = \sigma_{\max}/\sigma_{\min}$; singular values of $U$ unitary $= 1$; singular values of $A^{-1}$ $= 1/\sigma_i$.

\textbf{\color{brand}Step-by-step explanation.} $\kappa(A) = \lVert A\rVert_2\lVert A^{-1}\rVert_2 = \sigma_{\max}/\sigma_{\min}$.
\textbf{I is true:} a unitary $U$ preserves lengths, so all its singular values equal $1$ and $\kappa = 1/1 = 1$.
\textbf{II is true:} the singular values of $A^{-1}$ are $1/\sigma_i$, so $\kappa(A^{-1}) = \frac{1/\sigma_{\min}}{1/\sigma_{\max}} = \kappa(A)$
(directly: $\lVert A^{-1}\rVert\,\lVert (A^{-1})^{-1}\rVert = \lVert A^{-1}\rVert\,\lVert A\rVert$). Answer: (A).

\textbf{\color{brand}Why the other options are incorrect.}
\begin{itemize}[leftmargin=2.4em, itemsep=1pt, topsep=2pt]
\item[\textbf{(B)}] ``both false'' fails because both statements follow directly from the singular values.
\item[\textbf{(C)}] rejects II, but inverting swaps the largest and smallest singular values and leaves the ratio unchanged.
\item[\textbf{(D)}] rejects I, but a unitary preserves all lengths: $\sigma_{\max} = \sigma_{\min} = 1$.
\end{itemize}
\textbf{\color{brand}Key takeaway.} Unitaries are perfectly conditioned; inverting a matrix does not change its conditioning.
\end{solbox}
\qsep

\needspace{13\baselineskip}
\begin{qitem}{9}
\qmeta{Single correct}{Sparsity of the 1-D Laplacian}
The $8\times8$ matrix of the 1-D discrete Laplacian has $2$ on the diagonal, $-1$ on the first upper and lower
off-diagonals and $0$ elsewhere. Its sparsity $s$ (maximum non-zeros per row) and its total number of non-zero entries are
\par\smallskip\noindent\begin{tabularx}{\linewidth}{@{}*{4}{L}@{}}
\optlabel{A}\,$s = 8$; $\ 64$ & \correctlabel{B}\,$s = 3$; $\ 22$ & \optlabel{C}\,$s = 2$; $\ 22$ & \optlabel{D}\,$s = 3$; $\ 24$
\end{tabularx}
\end{qitem}
\begin{solbox}{B}
\textbf{\color{brand}Correct answer.} \textcolor{okgreen}{\textbf{(B)}} $s = 3$; $\ 22$

\textbf{\color{brand}Concepts required.} Sparsity $s$ counts the maximum number of non-zero entries per row; banded matrices have constant $s$.

\textbf{\color{brand}Relevant formulae.} Tridiagonal $n\times n$: $\mathrm{nnz} = n + 2(n-1) = 3n - 2$; $\ s = 3$ for $n \ge 3$.

\textbf{\color{brand}Step-by-step explanation.} Interior rows $2, \dots, 7$ contain $-1, 2, -1$: three non-zeros. The first and last rows contain only two ($2, -1$ and $-1, 2$).
So $s = 3$. Total: diagonal $8$ + upper off-diagonal $7$ + lower off-diagonal $7$ $= 22$ (not $3\times8 = 24$, because the boundary rows have
only two entries). Such matrices are \emph{sparse}: $s$ stays $3$ for any size $d$, which is the setting in which HHL's cost
$\propto s$ (or $s^2$) is small.

\textbf{\color{brand}Why the other options are incorrect.}
\begin{itemize}[leftmargin=2.4em, itemsep=1pt, topsep=2pt]
\item[\textbf{(A)}] $s = 8$, $64$ treats the matrix as dense; most entries are $0$.
\item[\textbf{(C)}] $s = 2$ counts only the off-diagonals and forgets the diagonal $2$.
\item[\textbf{(D)}] $24 = 3\times8$ assumes every row has three non-zeros, but the first and last rows have only two.
\end{itemize}
\textbf{\color{brand}Key takeaway.} Sparse means $s \ll d$ non-zeros per row; HHL's cost depends on $s$, not on $d$.
\end{solbox}
\qsep

\needspace{13\baselineskip}
\begin{qitem}{10}
\qmeta{Single correct}{Spectrum and $\kappa$ of the Laplacian}
The eigenvalues of the $n\times n$ Laplacian of Q.9 (with $2$ on the diagonal) are $\lambda_k = 2 - 2\cos\frac{k\pi}{n+1}$, $k = 1, \dots, n$.
For $n = 3$ the condition number is
\par\smallskip\noindent\begin{tabularx}{\linewidth}{@{}*{4}{L}@{}}
\correctlabel{A}\,$3 + 2\sqrt2 \approx 5.83$ & \optlabel{B}\,$2$ & \optlabel{C}\,$\sqrt2 \approx 1.41$ & \optlabel{D}\,$1 + \sqrt2 \approx 2.41$
\end{tabularx}
\end{qitem}
\begin{solbox}{A}
\textbf{\color{brand}Correct answer.} \textcolor{okgreen}{\textbf{(A)}} $3 + 2\sqrt2 \approx 5.83$

\textbf{\color{brand}Concepts required.} The discrete Laplacian has known eigenvalues; its smallest eigenvalue shrinks like $1/n^2$ as the grid is refined.

\textbf{\color{brand}Relevant formulae.} $\lambda_k = 2 - 2\cos\frac{k\pi}{n+1}$; $\ \kappa = \frac{\lambda_n}{\lambda_1}$; $\ \cos\frac\pi4 = \frac{\sqrt2}{2}$, $(2+\sqrt2)(2-\sqrt2) = 2$.

\textbf{\color{brand}Step-by-step explanation.} For $n = 3$: $\frac{k\pi}{4}$ with $k = 1, 2, 3$ gives
$\lambda_1 = 2 - 2\cdot\frac{\sqrt2}{2} = 2 - \sqrt2 \approx 0.586$, $\lambda_2 = 2 - 0 = 2$, $\lambda_3 = 2 + \sqrt2 \approx 3.414$. Hence
\[ \kappa = \frac{2+\sqrt2}{2-\sqrt2} = \frac{(2+\sqrt2)^2}{4 - 2} = \frac{6 + 4\sqrt2}{2} = 3 + 2\sqrt2 \approx 5.83. \]
For large $n$, $\lambda_1 \approx (\pi/(n+1))^2$ and $\kappa \approx 4(n+1)^2/\pi^2$: discretised differential operators become badly
conditioned as the grid is refined, which limits HHL's advantage for such problems.

\textbf{\color{brand}Why the other options are incorrect.}
\begin{itemize}[leftmargin=2.4em, itemsep=1pt, topsep=2pt]
\item[\textbf{(B)}] $2$ is the middle eigenvalue $\lambda_2$, not a ratio.
\item[\textbf{(C)}] $\sqrt2$ is $\lambda_2/\lambda_1$ evaluated wrongly; the ratio of the extreme eigenvalues is needed.
\item[\textbf{(D)}] $1 + \sqrt2$ is $\sqrt{\kappa}$; it misses the square that appears when rationalising $\frac{2+\sqrt2}{2-\sqrt2}$.
\end{itemize}
\textbf{\color{brand}Key takeaway.} Finer discretisations of differential equations have $\kappa \propto n^2$: the $\kappa$ factor of HHL can eat the speed-up.
\end{solbox}
\qsep

\needspace{13\baselineskip}
\begin{qitem}{11}
\qmeta{Single correct}{Normalisation and amplitude encoding}
The vector $\vec b = (3, 0, 4, 0)$ is amplitude-encoded in two qubits (index $= $ binary number $q_0q_1$). Measuring both
qubits, the probability of the outcome $\ket{10}$ is
\par\smallskip\noindent\begin{tabularx}{\linewidth}{@{}*{4}{L}@{}}
\optlabel{A}\,$0.16$ & \correctlabel{B}\,$0.64$ & \optlabel{C}\,$0.36$ & \optlabel{D}\,$0.8$
\end{tabularx}
\end{qitem}
\begin{solbox}{B}
\textbf{\color{brand}Correct answer.} \textcolor{okgreen}{\textbf{(B)}} $0.64$

\textbf{\color{brand}Concepts required.} Amplitude encoding stores the normalised vector $\vec b/\lVert\vec b\rVert$ as amplitudes; probabilities are squared magnitudes.

\textbf{\color{brand}Relevant formulae.} $\ket b = \sum_i\frac{b_i}{\lVert\vec b\rVert}\ket i$; $\ P(i) = \frac{|b_i|^2}{\lVert\vec b\rVert^2}$; index $i$ in binary $q_0q_1$.

\textbf{\color{brand}Step-by-step explanation.} Normalise: $\lVert\vec b\rVert = \sqrt{9 + 16} = 5$, so
$\ket b = \tfrac35\ket{00} + \tfrac45\ket{10}$ (index $2 = 10_2$ has amplitude $4/5$).
Born rule: $P(10) = |4/5|^2 = 16/25 = 0.64$. The norm $\lVert\vec b\rVert = 5$ is \emph{not} stored in the state; it must be kept classically
(the notebook multiplies it back at the end).

\textbf{\color{brand}Why the other options are incorrect.}
\begin{itemize}[leftmargin=2.4em, itemsep=1pt, topsep=2pt]
\item[\textbf{(A)}] $0.16 = 0.4^2$ divides the entry $4$ by $10$ instead of by the norm $5$ before squaring.
\item[\textbf{(C)}] $0.36$ is the probability of $\ket{00}$ (amplitude $3/5$), not of $\ket{10}$.
\item[\textbf{(D)}] $0.8$ is the amplitude $4/5$, not the probability; the Born rule squares it.
\end{itemize}
\textbf{\color{brand}Key takeaway.} Normalise first, then square: $P(i) = |b_i|^2/\lVert\vec b\rVert^2$. The norm itself is not in the state.
\end{solbox}
\qsep

\needspace{13\baselineskip}
\begin{qitem}{12}
\qmeta{Single correct}{Solving $A\vec x=\vec b$ in the eigenbasis}
With the eigenpairs of $A = \begin{pmatrix}2&1\\1&2\end{pmatrix}$ given in Q.6 and $\vec b = (1, 0)$, write $\vec b = \sum_j\beta_j\ket{u_j}$. Then
$\vec x = \sum_j(\beta_j/\lambda_j)\ket{u_j}$ equals
\par\smallskip\noindent\begin{tabularx}{\linewidth}{@{}*{4}{L}@{}}
\optlabel{A}\,$\left(2,\ -1\right)$ & \optlabel{B}\,$\left(\tfrac12,\ -\tfrac12\right)$ & \correctlabel{C}\,$\left(\tfrac23,\ -\tfrac13\right)$ & \optlabel{D}\,$\left(\tfrac23,\ \tfrac13\right)$
\end{tabularx}
\end{qitem}
\begin{solbox}{C}
\textbf{\color{brand}Correct answer.} \textcolor{okgreen}{\textbf{(C)}} $\left(\tfrac23,\ -\tfrac13\right)$

\textbf{\color{brand}Concepts required.} Solving through the eigenbasis: expand $\vec b$, divide each coefficient by its eigenvalue, recombine.

\textbf{\color{brand}Relevant formulae.} $\beta_j = \braket{u_j|b}$; $\ \vec x = A^{-1}\vec b = \sum_j\frac{\beta_j}{\lambda_j}\ket{u_j}$.

\textbf{\color{brand}Step-by-step explanation.} $\beta_1 = \braket{u_1|b} = \tfrac{1}{\sqrt2}$, $\beta_2 = \braket{u_2|b} = \tfrac{1}{\sqrt2}$. Divide by the eigenvalues $3$ and $1$:
\[ \vec x = \frac{1}{3\sqrt2}\cdot\frac{1}{\sqrt2}\begin{pmatrix}1\\1\end{pmatrix} + \frac{1}{\sqrt2}\cdot\frac{1}{\sqrt2}\begin{pmatrix}1\\-1\end{pmatrix}
          = \begin{pmatrix}\tfrac16 + \tfrac12\\ \tfrac16 - \tfrac12\end{pmatrix} = \begin{pmatrix}\tfrac23\\ -\tfrac13\end{pmatrix}. \]
\emph{Check:} $A\vec x = (\tfrac43 - \tfrac13,\ \tfrac23 - \tfrac23) = (1, 0)$. HHL performs these three steps (expand in the eigenbasis, divide
by $\lambda_j$, recombine) coherently: phase estimation, controlled rotation, uncomputation.

\textbf{\color{brand}Why the other options are incorrect.}
\begin{itemize}[leftmargin=2.4em, itemsep=1pt, topsep=2pt]
\item[\textbf{(A)}] $(2, -1)$ multiplies by $\lambda$ instead of dividing, or forgets the normalisation of the eigenvectors; $A(2,-1) = (3, 0)$.
\item[\textbf{(B)}] $(\tfrac12, -\tfrac12)$ keeps only the $\lambda = 1$ component and drops the $\lambda = 3$ one.
\item[\textbf{(D)}] $(\tfrac23, \tfrac13)$ has the wrong sign in the second component: $A(\tfrac23, \tfrac13) = (\tfrac53, \tfrac43) \neq (1, 0)$.
\end{itemize}
\textbf{\color{brand}Key takeaway.} $A^{-1}\vec b = \sum_j(\beta_j/\lambda_j)\ket{u_j}$ is the formula HHL realises with phase estimation and a controlled rotation.
\end{solbox}
\qsep

\needspace{13\baselineskip}
\begin{qitem}{13}
\qmeta{Matching lists}{Norms, $\kappa$ and determinant}
For $A = \operatorname{diag}(4,\ 1,\ 0.5,\ 3)$ match each quantity (List-I) with its value (List-II).
\par\smallskip\noindent{\renewcommand{\arraystretch}{1.35}%
\begin{tabularx}{\linewidth}{@{}c>{\raggedright\arraybackslash}p{0.33\linewidth}@{\hspace{1.2em}}cL@{}}
\toprule
\multicolumn{2}{@{}l}{\textbf{List-I}} & \multicolumn{2}{l@{}}{\textbf{List-II}} \\ \midrule
\textbf{(P)} & spectral norm $\lVert A\rVert_2$ & \textbf{(1)} & $8$ \\
\textbf{(Q)} & spectral norm $\lVert A^{-1}\rVert_2$ & \textbf{(2)} & $4$ \\
\textbf{(R)} & condition number $\kappa(A)$ & \textbf{(3)} & $6$ \\
\textbf{(S)} & determinant $\det A$ & \textbf{(4)} & $2$ \\
\bottomrule
\end{tabularx}}
\par\smallskip\noindent\begin{tabularx}{\linewidth}{@{}*{2}{L}@{}}
\optlabel{A}\,P$\to$2, Q$\to$4, R$\to$3, S$\to$1 & \optlabel{B}\,P$\to$2, Q$\to$3, R$\to$1, S$\to$4 \\[3pt]
\optlabel{C}\,P$\to$1, Q$\to$4, R$\to$2, S$\to$3 & \correctlabel{D}\,P$\to$2, Q$\to$4, R$\to$1, S$\to$3
\end{tabularx}
\end{qitem}
\begin{solbox}{D}
\textbf{\color{brand}Correct answer.} \textcolor{okgreen}{\textbf{(D)}} P$\to$2, Q$\to$4, R$\to$1, S$\to$3

\textbf{\color{brand}Concepts required.} For a diagonal matrix the spectral norm is the largest $|$entry$|$; norms of $A^{-1}$ use reciprocals; $\det$ is the product.

\textbf{\color{brand}Relevant formulae.} $\lVert A\rVert_2 = \max|\lambda|$, $\ \lVert A^{-1}\rVert_2 = 1/\min|\lambda|$, $\ \kappa = \lVert A\rVert_2\lVert A^{-1}\rVert_2$, $\ \det A = \prod\lambda_i$.

\textbf{\color{brand}Step-by-step explanation.} \begin{itemize}\itemsep0pt
\item P: $\lVert A\rVert_2 = $ largest $|\lambda| = 4$ (2).
\item Q: $A^{-1} = \operatorname{diag}(\tfrac14, 1, 2, \tfrac13)$, so $\lVert A^{-1}\rVert_2 = 1/|\lambda|_{\min} = 1/0.5 = 2$ (4).
\item R: $\kappa = \lVert A\rVert_2\lVert A^{-1}\rVert_2 = 4\times2 = 8$ (1).
\item S: $\det A = 4\cdot1\cdot0.5\cdot3 = 6$ (3).
\end{itemize}
Note that the determinant says nothing about conditioning: $\det(0.1\,I_{10}) = 10^{-10}$ is tiny while $\kappa = 1$.

\textbf{\color{brand}Why the other options are incorrect.}
\begin{itemize}[leftmargin=2.4em, itemsep=1pt, topsep=2pt]
\item[\textbf{(A)}] R$\to$3 gives $\kappa = 6$ (the determinant) and S$\to$1 gives $\det A = 8$.
\item[\textbf{(B)}] Q$\to$3 gives $\lVert A^{-1}\rVert_2 = 6$ and S$\to$4 gives $\det A = 2$: the determinant and the inverse norm are exchanged.
\item[\textbf{(C)}] P$\to$1 gives $\lVert A\rVert_2 = 8$ (that is $\kappa$) and R$\to$2 gives $\kappa = 4$ (that is $\lVert A\rVert_2$).
\end{itemize}
\textbf{\color{brand}Key takeaway.} $\kappa = \lVert A\rVert\,\lVert A^{-1}\rVert$; the determinant is unrelated to the condition number.
\end{solbox}
\qsep

\needspace{13\baselineskip}
\begin{qitem}{14}
\qmeta{Single correct}{Hermitian embedding of a non-Hermitian matrix}
HHL needs a Hermitian matrix, so the non-Hermitian $A = \begin{pmatrix}1&2\\0&1\end{pmatrix}$ is embedded as
$\tilde A = \begin{pmatrix}0 & A\\ A^\dagger & 0\end{pmatrix}$ ($4\times4$). The eigenvalues of $\tilde A$ are
\par\smallskip\noindent\begin{tabularx}{\linewidth}{@{}*{2}{L}@{}}
\optlabel{A}\,$3 \pm 2\sqrt2$ & \optlabel{B}\,$\pm1$, each twice \\[3pt]
\optlabel{C}\,$1, 1, 1, 1$ & \correctlabel{D}\,$\pm(\sqrt2 + 1)$ and $\pm(\sqrt2 - 1)$
\end{tabularx}
\end{qitem}
\begin{solbox}{D}
\textbf{\color{brand}Correct answer.} \textcolor{okgreen}{\textbf{(D)}} $\pm(\sqrt2 + 1)$ and $\pm(\sqrt2 - 1)$

\textbf{\color{brand}Concepts required.} The Hermitian embedding $\tilde A = \begin{psmallmatrix}0&A\\A^\dagger&0\end{psmallmatrix}$ has eigenvalues $\pm\sigma_i(A)$.

\textbf{\color{brand}Relevant formulae.} $\tilde A^2 = \operatorname{diag}(AA^\dagger, A^\dagger A)$; $\ \sigma_i = \sqrt{\operatorname{eig}(A^\dagger A)}$; $\ 3 \pm 2\sqrt2 = (\sqrt2 \pm 1)^2$.

\textbf{\color{brand}Step-by-step explanation.} $\tilde A^2 = \begin{psmallmatrix}AA^\dagger & 0\\ 0 & A^\dagger A\end{psmallmatrix}$, so the eigenvalues of $\tilde A$ are $\pm\sigma_i$, the singular values of $A$
with both signs. $A^\dagger A = \begin{psmallmatrix}1&2\\2&5\end{psmallmatrix}$ has trace $6$ and determinant $1$: eigenvalues $3 \pm 2\sqrt2 = (\sqrt2 \pm 1)^2$.
Hence $\sigma = \sqrt2 \pm 1$ and the spectrum of $\tilde A$ is $\{\pm(\sqrt2 + 1), \pm(\sqrt2 - 1)\} \approx \{\pm2.414, \pm0.414\}$.
The eigenvalues of $A$ itself ($1, 1$) are irrelevant here; $\kappa(\tilde A) = \frac{\sqrt2+1}{\sqrt2-1} = 3 + 2\sqrt2 \approx 5.83$, and HHL must handle
\emph{negative} eigenvalues.

\textbf{\color{brand}Why the other options are incorrect.}
\begin{itemize}[leftmargin=2.4em, itemsep=1pt, topsep=2pt]
\item[\textbf{(A)}] $3 \pm 2\sqrt2$ are the eigenvalues of $A^\dagger A$, i.e.\ the \emph{squares} of the singular values (and they are never negative).
\item[\textbf{(B)}] $\pm1$ twice again uses the eigenvalues of $A$; the embedding is governed by singular values.
\item[\textbf{(C)}] $1, 1, 1, 1$ uses the eigenvalues of $A$ (a defective matrix), not its singular values, and forgets the $\pm$ pairs.
\end{itemize}
\textbf{\color{brand}Key takeaway.} Embedding a non-Hermitian $A$: dimension doubles, eigenvalues become $\pm\sigma_i$, and $\kappa = \sigma_{\max}/\sigma_{\min}$.
\end{solbox}
\qsep

\needspace{13\baselineskip}
\begin{qitem}{15}
\qmeta{Statement I/II}{When is $e^{iAt}$ unitary?}
\textbf{Statement I:} If $A$ is Hermitian, then $U = e^{iAt}$ is unitary for every real $t$.\\[1mm]
\textbf{Statement II:} If $A$ is any real square matrix, then $e^{iAt}$ is unitary for every real $t$.
\begin{itemize}[leftmargin=2.6em, labelsep=0.4em, itemsep=2pt, topsep=3pt]
\item[\optlabel{A}] Both Statement I and Statement II are true
\item[\optlabel{B}] Both Statement I and Statement II are false
\item[\correctlabel{C}] Statement I is true, but Statement II is false
\item[\optlabel{D}] Statement I is false, but Statement II is true
\end{itemize}
\end{qitem}
\begin{solbox}{C}
\textbf{\color{brand}Correct answer.} \textcolor{okgreen}{\textbf{(C)}} Statement I is true, but Statement II is false

\textbf{\color{brand}Concepts required.} $e^{iAt}$ is unitary iff $(e^{iAt})^\dagger e^{iAt} = I$; this holds for every $t$ exactly when $A$ is Hermitian.

\textbf{\color{brand}Relevant formulae.} $(e^{iAt})^\dagger = e^{-iA^\dagger t}$; $\ A = V\Lambda V^\dagger \Rightarrow e^{iAt} = Ve^{i\Lambda t}V^\dagger$.

\textbf{\color{brand}Step-by-step explanation.} \textbf{I is true:} with $A = V\Lambda V^\dagger$, $e^{iAt} = V\operatorname{diag}(e^{i\lambda_jt})V^\dagger$ and $|e^{i\lambda_jt}| = 1$ because $\lambda_j$ is real;
equivalently $(e^{iAt})^\dagger = e^{-iA^\dagger t} = e^{-iAt}$.
\textbf{II is false:} realness is not enough. For $A = \begin{psmallmatrix}0&1\\0&0\end{psmallmatrix}$, $e^{iAt} = I + iAt = \begin{psmallmatrix}1 & it\\ 0&1\end{psmallmatrix}$,
which maps $(0, 1)$ to $(it, 1)$ of length $\sqrt{1 + t^2} \neq 1$. Only real \emph{symmetric} matrices are guaranteed to work. This is why HHL embeds
non-Hermitian matrices first. Answer: (C).

\textbf{\color{brand}Why the other options are incorrect.}
\begin{itemize}[leftmargin=2.4em, itemsep=1pt, topsep=2pt]
\item[\textbf{(A)}] accepts II, but a real non-symmetric matrix such as $\begin{psmallmatrix}0&1\\0&0\end{psmallmatrix}$ gives a non-unitary $e^{iAt}$.
\item[\textbf{(B)}] rejects I, but Hermitian $A$ has real eigenvalues, so $e^{i\lambda t}$ are pure phases.
\item[\textbf{(D)}] rejects I and accepts II --- both judgements are reversed.
\end{itemize}
\textbf{\color{brand}Key takeaway.} Hermitian $\Leftrightarrow$ $e^{iAt}$ unitary. This is the reason HHL requires (or embeds into) a Hermitian matrix.
\end{solbox}
\qsep

\needspace{13\baselineskip}
\begin{qitem}{16}
\qmeta{Single correct}{Error amplification by $\kappa$}
A linear system has $\kappa(A) = 1000$. The right-hand side is known with relative error $\lVert\delta\vec b\rVert/\lVert\vec b\rVert = 10^{-6}$.
The best general upper bound on the relative error $\lVert\delta\vec x\rVert/\lVert\vec x\rVert$ of the solution is
\par\smallskip\noindent\begin{tabularx}{\linewidth}{@{}*{2}{L}@{}}
\correctlabel{A}\,$10^{-3}$ & \optlabel{B}\,$10^{-6}$ \\[3pt]
\optlabel{C}\,$10^{-9}$ & \optlabel{D}\,$10^{-6}/\sqrt{1000} \approx 3\times10^{-8}$
\end{tabularx}
\end{qitem}
\begin{solbox}{A}
\textbf{\color{brand}Correct answer.} \textcolor{okgreen}{\textbf{(A)}} $10^{-3}$

\textbf{\color{brand}Concepts required.} Perturbation theory for linear systems: the relative error in $\vec x$ can be $\kappa$ times the relative error in $\vec b$.

\textbf{\color{brand}Relevant formulae.} $\frac{\lVert\delta\vec x\rVert}{\lVert\vec x\rVert} \le \kappa(A)\,\frac{\lVert\delta\vec b\rVert}{\lVert\vec b\rVert}$.

\textbf{\color{brand}Step-by-step explanation.} From $A(\vec x + \delta\vec x) = \vec b + \delta\vec b$: $\delta\vec x = A^{-1}\delta\vec b$, so $\lVert\delta\vec x\rVert \le \lVert A^{-1}\rVert\,\lVert\delta\vec b\rVert$, and
$\lVert\vec b\rVert \le \lVert A\rVert\,\lVert\vec x\rVert$. Dividing:
\[ \frac{\lVert\delta\vec x\rVert}{\lVert\vec x\rVert} \le \lVert A\rVert\lVert A^{-1}\rVert\,\frac{\lVert\delta\vec b\rVert}{\lVert\vec b\rVert} = \kappa\,\frac{\lVert\delta\vec b\rVert}{\lVert\vec b\rVert} = 1000\times10^{-6} = 10^{-3}. \]
The bound is attained when $\vec b$ points along the largest and $\delta\vec b$ along the smallest singular direction. $\kappa$ measures the
intrinsic sensitivity of the problem --- no algorithm, classical or quantum, removes it.

\textbf{\color{brand}Why the other options are incorrect.}
\begin{itemize}[leftmargin=2.4em, itemsep=1pt, topsep=2pt]
\item[\textbf{(B)}] $10^{-6}$ assumes errors are not amplified, true only for $\kappa = 1$.
\item[\textbf{(C)}] $10^{-9}$ divides by $\kappa$ instead of multiplying.
\item[\textbf{(D)}] $10^{-6}/\sqrt{1000}$ has no basis: $\kappa$ multiplies the error bound.
\end{itemize}
\textbf{\color{brand}Key takeaway.} $\kappa$ is an intrinsic sensitivity: $\log_{10}\kappa$ digits can be lost, whatever the algorithm.
\end{solbox}
\qsep

\needspace{20\baselineskip}
\begin{passagebox}{Paragraph for Questions 17 and 18}
The notebook \texttt{HHL.ipynb} lists the $2\times2$ example
\[ A = \begin{pmatrix} 1 & -\tfrac13\\[2pt] -\tfrac13 & 1\end{pmatrix}, \qquad \vec b = \begin{pmatrix}1\\0\end{pmatrix}. \]
The questions below concern the classical linear algebra that HHL relies on for this system.
\end{passagebox}
\needspace{13\baselineskip}
\begin{qitem}{17}
\qmeta{Paragraph based}{$2\times2$ example: eigenvalues}
The eigenvalues of $A$ and its condition number are
\par\smallskip\noindent\begin{tabularx}{\linewidth}{@{}*{2}{L}@{}}
\optlabel{A}\,$\tfrac13$ and $1$; $\ \kappa = 3$ & \optlabel{B}\,$1$ and $-\tfrac13$; $\ \kappa = 3$ \\[3pt]
\optlabel{C}\,$\tfrac23$ and $\tfrac43$; $\ \kappa = \tfrac43$ & \correctlabel{D}\,$\tfrac23$ and $\tfrac43$; $\ \kappa = 2$
\end{tabularx}
\end{qitem}
\begin{solbox}{D}
\textbf{\color{brand}Correct answer.} \textcolor{okgreen}{\textbf{(D)}} $\tfrac23$ and $\tfrac43$; $\ \kappa = 2$

\textbf{\color{brand}Concepts required.} For $\begin{psmallmatrix}a&c\\c&a\end{psmallmatrix}$ the eigenvalues are $a \pm c$ with eigenvectors $\frac{1}{\sqrt2}(1, \pm1)$.

\textbf{\color{brand}Relevant formulae.} $\lambda = a \pm c$; $\ \kappa = \frac{|\lambda|_{\max}}{|\lambda|_{\min}}$; here $a = 1$, $c = -\frac13$.

\textbf{\color{brand}Step-by-step explanation.} $\operatorname{tr}A = 2$, $\det A = 1 - \tfrac19 = \tfrac89$. The eigenvalues solve $\lambda^2 - 2\lambda + \tfrac89 = 0$:
$\lambda = 1 \pm \sqrt{1 - \tfrac89} = 1 \pm \tfrac13$, i.e.\ $\tfrac23$ (eigenvector $\tfrac{1}{\sqrt2}(1, 1)$) and $\tfrac43$ (eigenvector $\tfrac{1}{\sqrt2}(1, -1)$).
\[ \kappa = \frac{4/3}{2/3} = 2. \]
(The diagonal $1$ and the off-diagonal $-\tfrac13$ are entries, not eigenvalues.) With $\kappa = 2$ the notebook's ``auto'' time maps $\tfrac23 \to$
clock value $1$ and $\tfrac43 \to$ clock value $2$, so HHL is exact for this matrix.

\textbf{\color{brand}Why the other options are incorrect.}
\begin{itemize}[leftmargin=2.4em, itemsep=1pt, topsep=2pt]
\item[\textbf{(A)}] $\tfrac13$ and $1$ give the trace $\tfrac43 \neq 2$; they are not roots of the characteristic polynomial.
\item[\textbf{(B)}] $1$ and $-\tfrac13$ are the matrix entries, not eigenvalues (their product $-\tfrac13 \neq \det A = \tfrac89$).
\item[\textbf{(C)}] $\kappa = \tfrac43$ divides the largest eigenvalue by $1$ (the diagonal entry) instead of by $\lambda_{\min} = \tfrac23$.
\end{itemize}
\textbf{\color{brand}Key takeaway.} Symmetric $2\times2$ with equal diagonal: eigenvalues $a \pm c$ --- a two-second calculation worth memorising.
\end{solbox}
\qsep

\needspace{13\baselineskip}
\begin{qitem}{18}
\qmeta{Paragraph based}{$2\times2$ example: normalised solution}
The solution $\vec x = A^{-1}\vec b$ and the probability of measuring the system qubit as $\ket0$ in the normalised state $\ket x$ are
\par\smallskip\noindent\begin{tabularx}{\linewidth}{@{}*{2}{L}@{}}
\optlabel{A}\,$\vec x = \left(1,\ \tfrac13\right)$; $\ \tfrac{1}{\sqrt{10}}$ & \optlabel{B}\,$\vec x = \left(\tfrac98,\ \tfrac38\right)$; $\ 0.75$ \\[3pt]
\correctlabel{C}\,$\vec x = \left(\tfrac98,\ \tfrac38\right)$; $\ 0.9$ & \optlabel{D}\,$\vec x = \left(\tfrac98,\ -\tfrac38\right)$; $\ 0.9$
\end{tabularx}
\end{qitem}
\begin{solbox}{C}
\textbf{\color{brand}Correct answer.} \textcolor{okgreen}{\textbf{(C)}} $\vec x = \left(\tfrac98,\ \tfrac38\right)$; $\ 0.9$

\textbf{\color{brand}Concepts required.} The normalised quantum state $\ket x$ keeps the ratios of the solution components; probabilities are squared normalised components.

\textbf{\color{brand}Relevant formulae.} $\vec x = A^{-1}\vec b$; $\ \ket x = \vec x/\lVert\vec x\rVert$; $\ P(i) = x_i^2/\lVert\vec x\rVert^2$.

\textbf{\color{brand}Step-by-step explanation.} $A^{-1} = \frac{1}{\det A}\begin{psmallmatrix}1&\frac13\\ \frac13&1\end{psmallmatrix} = \frac98\begin{psmallmatrix}1&\frac13\\ \frac13&1\end{psmallmatrix}$, so
$\vec x = A^{-1}(1, 0)^T = \left(\tfrac98, \tfrac38\right)$. Check: $\tfrac98 - \tfrac13\cdot\tfrac38 = \tfrac98 - \tfrac18 = 1$ and $-\tfrac13\cdot\tfrac98 + \tfrac38 = 0$.
Normalised: $\vec x \propto (3, 1)$, $\ket x = \frac{1}{\sqrt{10}}(3\ket0 + \ket1)$, so $P(0) = \frac{9}{10} = 0.9$.
The \emph{ratio} $x_1/x_2 = 3$ survives in the quantum state; the overall scale $\lVert\vec x\rVert = \tfrac38\sqrt{10}$ does not.

\textbf{\color{brand}Why the other options are incorrect.}
\begin{itemize}[leftmargin=2.4em, itemsep=1pt, topsep=2pt]
\item[\textbf{(A)}] $(1, \tfrac13)$ is $\vec x$ up to the factor $\tfrac98$, and $\tfrac{1}{\sqrt{10}}$ is an amplitude, not the probability of $\ket0$.
\item[\textbf{(B)}] $0.75$ is $x_0/(x_0 + x_1) = \frac{9/8}{12/8}$: it normalises the components linearly instead of by squares.
\item[\textbf{(D)}] $(\tfrac98, -\tfrac38)$ has the wrong sign: the off-diagonal $-\tfrac13$ of $A$ becomes $+\tfrac13$ in $A^{-1}$.
\end{itemize}
\textbf{\color{brand}Key takeaway.} Measuring $\ket x$ gives $x_i^2/\lVert\vec x\rVert^2$: ratios survive, the scale does not.
\end{solbox}
\qsep

\needspace{13\baselineskip}
\begin{qitem}{19}
\qmeta{Single correct}{Embedding + padding: qubit count}
A $5\times5$ \emph{non-Hermitian} matrix is prepared for HHL exactly as in the notebook: first the Hermitian embedding
$\begin{psmallmatrix}0&A\\A^\dagger&0\end{psmallmatrix}$, then padding with an identity block to the next power of $2$. The number of system qubits $n_b$ is
\par\smallskip\noindent\begin{tabularx}{\linewidth}{@{}*{4}{L}@{}}
\optlabel{A}\,$5$ & \correctlabel{B}\,$4$ & \optlabel{C}\,$10$ & \optlabel{D}\,$3$
\end{tabularx}
\end{qitem}
\begin{solbox}{B}
\textbf{\color{brand}Correct answer.} \textcolor{okgreen}{\textbf{(B)}} $4$

\textbf{\color{brand}Concepts required.} Two preparation steps change the size: Hermitian embedding ($d \to 2d$) and padding to a power of two.

\textbf{\color{brand}Relevant formulae.} $n_b = \lceil\log_2(2d)\rceil$ for non-Hermitian $A$; $\ n_b = \lceil\log_2 d\rceil$ for Hermitian $A$.

\textbf{\color{brand}Step-by-step explanation.} Embedding doubles the dimension: $2\times5 = 10$. The next power of two is $16 = 2^4$; the extra $6$ dimensions are filled with an
identity block (eigenvalue $1$, decoupled from the system because $\vec b$ is padded with zeros). Hence $n_b = \log_2 16 = 4$ system qubits.
(Without the embedding $5 \to 8$ would need $3$ qubits.) In general $n_b = \lceil\log_2(2d)\rceil$ for a non-Hermitian $d\times d$ matrix.

\textbf{\color{brand}Why the other options are incorrect.}
\begin{itemize}[leftmargin=2.4em, itemsep=1pt, topsep=2pt]
\item[\textbf{(A)}] $5$ counts one qubit per row; qubit counts are logarithmic in the dimension.
\item[\textbf{(C)}] $10$ is the embedded dimension, not the number of qubits.
\item[\textbf{(D)}] $3$ pads $5 \to 8$ but forgets the doubling of the embedding.
\end{itemize}
\textbf{\color{brand}Key takeaway.} Qubits $= \log_2$(dimension after embedding and padding).
\end{solbox}
\qsep

\needspace{13\baselineskip}
\begin{qitem}{20}
\qmeta{Assertion--Reason}{$\kappa$ from singular values}
\textbf{Assertion (A):} For a Hermitian positive-definite matrix the condition number equals $\lambda_{\max}/\lambda_{\min}$.\\[1mm]
\textbf{Reason (R):} The singular values of a Hermitian matrix are the absolute values of its eigenvalues.
\begin{itemize}[leftmargin=2.6em, labelsep=0.4em, itemsep=2pt, topsep=3pt]
\item[\correctlabel{A}] Both Assertion (A) and Reason (R) are true, and R is the correct explanation of A
\item[\optlabel{B}] Both A and R are true, but R is \emph{not} the correct explanation of A
\item[\optlabel{C}] A is true, but R is false
\item[\optlabel{D}] A is false, but R is true
\end{itemize}
\end{qitem}
\begin{solbox}{A}
\textbf{\color{brand}Correct answer.} \textcolor{okgreen}{\textbf{(A)}} Both Assertion (A) and Reason (R) are true, and R is the correct explanation of A

\textbf{\color{brand}Concepts required.} Singular values are the square roots of the eigenvalues of $A^\dagger A$; for Hermitian $A$ they are $|\lambda_j|$.

\textbf{\color{brand}Relevant formulae.} $\kappa = \sigma_{\max}/\sigma_{\min}$; $\ A = V\Lambda V^\dagger \Rightarrow \sigma_j = |\lambda_j|$; positive definite: $\lambda_j > 0$.

\textbf{\color{brand}Step-by-step explanation.} By definition $\kappa = \sigma_{\max}/\sigma_{\min}$ (spectral norm of $A$ times that of $A^{-1}$).
\textbf{R is true:} for $A = V\Lambda V^\dagger$, $A^\dagger A = V\Lambda^2V^\dagger$, so $\sigma_j = \sqrt{\lambda_j^2} = |\lambda_j|$.
\textbf{A is true}, and it follows from R: for a positive-definite matrix all $\lambda_j > 0$, so $|\lambda_j| = \lambda_j$ and
$\kappa = \sigma_{\max}/\sigma_{\min} = \lambda_{\max}/\lambda_{\min}$. R is the correct explanation. Answer: (A).

\textbf{\color{brand}Why the other options are incorrect.}
\begin{itemize}[leftmargin=2.4em, itemsep=1pt, topsep=2pt]
\item[\textbf{(B)}] denies the explanatory link, but A is derived directly from R plus positivity.
\item[\textbf{(C)}] rejects R, but $A^\dagger A = A^2 = V\Lambda^2V^\dagger$ gives $\sigma_j = |\lambda_j|$.
\item[\textbf{(D)}] rejects A, but for positive eigenvalues $|\lambda| = \lambda$ and the ratio is $\lambda_{\max}/\lambda_{\min}$.
\end{itemize}
\textbf{\color{brand}Key takeaway.} For Hermitian matrices $\kappa = |\lambda|_{\max}/|\lambda|_{\min}$, the definition used by the notebook.
\end{solbox}
\qsep

\newpage

\begin{sectionbanner}SECTION B \enspace$\cdot$\enspace Quantum Computing Foundations \hfill Q.21--40\end{sectionbanner}

\needspace{13\baselineskip}
\begin{qitem}{21}
\qmeta{Single correct}{Born rule}
A qubit is in the state $\ket\psi = \frac{1}{\sqrt5}\left(\ket0 + 2i\ket1\right)$. The probability of measuring $1$ is
\par\smallskip\noindent\begin{tabularx}{\linewidth}{@{}*{4}{L}@{}}
\optlabel{A}\,$\tfrac{2}{\sqrt5}$ & \optlabel{B}\,$\tfrac25$ & \correctlabel{C}\,$\tfrac45$ & \optlabel{D}\,$-\tfrac45$
\end{tabularx}
\end{qitem}
\begin{solbox}{C}
\textbf{\color{brand}Correct answer.} \textcolor{okgreen}{\textbf{(C)}} $\tfrac45$

\textbf{\color{brand}Concepts required.} Born rule: the probability of an outcome is the squared magnitude of its amplitude.

\textbf{\color{brand}Relevant formulae.} $P(k) = |\braket{k|\psi}|^2$; $\ |a + ib|^2 = a^2 + b^2$; $\ \sum_k P(k) = 1$.

\textbf{\color{brand}Step-by-step explanation.} Born rule: $P(1) = |\braket{1|\psi}|^2 = \left|\frac{2i}{\sqrt5}\right|^2 = \frac{|2i|^2}{5} = \frac{4}{5}$.
Check the normalisation: $\frac{1}{5}(|1|^2 + |2i|^2) = \frac{1 + 4}{5} = 1$. The factor $i$ is a \emph{relative phase}: it does not change
the computational-basis probabilities but would change the outcome of a measurement in the $X$ basis.

\textbf{\color{brand}Why the other options are incorrect.}
\begin{itemize}[leftmargin=2.4em, itemsep=1pt, topsep=2pt]
\item[\textbf{(A)}] $\tfrac{2}{\sqrt5}$ is the magnitude of the amplitude, not the probability.
\item[\textbf{(B)}] $\tfrac25$ uses $|2i|$ instead of $|2i|^2$ in the numerator.
\item[\textbf{(D)}] $-\tfrac45$ squares $2i$ as $(2i)^2 = -4$ instead of taking $|2i|^2 = 4$; probabilities are never negative.
\end{itemize}
\textbf{\color{brand}Key takeaway.} Probability $=|$amplitude$|^2$; phases such as $i$ drop out of computational-basis statistics.
\end{solbox}
\qsep

\needspace{13\baselineskip}
\begin{qitem}{22}
\qmeta{Single correct}{State-vector size and memory}
The default HHL circuit uses $7$ qubits. A state-vector simulator stores one complex double-precision number
($16$ bytes) per amplitude. The number of amplitudes and the memory of the state are
\par\smallskip\noindent\begin{tabularx}{\linewidth}{@{}*{2}{L}@{}}
\correctlabel{A}\,$128$ amplitudes; $\ 2048$ bytes & \optlabel{B}\,$7$ amplitudes; $\ 112$ bytes \\[3pt]
\optlabel{C}\,$128$ amplitudes; $\ 1024$ bytes & \optlabel{D}\,$49$ amplitudes; $\ 784$ bytes
\end{tabularx}
\end{qitem}
\begin{solbox}{A}
\textbf{\color{brand}Correct answer.} \textcolor{okgreen}{\textbf{(A)}} $128$ amplitudes; $\ 2048$ bytes

\textbf{\color{brand}Concepts required.} An $n$-qubit state has $2^n$ complex amplitudes; a simulator stores all of them.

\textbf{\color{brand}Relevant formulae.} memory $= 2^n\times16$ bytes (fp64 complex), $2^n\times8$ bytes (fp32).

\textbf{\color{brand}Step-by-step explanation.} An $n$-qubit state has $2^n$ amplitudes: $2^7 = 128$. Memory $= 16 \times 2^7 = 2048$ bytes $= 2$ KiB.
In general $16\times2^n$ bytes: every extra qubit doubles the memory, $30$ qubits need $17.2$ GB and $40$ qubits $17.6$ TB.
(A real \texttt{complex128} has $8$ bytes for the real part and $8$ for the imaginary part; $1024$ bytes would be single precision.)

\textbf{\color{brand}Why the other options are incorrect.}
\begin{itemize}[leftmargin=2.4em, itemsep=1pt, topsep=2pt]
\item[\textbf{(B)}] $7$ amplitudes treats qubits like classical bits with one number each.
\item[\textbf{(C)}] $1024$ bytes assumes $8$ bytes per amplitude (single precision), not $16$.
\item[\textbf{(D)}] $49 = 7^2$ has no meaning here; the count is $2^7$.
\end{itemize}
\textbf{\color{brand}Key takeaway.} Memory doubles with every qubit: this, not the gate count, limits classical simulation of HHL.
\end{solbox}
\qsep

\needspace{13\baselineskip}
\begin{qitem}{23}
\qmeta{Single correct}{Gate identity $HZH = X$}
The product $HZH$ equals
\par\smallskip\noindent\begin{tabularx}{\linewidth}{@{}*{4}{L}@{}}
\optlabel{A}\,$Y$ & \optlabel{B}\,$Z$ & \correctlabel{C}\,$X$ & \optlabel{D}\,$-X$
\end{tabularx}
\end{qitem}
\begin{solbox}{C}
\textbf{\color{brand}Correct answer.} \textcolor{okgreen}{\textbf{(C)}} $X$

\textbf{\color{brand}Concepts required.} Conjugating by $H$ exchanges the $X$ and $Z$ axes.

\textbf{\color{brand}Relevant formulae.} $HZH = X$, $\ HXH = Z$, $\ HYH = -Y$, $\ H^2 = I$.

\textbf{\color{brand}Step-by-step explanation.} \[ HZH = \frac12\begin{pmatrix}1&1\\1&-1\end{pmatrix}\begin{pmatrix}1&0\\0&-1\end{pmatrix}\begin{pmatrix}1&1\\1&-1\end{pmatrix}
         = \frac12\begin{pmatrix}1&-1\\1&1\end{pmatrix}\begin{pmatrix}1&1\\1&-1\end{pmatrix} = \frac12\begin{pmatrix}0&2\\2&0\end{pmatrix} = X. \]
The Hadamard exchanges the $Z$ and $X$ bases ($H\ket0 = \ket+$, $H\ket1 = \ket-$), so it maps $Z$ to $X$ and vice versa ($HXH = Z$).

\textbf{\color{brand}Why the other options are incorrect.}
\begin{itemize}[leftmargin=2.4em, itemsep=1pt, topsep=2pt]
\item[\textbf{(A)}] $Y$ is the image of neither $X$ nor $Z$ under $H$-conjugation ($HYH = -Y$).
\item[\textbf{(B)}] $Z$ would require $H$ to commute with $Z$, which it does not.
\item[\textbf{(D)}] $-X$ has the wrong sign: the explicit product gives $+X$.
\end{itemize}
\textbf{\color{brand}Key takeaway.} $H$ switches between the $Z$ basis and the $X$ basis; this is the basis of phase kickback and QPE.
\end{solbox}
\qsep

\needspace{13\baselineskip}
\begin{qitem}{24}
\qmeta{Matching lists}{$R_Y$ rotations of $\ket0$}
Match each rotation applied to $\ket0$ (List-I) with the resulting state (List-II).
\par\smallskip\noindent{\renewcommand{\arraystretch}{1.35}%
\begin{tabularx}{\linewidth}{@{}c>{\raggedright\arraybackslash}p{0.33\linewidth}@{\hspace{1.2em}}cL@{}}
\toprule
\multicolumn{2}{@{}l}{\textbf{List-I}} & \multicolumn{2}{l@{}}{\textbf{List-II}} \\ \midrule
\textbf{(P)} & $R_Y(\pi)\ket0$ & \textbf{(1)} & $\ket+$ \\
\textbf{(Q)} & $R_Y(\pi/2)\ket0$ & \textbf{(2)} & $\frac{\sqrt3}{2}\ket0 + \frac12\ket1$ \\
\textbf{(R)} & $R_Y(-\pi/2)\ket0$ & \textbf{(3)} & $\ket1$ \\
\textbf{(S)} & $R_Y(\pi/3)\ket0$ & \textbf{(4)} & $\ket-$ \\
\bottomrule
\end{tabularx}}
\par\smallskip\noindent\begin{tabularx}{\linewidth}{@{}*{2}{L}@{}}
\optlabel{A}\,P$\to$1, Q$\to$3, R$\to$4, S$\to$2 & \correctlabel{B}\,P$\to$3, Q$\to$1, R$\to$4, S$\to$2 \\[3pt]
\optlabel{C}\,P$\to$3, Q$\to$4, R$\to$1, S$\to$2 & \optlabel{D}\,P$\to$3, Q$\to$1, R$\to$2, S$\to$4
\end{tabularx}
\end{qitem}
\begin{solbox}{B}
\textbf{\color{brand}Correct answer.} \textcolor{okgreen}{\textbf{(B)}} P$\to$3, Q$\to$1, R$\to$4, S$\to$2

\textbf{\color{brand}Concepts required.} $R_Y(\theta)$ rotates the Bloch vector in the $x$--$z$ plane; on $\ket0$ it gives real amplitudes $\cos\frac\theta2$ and $\sin\frac\theta2$.

\textbf{\color{brand}Relevant formulae.} $R_Y(\theta)\ket0 = \cos\frac\theta2\ket0 + \sin\frac\theta2\ket1$.

\textbf{\color{brand}Step-by-step explanation.} $R_Y(\theta)\ket0 = \cos\frac\theta2\ket0 + \sin\frac\theta2\ket1$ (first column of the matrix in the conventions).
\begin{itemize}\itemsep0pt
\item P: $\theta = \pi$: $\cos\frac\pi2\ket0 + \sin\frac\pi2\ket1 = \ket1$ (3).
\item Q: $\theta = \frac\pi2$: $\frac{1}{\sqrt2}(\ket0 + \ket1) = \ket+$ (1).
\item R: $\theta = -\frac\pi2$: $\cos(-\frac\pi4) = \frac{1}{\sqrt2}$, $\sin(-\frac\pi4) = -\frac{1}{\sqrt2}$, giving $\ket-$ (4).
\item S: $\theta = \frac\pi3$: $\cos\frac\pi6\ket0 + \sin\frac\pi6\ket1 = \frac{\sqrt3}{2}\ket0 + \frac12\ket1$ (2).
\end{itemize}
In HHL the angle $\theta = 2\arcsin(C/\lambda)$ is chosen so that the amplitude of $\ket1$ is exactly $\sin\frac\theta2 = C/\lambda$.

\textbf{\color{brand}Why the other options are incorrect.}
\begin{itemize}[leftmargin=2.4em, itemsep=1pt, topsep=2pt]
\item[\textbf{(A)}] P$\to$1, Q$\to$3 exchanges $R_Y(\pi)$ and $R_Y(\pi/2)$: the half-angle means $\theta = \pi$ is a full flip.
\item[\textbf{(C)}] Q$\to$4, R$\to$1 exchanges $\ket+$ and $\ket-$: a positive angle gives a positive $\ket1$ amplitude.
\item[\textbf{(D)}] R$\to$2, S$\to$4 assigns $\ket-$ to $\theta = \pi/3$ and the $\frac{\sqrt3}2, \frac12$ state to $\theta = -\pi/2$.
\end{itemize}
\textbf{\color{brand}Key takeaway.} The half-angle: $R_Y(\theta)$ puts amplitude $\sin\frac\theta2$ on $\ket1$, hence HHL uses $\theta = 2\arcsin(C/\lambda)$.
\end{solbox}
\qsep

\needspace{13\baselineskip}
\begin{qitem}{25}
\qmeta{Single correct}{Entanglement test}
Which of the following two-qubit states is \emph{entangled}?
\par\smallskip\noindent\begin{tabularx}{\linewidth}{@{}*{2}{L}@{}}
\optlabel{A}\,$\tfrac12(\ket{00} + \ket{01} + \ket{10} + \ket{11})$ & \optlabel{B}\,$\tfrac12(\ket{00} - \ket{01} + \ket{10} - \ket{11})$ \\[3pt]
\correctlabel{C}\,$\tfrac12(\ket{00} + \ket{01} + \ket{10} - \ket{11})$ & \optlabel{D}\,$\tfrac12(\ket{00} + \ket{01} - \ket{10} - \ket{11})$
\end{tabularx}
\end{qitem}
\begin{solbox}{C}
\textbf{\color{brand}Correct answer.} \textcolor{okgreen}{\textbf{(C)}} $\tfrac12(\ket{00} + \ket{01} + \ket{10} - \ket{11})$

\textbf{\color{brand}Concepts required.} A two-qubit pure state is a product state iff its $2\times2$ coefficient matrix has rank 1 (determinant $0$).

\textbf{\color{brand}Relevant formulae.} $\ket\psi = \sum a_{ij}\ket{ij}$ is a product $\iff a_{00}a_{11} - a_{01}a_{10} = 0$.

\textbf{\color{brand}Step-by-step explanation.} Write $\ket\psi = \sum a_{ij}\ket{ij}$. It is a product state iff the coefficient matrix $\begin{psmallmatrix}a_{00}&a_{01}\\a_{10}&a_{11}\end{psmallmatrix}$
has rank 1, i.e.\ $a_{00}a_{11} - a_{01}a_{10} = 0$.
\begin{itemize}\itemsep0pt
\item $(1, 1, 1, -1)/2$: $1\cdot(-1) - 1\cdot1 = -2 \neq 0$: \textbf{entangled}.
\item $(1, 1, 1, 1)/2 = \ket+\ket+$; $(1, -1, 1, -1)/2 = \ket+\ket-$; $(1, 1, -1, -1)/2 = \ket-\ket+$: determinant $0$, product states.
\end{itemize}
In HHL, phase estimation entangles the clock with the system ($\sum_j\beta_j\ket{u_j}\ket{k_j}$); uncomputation removes this entanglement again.

\textbf{\color{brand}Why the other options are incorrect.}
\begin{itemize}[leftmargin=2.4em, itemsep=1pt, topsep=2pt]
\item[\textbf{(A)}] $\tfrac12(1,1,1,1) = \ket+\ket+$: determinant $1 - 1 = 0$, a product state.
\item[\textbf{(B)}] $\tfrac12(1,-1,1,-1) = \ket+\ket-$: determinant $-1 + 1 = 0$, a product state.
\item[\textbf{(D)}] $\tfrac12(1,1,-1,-1) = \ket-\ket+$: determinant $-1 + 1 = 0$, a product state.
\end{itemize}
\textbf{\color{brand}Key takeaway.} Entanglement test for two qubits: $a_{00}a_{11} \neq a_{01}a_{10}$.
\end{solbox}
\qsep

\needspace{13\baselineskip}
\begin{qitem}{26}
\qmeta{Single correct}{Phase kickback}
The control qubit is in $\ket+$ and the target in an eigenstate $\ket u$ of $U$ with $U\ket u = e^{i\phi}\ket u$. After a controlled-$U$
the joint state is
\begin{itemize}[leftmargin=2.6em, labelsep=0.4em, itemsep=2pt, topsep=3pt]
\item[\optlabel{A}] $\ket+\otimes e^{i\phi}\ket u$ (the phase stays on the target)
\item[\optlabel{B}] $\frac{1}{\sqrt2}\left(\ket0\ket u + \ket1 U^\dagger\ket u\right)$, entangled
\item[\optlabel{C}] $\frac{1}{\sqrt2}\left(e^{i\phi}\ket0 + \ket1\right)\otimes\ket u$
\item[\correctlabel{D}] $\frac{1}{\sqrt2}\left(\ket0 + e^{i\phi}\ket1\right)\otimes\ket u$
\end{itemize}
\end{qitem}
\begin{solbox}{D}
\textbf{\color{brand}Correct answer.} \textcolor{okgreen}{\textbf{(D)}} $\frac{1}{\sqrt2}\left(\ket0 + e^{i\phi}\ket1\right)\otimes\ket u$

\textbf{\color{brand}Concepts required.} Phase kickback: a controlled-$U$ acting on an eigenstate of $U$ moves the eigenvalue phase onto the control qubit.

\textbf{\color{brand}Relevant formulae.} $\text{C-}U(\alpha\ket0 + \beta\ket1)\ket u = (\alpha\ket0 + \beta e^{i\phi}\ket1)\ket u$ if $U\ket u = e^{i\phi}\ket u$.

\textbf{\color{brand}Step-by-step explanation.} Linearity: $\text{C-}U\,\frac{1}{\sqrt2}(\ket0 + \ket1)\ket u = \frac{1}{\sqrt2}\left(\ket0\ket u + \ket1U\ket u\right) = \frac{1}{\sqrt2}\left(\ket0\ket u + e^{i\phi}\ket1\ket u\right)$.
Because $U\ket u$ is $\ket u$ times a number, the number can be attached to the control: $\frac{1}{\sqrt2}(\ket0 + e^{i\phi}\ket1)\otimes\ket u$.
This is \textbf{phase kickback}: the eigenvalue becomes a \emph{relative} phase of the control, the target is unchanged and no entanglement
is created. Phase estimation reads these kicked-back phases with the inverse QFT.

\textbf{\color{brand}Why the other options are incorrect.}
\begin{itemize}[leftmargin=2.4em, itemsep=1pt, topsep=2pt]
\item[\textbf{(A)}] ``The phase stays on the target'': $e^{i\phi}$ multiplies only the $\ket1$-control branch, so it is a \emph{relative} phase of the control, not a global factor.
\item[\textbf{(B)}] $U^\dagger$ appears nowhere and, since $\ket u$ is an eigenstate, no entanglement is created.
\item[\textbf{(C)}] $(e^{i\phi}\ket0 + \ket1)$ puts the phase on the wrong branch: the $\ket0$ branch is untouched by a controlled gate.
\end{itemize}
\textbf{\color{brand}Key takeaway.} Eigenvalue phases are read off the control register --- the core of phase estimation and hence of HHL.
\end{solbox}
\qsep

\needspace{13\baselineskip}
\begin{qitem}{27}
\qmeta{Single correct}{Inverse QFT of a phase state}
The inverse QFT on two qubits is applied to $\ket\psi = \frac12\sum_{y=0}^{3}e^{2\pi i\,3y/4}\ket y$. The result is
\par\smallskip\noindent\begin{tabularx}{\linewidth}{@{}*{4}{L}@{}}
\optlabel{A}\,$\frac12\sum_{y}\ket y$ & \optlabel{B}\,$\ket{10}$ & \optlabel{C}\,$\ket{01}$ & \correctlabel{D}\,$\ket{11}$
\end{tabularx}
\end{qitem}
\begin{solbox}{D}
\textbf{\color{brand}Correct answer.} \textcolor{okgreen}{\textbf{(D)}} $\ket{11}$

\textbf{\color{brand}Concepts required.} The QFT maps $\ket x$ to a state whose phases advance by $2\pi x/N$ per step; the inverse QFT reads $x$ back.

\textbf{\color{brand}Relevant formulae.} $\mathrm{QFT}\ket x = \frac{1}{\sqrt N}\sum_y e^{2\pi ixy/N}\ket y$; $\ \mathrm{QFT}^{-1}\mathrm{QFT}\ket x = \ket x$.

\textbf{\color{brand}Step-by-step explanation.} By definition $\mathrm{QFT}\ket x = \frac12\sum_y e^{2\pi i\,xy/4}\ket y$; for $x = 3$ this is exactly $\ket\psi$. Hence
$\mathrm{QFT}^{-1}\ket\psi = \ket3 = \ket{11}$. Explicitly the amplitudes are $\frac12(1, -i, -1, i)$: the phase advances by $3\cdot\frac{2\pi}{4}$ per step.
(The equivalent reading $e^{2\pi i\,3y/4} = e^{-2\pi i\,y/4}$ means ``$-1$'', which is $3$ modulo $4$ --- not $1$.)

\textbf{\color{brand}Why the other options are incorrect.}
\begin{itemize}[leftmargin=2.4em, itemsep=1pt, topsep=2pt]
\item[\textbf{(A)}] The inverse QFT outputs the uniform superposition only for the input $\ket{00}$ ($\mathrm{QFT}^{-1}\ket{00} = \frac12\sum_y\ket y$); a phase gradient maps to a single basis state.
\item[\textbf{(B)}] $\ket{10} = \ket2$ would need amplitudes $\frac12(1, -1, 1, -1)$.
\item[\textbf{(C)}] $\ket{01}$ reads $e^{2\pi i\,3y/4} = e^{-2\pi iy/4}$ as ``$+1$''; it is $-1 \equiv 3 \pmod 4$.
\end{itemize}
\textbf{\color{brand}Key takeaway.} A linear phase gradient $e^{2\pi ixy/N}$ is $\mathrm{QFT}\ket x$; the inverse QFT turns it into $\ket x$.
\end{solbox}
\qsep

\needspace{13\baselineskip}
\begin{qitem}{28}
\qmeta{Single correct}{QPE with an exact 3-bit phase}
Phase estimation with $n_c = 3$ clock qubits is run on an eigenstate with $\varphi = 0.375$. The measured clock bit string
($c_0$ = most significant bit) is
\par\smallskip\noindent\begin{tabularx}{\linewidth}{@{}*{2}{L}@{}}
\optlabel{A}\,$100$ with probability $1$ & \correctlabel{B}\,$011$ with probability $1$ \\[3pt]
\optlabel{C}\,$011$ with probability $\approx 0.81$ & \optlabel{D}\,$110$ with probability $1$
\end{tabularx}
\end{qitem}
\begin{solbox}{B}
\textbf{\color{brand}Correct answer.} \textcolor{okgreen}{\textbf{(B)}} $011$ with probability $1$

\textbf{\color{brand}Concepts required.} QPE is exact when $\varphi\,2^{n_c}$ is an integer: the clock ends in a single basis state.

\textbf{\color{brand}Relevant formulae.} $k = \varphi\,2^{n_c}$; $\ 0.375 = 0.011_2$; $\ \tilde\varphi = k/2^{n_c}$.

\textbf{\color{brand}Step-by-step explanation.} $\varphi = 0.375 = \tfrac38 = 0.011_2$ is an exact 3-bit binary fraction. After the controlled powers the clock holds
$\frac{1}{\sqrt8}\sum_k e^{2\pi i\,3k/8}\ket k = \mathrm{QFT}\ket3$, so the inverse QFT returns $\ket{3} = \ket{011}$ with certainty and
$\tilde\varphi = 3/8$. Probabilities below $1$ (like the $4/\pi^2 \approx 0.405$ bound) occur only when $\varphi\,2^{n_c}$ is \emph{not} an integer.
$110$ is the bit-reversed string (a classic qubit-ordering slip); $100 = 4$ would mean $\varphi = 0.5$.

\textbf{\color{brand}Why the other options are incorrect.}
\begin{itemize}[leftmargin=2.4em, itemsep=1pt, topsep=2pt]
\item[\textbf{(A)}] $100 = 4$ corresponds to $\varphi = 0.5$.
\item[\textbf{(C)}] $0.81$ (or $4/\pi^2$-type values) apply only to non-representable phases; here the phase is exact.
\item[\textbf{(D)}] $110$ reverses the bit order: $0.011_2$ read backwards.
\end{itemize}
\textbf{\color{brand}Key takeaway.} Exact binary phase $\Rightarrow$ deterministic QPE outcome. HHL is exact when every $\lambda_jt2^{n_c}/2\pi$ is an integer.
\end{solbox}
\qsep

\needspace{13\baselineskip}
\begin{qitem}{29}
\qmeta{Single correct}{QPE with a non-representable phase}
Phase estimation with $n_c = 3$ clock qubits is run on an eigenstate with $\varphi = \tfrac13$. The most likely outcome $k$ and its
probability are
\par\smallskip\noindent\begin{tabularx}{\linewidth}{@{}*{4}{L}@{}}
\correctlabel{A}\,$k = 3$; $\ \approx 0.69$ & \optlabel{B}\,$k = 3$; $\ \approx 0.41$ & \optlabel{C}\,$k = 3$; $\ 1$ & \optlabel{D}\,$k = 2$; $\ \approx 0.69$
\end{tabularx}
\end{qitem}
\begin{solbox}{A}
\textbf{\color{brand}Correct answer.} \textcolor{okgreen}{\textbf{(A)}} $k = 3$; $\ \approx 0.69$

\textbf{\color{brand}Concepts required.} For non-representable phases QPE outputs a distribution peaked at the nearest integer(s) to $\varphi\,2^{n_c}$.

\textbf{\color{brand}Relevant formulae.} $P(k) = \frac{\sin^2(\pi2^{n_c}\delta)}{2^{2n_c}\sin^2(\pi\delta)}$, $\delta = \varphi - k/2^{n_c}$; worst case $P_{\text{best}} \ge 4/\pi^2$.

\textbf{\color{brand}Step-by-step explanation.} $\varphi\,2^{n_c} = 8/3 \approx 2.667$ is not an integer, so the outcome is random, peaked at the nearest integer $k = 3$.
For outcome $k$, with $\delta = \varphi - k/8$:
\[ P(k) = \left|\frac18\sum_{m=0}^{7}e^{2\pi i\,m\delta}\right|^2 = \frac{\sin^2(8\pi\delta)}{64\sin^2(\pi\delta)}. \]
For $k = 3$: $\delta = \frac13 - \frac38 = -\frac1{24}$, $\sin^2(\pi/3) = 0.75$, $64\sin^2(\pi/24) = 64\times0.017037 = 1.0904$, so $P(3) \approx 0.688$.
(For $k = 2$: $\delta = \frac1{12}$, $P(2) = \frac{\sin^2(2\pi/3)}{64\sin^2(\pi/12)} = \frac{0.75}{4.287} \approx 0.175$.) The bound $4/\pi^2 \approx 0.405$
is only a worst-case guarantee. In HHL such rounding of $\lambda t2^{n_c}/2\pi$ is the main source of error.

\textbf{\color{brand}Why the other options are incorrect.}
\begin{itemize}[leftmargin=2.4em, itemsep=1pt, topsep=2pt]
\item[\textbf{(B)}] $0.41 \approx 4/\pi^2$ is the worst-case lower bound (reached when $\varphi2^{n_c}$ is half-way between integers), not the value here.
\item[\textbf{(C)}] Probability $1$ requires $8/3$ to be an integer; it is not.
\item[\textbf{(D)}] $k = 2$ is the second-nearest value ($|8/3 - 2| = 0.67 > |8/3 - 3| = 0.33$); its probability is only $0.175$.
\end{itemize}
\textbf{\color{brand}Key takeaway.} Rounding in QPE spreads probability over neighbouring clock values; in HHL this means applying slightly wrong $1/\lambda$ factors.
\end{solbox}
\qsep

\needspace{13\baselineskip}
\begin{qitem}{30}
\qmeta{Statement I/II}{Unitarity vs.\ measurement}
\textbf{Statement I:} A unitary gate preserves the norm of every state vector.\\[1mm]
\textbf{Statement II:} A projective measurement is reversible: the pre-measurement state can always be recovered from the post-measurement state.
\begin{itemize}[leftmargin=2.6em, labelsep=0.4em, itemsep=2pt, topsep=3pt]
\item[\optlabel{A}] Both Statement I and Statement II are true
\item[\optlabel{B}] Both Statement I and Statement II are false
\item[\correctlabel{C}] Statement I is true, but Statement II is false
\item[\optlabel{D}] Statement I is false, but Statement II is true
\end{itemize}
\end{qitem}
\begin{solbox}{C}
\textbf{\color{brand}Correct answer.} \textcolor{okgreen}{\textbf{(C)}} Statement I is true, but Statement II is false

\textbf{\color{brand}Concepts required.} Unitary evolution is reversible and norm-preserving; projective measurement is irreversible.

\textbf{\color{brand}Relevant formulae.} $U^\dagger U = I \Rightarrow \lVert U\psi\rVert = \lVert\psi\rVert$; measurement: $\ket\psi \to \frac{P_k\ket\psi}{\lVert P_k\ket\psi\rVert}$.

\textbf{\color{brand}Step-by-step explanation.} \textbf{I is true:} $\lVert U\psi\rVert^2 = \braket{\psi|U^\dagger U|\psi} = \braket{\psi|\psi}$.
\textbf{II is false:} measurement maps $\sum_k a_k\ket k$ to a single $\ket k$; all the other amplitudes are lost and different states
(e.g.\ $\ket+$ and $\ket-$) give the same post-measurement states. This irreversibility is why HHL's ancilla measurement must be the
\emph{last} step on the system, and why a failed run (ancilla $= 0$) cannot be repaired --- it must be repeated. Answer: (C).

\textbf{\color{brand}Why the other options are incorrect.}
\begin{itemize}[leftmargin=2.4em, itemsep=1pt, topsep=2pt]
\item[\textbf{(A)}] accepts II, but different states ($\ket+$, $\ket-$) give the same post-measurement states, so measurement cannot be undone.
\item[\textbf{(B)}] rejects I, but unitarity is exactly norm preservation.
\item[\textbf{(D)}] rejects I and accepts II --- both reversed.
\end{itemize}
\textbf{\color{brand}Key takeaway.} Measure last. In HHL the ancilla measurement ends the computation; a failure means starting again.
\end{solbox}
\qsep

\needspace{13\baselineskip}
\begin{qitem}{31}
\qmeta{Assertion--Reason}{Global vs.\ relative phase}
\textbf{Assertion (A):} The states $\ket\psi$ and $e^{i\alpha}\ket\psi$ give identical statistics for every measurement.\\[1mm]
\textbf{Reason (R):} The relative phase between the components of a superposition affects interference and hence measurement results.
\begin{itemize}[leftmargin=2.6em, labelsep=0.4em, itemsep=2pt, topsep=3pt]
\item[\optlabel{A}] Both Assertion (A) and Reason (R) are true, and R is the correct explanation of A
\item[\correctlabel{B}] Both A and R are true, but R is \emph{not} the correct explanation of A
\item[\optlabel{C}] A is true, but R is false
\item[\optlabel{D}] A is false, but R is true
\end{itemize}
\end{qitem}
\begin{solbox}{B}
\textbf{\color{brand}Correct answer.} \textcolor{okgreen}{\textbf{(B)}} Both A and R are true, but R is \emph{not} the correct explanation of A

\textbf{\color{brand}Concepts required.} Global phases are unobservable; relative phases change interference.

\textbf{\color{brand}Relevant formulae.} $|e^{i\alpha}\braket{m|\psi}|^2 = |\braket{m|\psi}|^2$; $\ H\ket\pm = \ket0, \ket1$.

\textbf{\color{brand}Step-by-step explanation.} \textbf{A is true:} every probability has the form $|\braket{m|\psi}|^2$ and $|e^{i\alpha}\braket{m|\psi}|^2 = |\braket{m|\psi}|^2$ --- a \emph{global} phase is unobservable.
\textbf{R is true:} e.g.\ $\ket+$ and $\ket-$ differ only by a relative phase, yet $H$ maps them to $\ket0$ and $\ket1$.
R does not explain A: A is about a phase common to all components, R about phases \emph{between} components. Answer: (B).
(In HHL, the sign of $C/\lambda_j$ is a relative phase between eigencomponents and therefore matters.)

\textbf{\color{brand}Why the other options are incorrect.}
\begin{itemize}[leftmargin=2.4em, itemsep=1pt, topsep=2pt]
\item[\textbf{(A)}] claims R explains A, but R concerns relative phases and A global ones; A follows from $|e^{i\alpha}| = 1$.
\item[\textbf{(C)}] rejects R, but $\ket+$ and $\ket-$ differ only by a relative phase and are perfectly distinguishable.
\item[\textbf{(D)}] rejects A, but every probability is unchanged by a common phase.
\end{itemize}
\textbf{\color{brand}Key takeaway.} Global phase: irrelevant. Relative phase (e.g.\ the sign of $1/\lambda_j$ in HHL): essential.
\end{solbox}
\qsep

\needspace{13\baselineskip}
\begin{qitem}{32}
\qmeta{Single correct}{Controlled-$R_Y$ and ancilla probability}
A controlled-$R_Y(\theta)$ with $\theta = 2\arcsin\tfrac12$ acts on control $\ket1$ and target (ancilla) $\ket0$. The probability of then
measuring the ancilla as $1$ is
\par\smallskip\noindent\begin{tabularx}{\linewidth}{@{}*{4}{L}@{}}
\correctlabel{A}\,$\tfrac14$ & \optlabel{B}\,$\tfrac34$ & \optlabel{C}\,$\tfrac1{16}$ & \optlabel{D}\,$\tfrac12$
\end{tabularx}
\end{qitem}
\begin{solbox}{A}
\textbf{\color{brand}Correct answer.} \textcolor{okgreen}{\textbf{(A)}} $\tfrac14$

\textbf{\color{brand}Concepts required.} A controlled rotation acts only when the control is $\ket1$; then the ancilla amplitude on $\ket1$ is $\sin\frac\theta2$.

\textbf{\color{brand}Relevant formulae.} $\theta = 2\arcsin c \Rightarrow R_Y(\theta)\ket0 = \sqrt{1 - c^2}\ket0 + c\ket1$, $\ P(1) = c^2$.

\textbf{\color{brand}Step-by-step explanation.} $\theta = 2\arcsin\frac12 = 2\cdot\frac\pi6 = \frac\pi3$. The control is $\ket1$, so the rotation acts:
$R_Y(\pi/3)\ket0 = \cos\frac\pi6\ket0 + \sin\frac\pi6\ket1 = \frac{\sqrt3}{2}\ket0 + \frac12\ket1$. Hence $P(1) = (\frac12)^2 = \frac14$.
In HHL this is the eigenvalue inversion for $C/\lambda = \frac12$ (e.g.\ $C = 1$, $\lambda = 2$): the amplitude is $C/\lambda$ and the success
probability of that component is $(C/\lambda)^2$.

\textbf{\color{brand}Why the other options are incorrect.}
\begin{itemize}[leftmargin=2.4em, itemsep=1pt, topsep=2pt]
\item[\textbf{(B)}] $\tfrac34$ is the probability of the ancilla staying in $\ket0$.
\item[\textbf{(C)}] $\tfrac1{16}$ squares twice ($c^4$).
\item[\textbf{(D)}] $\tfrac12$ is the amplitude $c$, not the probability $c^2$.
\end{itemize}
\textbf{\color{brand}Key takeaway.} $\theta = 2\arcsin c$ loads amplitude $c$ and probability $c^2$ onto $\ket1$ --- HHL's eigenvalue inversion.
\end{solbox}
\qsep

\needspace{13\baselineskip}
\begin{qitem}{33}
\qmeta{Single correct}{Selecting a clock value with $X$ gates}
A 3-qubit clock register $c_0c_1c_2$ ($c_0$ = most significant bit) controls a multi-controlled $R_Y$ that fires only when all controls are $1$.
To make it act only on the clock value $k = 2$, the code places $X$ gates before and after the rotation. Which $X$ gates are needed?
\begin{itemize}[leftmargin=2.6em, labelsep=0.4em, itemsep=2pt, topsep=3pt]
\item[\optlabel{A}] $X$ on $c_1$ and $c_2$, before and after (4 gates)
\item[\optlabel{B}] $X$ on $c_1$ only, before and after (2 gates)
\item[\correctlabel{C}] $X$ on $c_0$ and $c_2$, before and after (4 gates)
\item[\optlabel{D}] none: $k = 2$ is selected by the rotation angle
\end{itemize}
\end{qitem}
\begin{solbox}{C}
\textbf{\color{brand}Correct answer.} \textcolor{okgreen}{\textbf{(C)}} $X$ on $c_0$ and $c_2$, before and after (4 gates)

\textbf{\color{brand}Concepts required.} Multi-controlled gates fire on all-ones; to condition on another pattern, flip the controls that should be $0$ (and flip them back).

\textbf{\color{brand}Relevant formulae.} $k = \sum_i c_i2^{n_c-1-i}$ ($c_0$ = MSB); flip every $c_i = 0$; gates $= 2\times$(number of zero bits).

\textbf{\color{brand}Step-by-step explanation.} $k = 2 = 010_2$: $c_0 = 0$, $c_1 = 1$, $c_2 = 0$. A multi-controlled gate fires on $\ket{111}$, so every control that must be $0$
is flipped by $X$ first: $X$ on $c_0$ and $c_2$ maps $\ket{010} \to \ket{111}$. After the rotation the same $X$ gates restore the clock.
Total $2\times2 = 4$ $X$ gates. This is exactly the \texttt{zero\_bits} loop in \texttt{HHL.ipynb}: \texttt{kernel.x(clock[i])} for every bit $i$ of $k$ that is $0$.

\textbf{\color{brand}Why the other options are incorrect.}
\begin{itemize}[leftmargin=2.4em, itemsep=1pt, topsep=2pt]
\item[\textbf{(A)}] $X$ on $c_1$ and $c_2$ selects $100_2 = 4$; $c_0$ is the most significant bit.
\item[\textbf{(B)}] $X$ on $c_1$ flips the bit that is already $1$ and selects $k = 101_2 = 5$ instead of $2$.
\item[\textbf{(D)}] The rotation angle does not select anything: without $X$ gates the gate fires only for $k = 7$.
\end{itemize}
\textbf{\color{brand}Key takeaway.} $X$--controlled gate--$X$ ``polarity'' trick: one multi-controlled rotation per clock value in \texttt{HHL.ipynb}.
\end{solbox}
\qsep

\needspace{13\baselineskip}
\begin{qitem}{34}
\qmeta{Single correct}{Post-selection on an ancilla}
An ancilla $a$ and a system register are in $\frac12\ket0_a\ket{\psi_0} + \frac{\sqrt3}{2}\ket1_a\ket{\psi_1}$ ($\ket{\psi_0}, \ket{\psi_1}$ normalised).
The ancilla is measured and the result $1$ is kept (post-selection). The probability of this and the resulting system state are
\par\smallskip\noindent\begin{tabularx}{\linewidth}{@{}*{2}{L}@{}}
\correctlabel{A}\,$\tfrac34$; $\ \ket{\psi_1}$ & \optlabel{B}\,$\tfrac{\sqrt3}{2}$; $\ \ket{\psi_1}$ \\[3pt]
\optlabel{C}\,$\tfrac14$; $\ \ket{\psi_0}$ & \optlabel{D}\,$\tfrac34$; $\ \frac{1}{\sqrt2}(\ket{\psi_0} + \ket{\psi_1})$
\end{tabularx}
\end{qitem}
\begin{solbox}{A}
\textbf{\color{brand}Correct answer.} \textcolor{okgreen}{\textbf{(A)}} $\tfrac34$; $\ \ket{\psi_1}$

\textbf{\color{brand}Concepts required.} Post-selection keeps only the runs with the desired measurement outcome; the kept state is the renormalised branch.

\textbf{\color{brand}Relevant formulae.} $P(1) = \lVert\Pi_1\ket\Psi\rVert^2$; $\ \ket{\Psi'} = \Pi_1\ket\Psi/\sqrt{P(1)}$.

\textbf{\color{brand}Step-by-step explanation.} $P(a = 1) = \left\lVert\frac{\sqrt3}{2}\ket{\psi_1}\right\rVert^2 = \frac34$. After the measurement only the $a = 1$ branch survives and is renormalised:
$\frac{\frac{\sqrt3}{2}\ket{\psi_1}}{\sqrt{3/4}} = \ket{\psi_1}$. HHL works in exactly this way: the ancilla-$1$ branch holds $\propto A^{-1}\ket b$, the ancilla-$0$
branch is discarded, and the probability of keeping it is the success probability.

\textbf{\color{brand}Why the other options are incorrect.}
\begin{itemize}[leftmargin=2.4em, itemsep=1pt, topsep=2pt]
\item[\textbf{(B)}] $\tfrac{\sqrt3}2$ is the amplitude of the branch, not its probability.
\item[\textbf{(C)}] $\tfrac14$ and $\ket{\psi_0}$ describe the other outcome ($a = 0$), which HHL discards.
\item[\textbf{(D)}] The $\ket{\psi_0}$ branch is removed by the measurement; no superposition of both branches survives.
\end{itemize}
\textbf{\color{brand}Key takeaway.} Post-selection probability $=$ squared norm of the kept branch; in HHL it is $P_{\text{success}} = C^2\lVert A^{-1}\ket b\rVert^2$.
\end{solbox}
\qsep

\needspace{13\baselineskip}
\begin{qitem}{35}
\qmeta{Single correct}{Inverse of a circuit}
A circuit applies $U_1$, then $U_2$, then $U_3$ (all unitary). The circuit that undoes it applies
\begin{itemize}[leftmargin=2.6em, labelsep=0.4em, itemsep=2pt, topsep=3pt]
\item[\optlabel{A}] $U_1^\dagger$, then $U_2^\dagger$, then $U_3^\dagger$
\item[\optlabel{B}] $U_3$, then $U_2$, then $U_1$
\item[\optlabel{C}] $U_1^\dagger U_2^\dagger U_3^\dagger$ applied as a single gate in the same order as the original
\item[\correctlabel{D}] $U_3^\dagger$, then $U_2^\dagger$, then $U_1^\dagger$
\end{itemize}
\end{qitem}
\begin{solbox}{D}
\textbf{\color{brand}Correct answer.} \textcolor{okgreen}{\textbf{(D)}} $U_3^\dagger$, then $U_2^\dagger$, then $U_1^\dagger$

\textbf{\color{brand}Concepts required.} The inverse of a product reverses the order: $(U_3U_2U_1)^{-1} = U_1^{-1}U_2^{-1}U_3^{-1}$.

\textbf{\color{brand}Relevant formulae.} $(AB)^\dagger = B^\dagger A^\dagger$; a circuit ``$U_1$ then $U_2$'' is the matrix $U_2U_1$.

\textbf{\color{brand}Step-by-step explanation.} The circuit implements the product $W = U_3U_2U_1$ (the first gate acts first, so it stands on the right).
$W^{-1} = W^\dagger = U_1^\dagger U_2^\dagger U_3^\dagger$; read right to left this means: apply $U_3^\dagger$ first, then $U_2^\dagger$, then $U_1^\dagger$.
Reverse the order \emph{and} invert every gate. HHL's uncomputation does exactly this with phase estimation: the QFT (inverse of the inverse QFT)
first, then the controlled $U^{-2^j}$ in reverse order, then the Hadamards.

\textbf{\color{brand}Why the other options are incorrect.}
\begin{itemize}[leftmargin=2.4em, itemsep=1pt, topsep=2pt]
\item[\textbf{(A)}] Inverting the gates without reversing the order gives $U_3^\dagger U_2^\dagger U_1^\dagger$, which is not $W^{-1}$ when the gates do not commute.
\item[\textbf{(B)}] Reversing the order without inverting the gates gives $U_1U_2U_3$, not the inverse.
\item[\textbf{(C)}] Combining the daggers into one gate in the original order is the same mistake as the first option.
\end{itemize}
\textbf{\color{brand}Key takeaway.} Undo = reverse order and dagger every gate. HHL's uncomputation is QPE run backwards.
\end{solbox}
\qsep

\needspace{13\baselineskip}
\begin{qitem}{36}
\qmeta{Matching lists}{Gate counts of a 4-qubit QFT}
The textbook QFT on $n = 4$ qubits (as used, inverted, on the clock register of the default HHL circuit). Match each quantity (List-I)
with its value (List-II).
\par\smallskip\noindent{\renewcommand{\arraystretch}{1.35}%
\begin{tabularx}{\linewidth}{@{}c>{\raggedright\arraybackslash}p{0.33\linewidth}@{\hspace{1.2em}}cL@{}}
\toprule
\multicolumn{2}{@{}l}{\textbf{List-I}} & \multicolumn{2}{l@{}}{\textbf{List-II}} \\ \midrule
\textbf{(P)} & number of Hadamard gates & \textbf{(1)} & $2$ \\
\textbf{(Q)} & number of controlled-phase gates & \textbf{(2)} & $\pi/8$ \\
\textbf{(R)} & number of SWAP gates & \textbf{(3)} & $6$ \\
\textbf{(S)} & magnitude of the smallest controlled-phase angle & \textbf{(4)} & $4$ \\
\bottomrule
\end{tabularx}}
\par\smallskip\noindent\begin{tabularx}{\linewidth}{@{}*{2}{L}@{}}
\optlabel{A}\,P$\to$3, Q$\to$4, R$\to$1, S$\to$2 & \optlabel{B}\,P$\to$4, Q$\to$1, R$\to$3, S$\to$2 \\[3pt]
\optlabel{C}\,P$\to$4, Q$\to$3, R$\to$2, S$\to$1 & \correctlabel{D}\,P$\to$4, Q$\to$3, R$\to$1, S$\to$2
\end{tabularx}
\end{qitem}
\begin{solbox}{D}
\textbf{\color{brand}Correct answer.} \textcolor{okgreen}{\textbf{(D)}} P$\to$4, Q$\to$3, R$\to$1, S$\to$2

\textbf{\color{brand}Concepts required.} The $n$-qubit QFT circuit: $n$ Hadamards, $n(n-1)/2$ controlled phases $R_k$, $\lfloor n/2\rfloor$ SWAPs.

\textbf{\color{brand}Relevant formulae.} $R_k = \operatorname{diag}(1, e^{2\pi i/2^k})$, $k = 2, \dots, n$; smallest angle $2\pi/2^n$.

\textbf{\color{brand}Step-by-step explanation.} \begin{itemize}\itemsep0pt
\item P: one Hadamard per qubit: $4$ (4).
\item Q: $n(n-1)/2 = 4\cdot3/2 = 6$ controlled phases (3).
\item R: $\lfloor n/2\rfloor = 2$ SWAPs to reverse the qubit order (1).
\item S: angles $2\pi/2^k$ for $k = 2, \dots, n$; the smallest is $2\pi/2^4 = \pi/8$ (2).
\end{itemize}
Consistency with \texttt{HHL.ipynb}: the recorded statistics show $12$ \texttt{cr1} and $4$ \texttt{swap} gates (inverse QFT + QFT for uncomputation:
$2\times6$ and $2\times2$) and $16$ \texttt{h} gates ($4$ initial + $4$ + $4$ in the two QFTs + $4$ final).

\textbf{\color{brand}Why the other options are incorrect.}
\begin{itemize}[leftmargin=2.4em, itemsep=1pt, topsep=2pt]
\item[\textbf{(A)}] P$\to$3, Q$\to$4: there is one Hadamard per qubit ($4$), and $6 = 4\cdot3/2$ controlled phases.
\item[\textbf{(B)}] Q$\to$1, R$\to$3 exchanges the counts: $6$ controlled phases and $2$ SWAPs, not $2$ and $6$.
\item[\textbf{(C)}] R$\to$2, S$\to$1 gives $\pi/8$ SWAPs and an angle of $2$ --- units mixed up.
\end{itemize}
\textbf{\color{brand}Key takeaway.} QFT cost is $O(n^2)$ gates; in HHL it is negligible compared with the controlled $e^{iAt2^j}$.
\end{solbox}
\qsep

\needspace{13\baselineskip}
\begin{qitem}{37}
\qmeta{Single correct}{Uniform superposition}
A Hadamard is applied to each of four qubits in $\ket{0000}$. The amplitude of $\ket{1011}$ is
\par\smallskip\noindent\begin{tabularx}{\linewidth}{@{}*{4}{L}@{}}
\optlabel{A}\,$\tfrac1{16}$ & \correctlabel{B}\,$\tfrac14$ & \optlabel{C}\,$-\tfrac14$ & \optlabel{D}\,$\tfrac{1}{\sqrt8}$
\end{tabularx}
\end{qitem}
\begin{solbox}{B}
\textbf{\color{brand}Correct answer.} \textcolor{okgreen}{\textbf{(B)}} $\tfrac14$

\textbf{\color{brand}Concepts required.} $H^{\otimes n}\ket{0\dots0}$ is the uniform superposition with all amplitudes $2^{-n/2}$.

\textbf{\color{brand}Relevant formulae.} $H^{\otimes n}\ket x = 2^{-n/2}\sum_y(-1)^{x\cdot y}\ket y$; for $x = 0$ all signs are $+$.

\textbf{\color{brand}Step-by-step explanation.} $H\ket0 = \frac{1}{\sqrt2}(\ket0 + \ket1)$ for each qubit, so $H^{\otimes4}\ket{0000} = \frac{1}{(\sqrt2)^4}\sum_{y=0}^{15}\ket y = \frac14\sum_y\ket y$.
Every one of the $16$ basis states has amplitude $\frac14$ (probability $\frac1{16}$), with a $+$ sign because the input is $\ket{0\dots0}$
(signs $(-1)^{x\cdot y}$ appear only for inputs $x \neq 0$). This is the first step of phase estimation on the HHL clock register.

\textbf{\color{brand}Why the other options are incorrect.}
\begin{itemize}[leftmargin=2.4em, itemsep=1pt, topsep=2pt]
\item[\textbf{(A)}] $\tfrac1{16}$ is the probability, not the amplitude.
\item[\textbf{(C)}] $-\tfrac14$: signs $(-1)^{x\cdot y}$ appear only for inputs $x \neq 0$.
\item[\textbf{(D)}] $\tfrac{1}{\sqrt8}$ corresponds to $3$ qubits.
\end{itemize}
\textbf{\color{brand}Key takeaway.} The clock register starts in the uniform superposition; QPE then writes phases into it.
\end{solbox}
\qsep

\needspace{13\baselineskip}
\begin{qitem}{38}
\qmeta{Statement I/II}{Product states, reduced states}
\textbf{Statement I:} Every two-qubit pure state can be written as $\ket a\otimes\ket b$ for suitable single-qubit states $\ket a$, $\ket b$.\\[1mm]
\textbf{Statement II:} For the Bell state $\frac{1}{\sqrt2}(\ket{00} + \ket{11})$ the reduced state of each qubit is the maximally mixed state $\frac12 I$.
\begin{itemize}[leftmargin=2.6em, labelsep=0.4em, itemsep=2pt, topsep=3pt]
\item[\optlabel{A}] Both Statement I and Statement II are true
\item[\optlabel{B}] Both Statement I and Statement II are false
\item[\optlabel{C}] Statement I is true, but Statement II is false
\item[\correctlabel{D}] Statement I is false, but Statement II is true
\end{itemize}
\end{qitem}
\begin{solbox}{D}
\textbf{\color{brand}Correct answer.} \textcolor{okgreen}{\textbf{(D)}} Statement I is false, but Statement II is true

\textbf{\color{brand}Concepts required.} Entangled states cannot be factorised; their subsystems are in mixed states.

\textbf{\color{brand}Relevant formulae.} $\rho_A = \operatorname{Tr}_B\ket\psi\!\bra\psi$; Bell state: $\rho_A = \frac12I$; product $\iff$ rank-1 coefficient matrix.

\textbf{\color{brand}Step-by-step explanation.} \textbf{I is false:} product states form a small subset; the Bell state is a counter-example (its coefficient matrix $\frac{1}{\sqrt2}I$ has rank 2).
\textbf{II is true:} tracing out the second qubit gives $\rho_A = \frac12(\ket0\!\bra0 + \ket1\!\bra1) = \frac12I$; a measurement of one qubit alone
gives $0$ or $1$ with probability $\frac12$ each.
In HHL, if the clock were \emph{not} uncomputed, the reduced state of the system register would be mixed like this instead of the pure $\ket x$.
Answer: (D).

\textbf{\color{brand}Why the other options are incorrect.}
\begin{itemize}[leftmargin=2.4em, itemsep=1pt, topsep=2pt]
\item[\textbf{(A)}] accepts I, but entangled states (e.g.\ Bell states) are not products.
\item[\textbf{(B)}] rejects II, but the reduced state of a Bell state is $\frac12I$.
\item[\textbf{(C)}] accepts I and rejects II --- both reversed.
\end{itemize}
\textbf{\color{brand}Key takeaway.} Leaving the clock entangled with the system would turn HHL's pure output $\ket x$ into a mixture.
\end{solbox}
\qsep

\needspace{20\baselineskip}
\begin{passagebox}{Paragraph for Questions 39 and 40}
A unitary $U$ has the eigenstate $\ket u$ with $U\ket u = e^{2\pi i\varphi}\ket u$, $\varphi = \tfrac14$. Phase estimation uses
$n_c = 2$ clock qubits $c_0c_1$ ($c_0$ = most significant bit). Both start in $\ket0$ and receive a Hadamard; $c_1$ (weight $2^0$) controls $U$
and $c_0$ (weight $2^1$) controls $U^2$, so that the state becomes $\frac12\sum_{k=0}^{3}e^{2\pi i\varphi k}\ket{k}\ket u$. Finally the inverse
QFT is applied to the clock and the clock is measured.
\end{passagebox}
\needspace{13\baselineskip}
\begin{qitem}{39}
\qmeta{Paragraph based}{QPE trace: phase state}
Before the inverse QFT, the amplitude of the clock state $\ket{11}$ (with $\ket u$ factored out) is
\par\smallskip\noindent\begin{tabularx}{\linewidth}{@{}*{4}{L}@{}}
\optlabel{A}\,$\tfrac12$ & \correctlabel{B}\,$-\tfrac i2$ & \optlabel{C}\,$-\tfrac12$ & \optlabel{D}\,$\tfrac i2$
\end{tabularx}
\end{qitem}
\begin{solbox}{B}
\textbf{\color{brand}Correct answer.} \textcolor{okgreen}{\textbf{(B)}} $-\tfrac i2$

\textbf{\color{brand}Concepts required.} After the controlled powers, clock value $k$ carries the phase $e^{2\pi i\varphi k}$.

\textbf{\color{brand}Relevant formulae.} $\frac{1}{\sqrt{2^{n_c}}}\sum_k e^{2\pi i\varphi k}\ket k$; $\ e^{2\pi i\cdot3/4} = e^{3\pi i/2} = -i$.

\textbf{\color{brand}Step-by-step explanation.} The clock amplitudes are $\frac12e^{2\pi i\varphi k}$ with $\varphi = \frac14$, i.e.\ $\frac12 i^k$. For $\ket{11}$, $k = 3$:
$\frac12 i^3 = \frac12(-i) = -\frac i2$. The four amplitudes are $\frac12(1, i, -1, -i)$: each controlled power adds the phase $2\pi\varphi\cdot 2^j$
to the clock qubit of weight $2^j$ (phase kickback, Q.26).

\textbf{\color{brand}Why the other options are incorrect.}
\begin{itemize}[leftmargin=2.4em, itemsep=1pt, topsep=2pt]
\item[\textbf{(A)}] $\tfrac12$ is the amplitude of $\ket{00}$; it ignores the phase.
\item[\textbf{(C)}] $-\tfrac12$ is the amplitude of $\ket{10}$ ($k = 2$).
\item[\textbf{(D)}] $\tfrac i2$ is the amplitude of $\ket{01}$ ($k = 1$), not of $\ket{11}$.
\end{itemize}
\textbf{\color{brand}Key takeaway.} Clock value $k$ accumulates $k$ times the eigenphase --- a phase gradient, i.e.\ a Fourier state.
\end{solbox}
\qsep

\needspace{13\baselineskip}
\begin{qitem}{40}
\qmeta{Paragraph based}{QPE trace: outcome}
After the inverse QFT the clock measurement gives
\par\smallskip\noindent\begin{tabularx}{\linewidth}{@{}*{2}{L}@{}}
\correctlabel{A}\,$01$ with probability $1$, so $\tilde\varphi = \tfrac14$ & \optlabel{B}\,$01$ or $11$ with probability $\tfrac12$ each \\[3pt]
\optlabel{C}\,$11$ with probability $1$, so $\tilde\varphi = \tfrac34$ & \optlabel{D}\,$10$ with probability $1$, so $\tilde\varphi = \tfrac12$
\end{tabularx}
\end{qitem}
\begin{solbox}{A}
\textbf{\color{brand}Correct answer.} \textcolor{okgreen}{\textbf{(A)}} $01$ with probability $1$, so $\tilde\varphi = \tfrac14$

\textbf{\color{brand}Concepts required.} If $\varphi2^{n_c}$ is an integer $k$, the clock is exactly $\mathrm{QFT}\ket k$ and the inverse QFT returns $\ket k$.

\textbf{\color{brand}Relevant formulae.} $\frac12\sum_y e^{2\pi i\,ky/4}\ket y = \mathrm{QFT}\ket k$; $\ \tilde\varphi = k/2^{n_c}$.

\textbf{\color{brand}Step-by-step explanation.} The clock state $\frac12\sum_k e^{2\pi i\,k/4}\ket k$ is $\mathrm{QFT}\ket1$ (definition with $N = 4$, $x = 1$). Hence the inverse QFT gives $\ket1 = \ket{01}$
with certainty, and $\tilde\varphi = 1/2^{n_c} = 1/4$ --- exact, because $\varphi\cdot2^{n_c} = 1$ is an integer. In HHL the same mechanism writes
$k = \lambda t2^{n_c}/2\pi$ into the clock, one value per eigencomponent of $\ket b$.

\textbf{\color{brand}Why the other options are incorrect.}
\begin{itemize}[leftmargin=2.4em, itemsep=1pt, topsep=2pt]
\item[\textbf{(B)}] An exactly representable phase gives a deterministic outcome, not a $50/50$ split.
\item[\textbf{(C)}] $11$ ($\tilde\varphi = \tfrac34$) corresponds to the conjugate phase $e^{-2\pi i/4}$.
\item[\textbf{(D)}] $10$ ($\tilde\varphi = \tfrac12$) would need the amplitudes $\frac12(1, -1, 1, -1)$.
\end{itemize}
\textbf{\color{brand}Key takeaway.} Exact phase, exact readout: this is why the default HHL example (integer clock values) has fidelity $1$.
\end{solbox}
\qsep

\newpage

\begin{sectionbanner}SECTION C \enspace$\cdot$\enspace HHL Algorithm Fundamentals \hfill Q.41--60\end{sectionbanner}

\needspace{13\baselineskip}
\begin{qitem}{41}
\qmeta{Single correct}{What HHL outputs}
For a Hermitian $A$ and a prepared $\ket b$, a successful run of the HHL algorithm produces
\begin{itemize}[leftmargin=2.6em, labelsep=0.4em, itemsep=2pt, topsep=3pt]
\item[\optlabel{A}] the eigenvalues of $A$ as a classical list
\item[\optlabel{B}] the classical list of all $d$ numbers $x_1, \dots, x_d$
\item[\correctlabel{C}] the quantum state $\ket x \propto A^{-1}\ket b$ in the system register
\item[\optlabel{D}] the matrix $A^{-1}$ stored in the clock register
\end{itemize}
\end{qitem}
\begin{solbox}{C}
\textbf{\color{brand}Correct answer.} \textcolor{okgreen}{\textbf{(C)}} the quantum state $\ket x \propto A^{-1}\ket b$ in the system register

\textbf{\color{brand}Concepts required.} HHL is a state-preparation algorithm: its output is a quantum state whose amplitudes are proportional to $\vec x$.

\textbf{\color{brand}Relevant formulae.} $\ket x = \frac{A^{-1}\ket b}{\lVert A^{-1}\ket b\rVert} = \frac{\sum_j(\beta_j/\lambda_j)\ket{u_j}}{\sqrt{\sum_j|\beta_j|^2/\lambda_j^2}}$.

\textbf{\color{brand}Step-by-step explanation.} HHL transforms $\ket b = \sum_j\beta_j\ket{u_j}$ into $\ket x = \frac{\sum_j(\beta_j/\lambda_j)\ket{u_j}}{\lVert\cdot\rVert}$, a normalised quantum state
whose amplitudes are proportional to the solution vector. The $d$ numbers are \emph{encoded} in $\log_2d$ qubits; reading them all out
needs $\Omega(d)$ repetitions (Section D), so HHL is useful when only a property of $\vec x$ (an expectation value, an overlap) is needed.
The eigenvalues are used only temporarily inside the circuit and are uncomputed; $A^{-1}$ is never stored.

\textbf{\color{brand}Why the other options are incorrect.}
\begin{itemize}[leftmargin=2.4em, itemsep=1pt, topsep=2pt]
\item[\textbf{(A)}] The eigenvalues are written into the clock only coherently and are erased again; they are not measured or output.
\item[\textbf{(B)}] A classical list of $d$ numbers would need $\Omega(d)$ measurements; one run yields one sample.
\item[\textbf{(D)}] $A^{-1}$ ($d^2$ numbers) is never formed; the clock holds eigenvalue estimates temporarily and is uncomputed.
\end{itemize}
\textbf{\color{brand}Key takeaway.} HHL outputs $\ket x$, not $\vec x$. Everything else about its usefulness follows from this.
\end{solbox}
\qsep

\needspace{13\baselineskip}
\begin{qitem}{42}
\qmeta{Single correct}{State after phase estimation}
With $\ket b = \sum_j\beta_j\ket{u_j}$ and exactly representable clock values $k_j$, the state after phase estimation (clock $\otimes$ system $\otimes$ ancilla) is
\par\smallskip\noindent\begin{tabularx}{\linewidth}{@{}*{2}{L}@{}}
\optlabel{A}\,$\ket{0\dots0}\ket x\ket1$ & \optlabel{B}\,$\left(\frac{1}{\sqrt{d}}\sum_j\ket{k_j}\right)\otimes\ket b\otimes\ket0$ \\[3pt]
\optlabel{C}\,$\sum_j\frac{\beta_j}{\lambda_j}\ket{k_j}\ket{u_j}\ket0$ & \correctlabel{D}\,$\sum_j\beta_j\ket{k_j}\ket{u_j}\ket0$
\end{tabularx}
\end{qitem}
\begin{solbox}{D}
\textbf{\color{brand}Correct answer.} \textcolor{okgreen}{\textbf{(D)}} $\sum_j\beta_j\ket{k_j}\ket{u_j}\ket0$

\textbf{\color{brand}Concepts required.} Phase estimation acts eigencomponent by eigencomponent (linearity) and writes each eigenvalue into the clock.

\textbf{\color{brand}Relevant formulae.} $\ket0\ket{u_j} \xrightarrow{\text{QPE}} \ket{k_j}\ket{u_j}$, $\ k_j = \lambda_jt2^{n_c}/2\pi$; $\ \ket b = \sum_j\beta_j\ket{u_j}$.

\textbf{\color{brand}Step-by-step explanation.} Phase estimation is linear. On each eigencomponent $\ket{u_j}$, $U = e^{iAt}$ acts as the phase $e^{i\lambda_jt} = e^{2\pi i\,\varphi_j}$ with
$\varphi_j = \lambda_jt/2\pi$, and QPE writes $k_j = \varphi_j2^{n_c}$ into the clock: $\ket0\ket{u_j} \mapsto \ket{k_j}\ket{u_j}$. Summing with weights $\beta_j$:
\[ \ket{0}\ket b\ket0 = \sum_j\beta_j\ket0\ket{u_j}\ket0 \;\longmapsto\; \sum_j\beta_j\ket{k_j}\ket{u_j}\ket0 . \]
The clock is now \emph{entangled} with the system; no factor $1/\lambda_j$ has appeared yet --- that is the job of the controlled rotation.

\textbf{\color{brand}Why the other options are incorrect.}
\begin{itemize}[leftmargin=2.4em, itemsep=1pt, topsep=2pt]
\item[\textbf{(A)}] $\ket{0\dots0}\ket x\ket1$ is the (ideal) state at the very end, after rotation, uncomputation and post-selection.
\item[\textbf{(B)}] A product state ignores that each eigenvalue is correlated with its own eigenvector: QPE entangles clock and system.
\item[\textbf{(C)}] The factors $1/\lambda_j$ appear only after the controlled rotation and post-selection, not after QPE.
\end{itemize}
\textbf{\color{brand}Key takeaway.} After QPE: $\sum_j\beta_j\ket{k_j}\ket{u_j}$ --- eigenvalues sit next to their eigenvectors, ready to be inverted.
\end{solbox}
\qsep

\needspace{13\baselineskip}
\begin{qitem}{43}
\qmeta{Single correct}{Eigenvalue $\to$ clock bit string}
In the default notebook run ($n_c = 4$, $t = \pi/8$), the eigenvalue $\lambda = 3$ is written into the clock register as the bit string
\par\smallskip\noindent\begin{tabularx}{\linewidth}{@{}*{4}{L}@{}}
\optlabel{A}\,$0010$ & \optlabel{B}\,$0110$ & \correctlabel{C}\,$0011$ & \optlabel{D}\,$1100$
\end{tabularx}
\end{qitem}
\begin{solbox}{C}
\textbf{\color{brand}Correct answer.} \textcolor{okgreen}{\textbf{(C)}} $0011$

\textbf{\color{brand}Concepts required.} The clock value is the eigenphase in units of $2\pi/2^{n_c}$, written in binary with the most significant bit first.

\textbf{\color{brand}Relevant formulae.} $k = \dfrac{\lambda t2^{n_c}}{2\pi}$; default $t = \pi/8$, $n_c = 4 \Rightarrow k = \lambda$.

\textbf{\color{brand}Step-by-step explanation.} $k = \dfrac{\lambda t\,2^{n_c}}{2\pi} = \dfrac{3\cdot\frac\pi8\cdot16}{2\pi} = \dfrac{6\pi}{2\pi} = 3$. With $n_c = 4$ bits and the most significant bit first, $3 = 0011_2$.
In fact $t = \pi/8$ was chosen (``auto'') so that $k = \lambda$ for every eigenvalue: $1, 2, 3, 4 \to 0001, 0010, 0011, 0100$, all exact, which is
why the recorded fidelity is $1$. $1100$ is the bit-reversed string; $0110 = 6$ comes from dividing by $\pi$ instead of $2\pi$.

\textbf{\color{brand}Why the other options are incorrect.}
\begin{itemize}[leftmargin=2.4em, itemsep=1pt, topsep=2pt]
\item[\textbf{(A)}] $0010 = 2$ is the clock value of $\lambda = 2$.
\item[\textbf{(B)}] $0110 = 6$ divides by $\pi$ instead of $2\pi$.
\item[\textbf{(D)}] $1100$ is $0011$ written in reverse bit order (least significant bit first).
\end{itemize}
\textbf{\color{brand}Key takeaway.} Choose $t$ so that every $\lambda t2^{n_c}/2\pi$ is an integer and QPE becomes exact.
\end{solbox}
\qsep

\needspace{13\baselineskip}
\begin{qitem}{44}
\qmeta{Single correct}{Constant $C$ and rotation angle}
For the default notebook run ($n_c = 4$, $t = \pi/8$) the constant $C$ and the rotation angle $\theta$ applied for the clock value of $\lambda = 2$ are
\par\smallskip\noindent\begin{tabularx}{\linewidth}{@{}*{2}{L}@{}}
\optlabel{A}\,$C = \tfrac1{16}$; $\ \theta = 2\arcsin\tfrac1{32}$ & \correctlabel{B}\,$C = 1$; $\ \theta = \pi/3$ \\[3pt]
\optlabel{C}\,$C = 1$; $\ \theta = \pi/6$ & \optlabel{D}\,$C = \pi/8$; $\ \theta = 2\arcsin(\pi/16)$
\end{tabularx}
\end{qitem}
\begin{solbox}{B}
\textbf{\color{brand}Correct answer.} \textcolor{okgreen}{\textbf{(B)}} $C = 1$; $\ \theta = \pi/3$

\textbf{\color{brand}Concepts required.} $C$ is the eigenvalue represented by clock value $1$; the rotation angle loads the amplitude $C/\tilde\lambda$.

\textbf{\color{brand}Relevant formulae.} $C = \dfrac{2\pi}{2^{n_c}t}$, $\ \tilde\lambda_k = kC$, $\ \theta_k = 2\arcsin\dfrac{C}{\tilde\lambda_k}$.

\textbf{\color{brand}Step-by-step explanation.} $C = \dfrac{2\pi}{2^{n_c}t} = \dfrac{2\pi}{16\cdot\pi/8} = \dfrac{2\pi}{2\pi} = 1$ --- the eigenvalue that maps to clock value $1$.
For $\lambda = 2$: $\theta = 2\arcsin\frac{C}{\lambda} = 2\arcsin\frac12 = 2\cdot\frac\pi6 = \frac\pi3$, so the ancilla becomes
$\cos\frac\pi6\ket0 + \sin\frac\pi6\ket1 = \frac{\sqrt3}{2}\ket0 + \frac12\ket1$: amplitude $C/\lambda = \frac12$ on $\ket1$. (Note $R_Y(\theta)$ rotates by $\theta/2$ in amplitude, hence the factor $2$.)

\textbf{\color{brand}Why the other options are incorrect.}
\begin{itemize}[leftmargin=2.4em, itemsep=1pt, topsep=2pt]
\item[\textbf{(A)}] $C = \frac1{16}$ uses $1/2^{n_c}$ and forgets the factor $2\pi/t$.
\item[\textbf{(C)}] $\theta = \pi/6$ forgets the factor $2$: $\arcsin\frac12 = \frac\pi6$ is the half-angle.
\item[\textbf{(D)}] $C = \pi/8$ confuses $C$ with the evolution time $t$.
\end{itemize}
\textbf{\color{brand}Key takeaway.} Default run: $C = 1$, so the ancilla amplitude for eigenvalue $\lambda$ is simply $1/\lambda$.
\end{solbox}
\qsep

\needspace{13\baselineskip}
\begin{qitem}{45}
\qmeta{Single correct}{Ancilla amplitude $C/\lambda$}
With $C = 1$, the eigencomponent with $\lambda = 4$ leaves the ancilla in $\sqrt{1 - C^2/\lambda^2}\,\ket0 + \frac C\lambda\ket1$. For this component the amplitude
on $\ket1$ and the probability of measuring the ancilla as $1$ are
\par\smallskip\noindent\begin{tabularx}{\linewidth}{@{}*{4}{L}@{}}
\optlabel{A}\,$4$ and $1$ & \optlabel{B}\,$\tfrac14$ and $\tfrac14$ & \optlabel{C}\,$\tfrac12$ and $\tfrac14$ & \correctlabel{D}\,$\tfrac14$ and $\tfrac1{16}$
\end{tabularx}
\end{qitem}
\begin{solbox}{D}
\textbf{\color{brand}Correct answer.} \textcolor{okgreen}{\textbf{(D)}} $\tfrac14$ and $\tfrac1{16}$

\textbf{\color{brand}Concepts required.} Eigenvalue inversion: the ancilla amplitude on $\ket1$ is $C/\lambda$; its square is the success probability of that eigencomponent.

\textbf{\color{brand}Relevant formulae.} $\ket0_a \to \sqrt{1 - \frac{C^2}{\lambda^2}}\ket0 + \frac C\lambda\ket1$; $\ P_j = C^2/\lambda_j^2$.

\textbf{\color{brand}Step-by-step explanation.} Amplitude $C/\lambda = 1/4$; probability $|1/4|^2 = 1/16 = 0.0625$. The $\ket0$ amplitude is $\sqrt{15/16} = 0.968$ so the ancilla is normalised.
Large eigenvalues are therefore strongly suppressed, exactly as $1/\lambda$ requires; the price is that components with large $\lambda$ rarely
pass post-selection. ($\frac12$ would correspond to $1/\sqrt\lambda$, and an amplitude $4$ is impossible --- that is why $C \le |\lambda|_{\min}$ is needed.)

\textbf{\color{brand}Why the other options are incorrect.}
\begin{itemize}[leftmargin=2.4em, itemsep=1pt, topsep=2pt]
\item[\textbf{(A)}] An amplitude of $4$ is impossible ($|\cdot| \le 1$); it inverts $C/\lambda$.
\item[\textbf{(B)}] $\tfrac14$ for the probability forgets to square the amplitude.
\item[\textbf{(C)}] $\tfrac12$ is $1/\sqrt\lambda$: the rotation implements $1/\lambda$, not $1/\sqrt\lambda$.
\end{itemize}
\textbf{\color{brand}Key takeaway.} Amplitude $C/\lambda$, probability $C^2/\lambda^2$: large eigenvalues are suppressed, small ones dominate $\ket x$.
\end{solbox}
\qsep

\needspace{13\baselineskip}
\begin{qitem}{46}
\qmeta{Matching lists}{Roles of the registers}
Match each part of the HHL circuit (List-I) with its role (List-II).
\par\smallskip\noindent{\renewcommand{\arraystretch}{1.35}%
\begin{tabularx}{\linewidth}{@{}c>{\raggedright\arraybackslash}p{0.33\linewidth}@{\hspace{1.2em}}cL@{}}
\toprule
\multicolumn{2}{@{}l}{\textbf{List-I}} & \multicolumn{2}{l@{}}{\textbf{List-II}} \\ \midrule
\textbf{(P)} & clock register ($n_c$ qubits) & \textbf{(1)} & is measured; outcome $1$ flags a successful inversion \\
\textbf{(Q)} & system register ($n_b$ qubits) & \textbf{(2)} & holds $\ket b$ at the start and $\ket x$ after success \\
\textbf{(R)} & ancilla qubit & \textbf{(3)} & holds the binary eigenvalue estimates $k_j$ \\
\textbf{(S)} & inverse phase estimation after the rotation & \textbf{(4)} & returns the clock to $\ket{0\dots0}$ and removes its entanglement with the system \\
\bottomrule
\end{tabularx}}
\par\smallskip\noindent\begin{tabularx}{\linewidth}{@{}*{2}{L}@{}}
\correctlabel{A}\,P$\to$3, Q$\to$2, R$\to$1, S$\to$4 & \optlabel{B}\,P$\to$2, Q$\to$3, R$\to$1, S$\to$4 \\[3pt]
\optlabel{C}\,P$\to$3, Q$\to$2, R$\to$4, S$\to$1 & \optlabel{D}\,P$\to$3, Q$\to$1, R$\to$2, S$\to$4
\end{tabularx}
\end{qitem}
\begin{solbox}{A}
\textbf{\color{brand}Correct answer.} \textcolor{okgreen}{\textbf{(A)}} P$\to$3, Q$\to$2, R$\to$1, S$\to$4

\textbf{\color{brand}Concepts required.} HHL uses three registers: clock (eigenvalues), system ($\ket b \to \ket x$) and one ancilla (success flag), plus uncomputation.

\textbf{\color{brand}Relevant formulae.} qubits $= n_c + n_b + 1$; $\ \ket0^{\otimes n_c}\ket b\ket0 \to \ket0^{\otimes n_c}\big(\sqrt{\cdot}\,\ket{\text{junk}}\ket0 + C A^{-1}\ket b\ket1\big)$.

\textbf{\color{brand}Step-by-step explanation.} \begin{itemize}\itemsep0pt
\item P: QPE writes $k_j = \lambda_jt2^{n_c}/2\pi$ into the clock (3).
\item Q: the system starts in $\ket b$ (amplitude encoding) and ends in $\ket x$ when the run succeeds (2).
\item R: the ancilla carries the amplitude $C/\tilde\lambda$ on $\ket1$; measuring $1$ post-selects the $A^{-1}$ branch (1).
\item S: uncomputation maps $\ket{k_j}\ket{u_j} \to \ket{0}\ket{u_j}$ so that the eigencomponents can interfere into $\ket x$ (4).
\end{itemize}
Total qubits: $n_c + n_b + 1$ ($4 + 2 + 1 = 7$ in the notebook).

\textbf{\color{brand}Why the other options are incorrect.}
\begin{itemize}[leftmargin=2.4em, itemsep=1pt, topsep=2pt]
\item[\textbf{(B)}] P$\to$2, Q$\to$3 swaps clock and system: the eigenvalues go into the clock, $\ket b$ into the system.
\item[\textbf{(C)}] R$\to$4, S$\to$1 swaps the ancilla's role with uncomputation; the ancilla is the qubit that is measured.
\item[\textbf{(D)}] Q$\to$1, R$\to$2 makes the system the success flag and the ancilla the solution register.
\end{itemize}
\textbf{\color{brand}Key takeaway.} Clock = eigenvalues, system = data, ancilla = flag; uncomputation cleans up the clock.
\end{solbox}
\qsep

\needspace{13\baselineskip}
\begin{qitem}{47}
\qmeta{Single correct}{Success probability}
HHL is run with $A = \operatorname{diag}(1, 2)$, $\ket b = \frac{1}{\sqrt2}(\ket0 + \ket1)$, $C = 1$ and exactly representable eigenvalues. The success probability is
\par\smallskip\noindent\begin{tabularx}{\linewidth}{@{}*{4}{L}@{}}
\optlabel{A}\,$0.75$ & \optlabel{B}\,$1$ & \optlabel{C}\,$0.5$ & \correctlabel{D}\,$0.625$
\end{tabularx}
\end{qitem}
\begin{solbox}{D}
\textbf{\color{brand}Correct answer.} \textcolor{okgreen}{\textbf{(D)}} $0.625$

\textbf{\color{brand}Concepts required.} The success probability is the squared norm of the post-selected branch $C\,A^{-1}\ket b$.

\textbf{\color{brand}Relevant formulae.} $P_{\text{success}} = \sum_j|\beta_j|^2\frac{C^2}{\lambda_j^2} = C^2\lVert A^{-1}\ket b\rVert^2$.

\textbf{\color{brand}Step-by-step explanation.} The eigenvectors are $\ket0, \ket1$ with $\lambda = 1, 2$ and $\beta_1 = \beta_2 = \frac{1}{\sqrt2}$.
\[ P_{\text{success}} = \sum_j|\beta_j|^2\frac{C^2}{\lambda_j^2} = \frac12\cdot\frac11 + \frac12\cdot\frac14 = \frac58 = 0.625. \]
Equivalently $P = C^2\lVert A^{-1}\ket b\rVert^2$ with $A^{-1}\ket b = \frac{1}{\sqrt2}(1, \frac12)$. ($0.75$ forgets to square $C/\lambda$.) The post-selected
state is $\propto \ket0 + \frac12\ket1$, i.e.\ $\ket x = \frac{1}{\sqrt5}(2\ket0 + \ket1)$.

\textbf{\color{brand}Why the other options are incorrect.}
\begin{itemize}[leftmargin=2.4em, itemsep=1pt, topsep=2pt]
\item[\textbf{(A)}] $0.75 = \frac12(1 + \frac12)$ uses $C/\lambda$ instead of $(C/\lambda)^2$.
\item[\textbf{(B)}] $1$ would need $C/\lambda_j = 1$ for every component, i.e.\ all eigenvalues equal to $C$.
\item[\textbf{(C)}] $0.5$ keeps only the $\lambda = 1$ component.
\end{itemize}
\textbf{\color{brand}Key takeaway.} $P_{\text{success}} \propto C^2/\lambda^2$: it decreases with the eigenvalues relative to $C$, typically like $1/\kappa^2$.
\end{solbox}
\qsep

\needspace{13\baselineskip}
\begin{qitem}{48}
\qmeta{Single correct}{Qubit count $n_b + n_c + 1$}
An $8\times8$ Hermitian system is solved with the notebook's circuit and $n_c = 5$ clock qubits. The total number of qubits is
\par\smallskip\noindent\begin{tabularx}{\linewidth}{@{}*{4}{L}@{}}
\optlabel{A}\,$6$ & \correctlabel{B}\,$9$ & \optlabel{C}\,$8$ & \optlabel{D}\,$14$
\end{tabularx}
\end{qitem}
\begin{solbox}{B}
\textbf{\color{brand}Correct answer.} \textcolor{okgreen}{\textbf{(B)}} $9$

\textbf{\color{brand}Concepts required.} Qubit counts are logarithmic in the matrix dimension.

\textbf{\color{brand}Relevant formulae.} $n_b = \log_2 d$; total $= n_c + n_b + 1$.

\textbf{\color{brand}Step-by-step explanation.} $n_b = \log_2 8 = 3$ system qubits, $n_c = 5$ clock qubits and $1$ ancilla: $3 + 5 + 1 = 9$. The matrix dimension enters only
logarithmically --- the source of HHL's potential advantage. ($14 = 8 + 5 + 1$ confuses the dimension $d$ with the number of qubits $\log_2d$.)

\textbf{\color{brand}Why the other options are incorrect.}
\begin{itemize}[leftmargin=2.4em, itemsep=1pt, topsep=2pt]
\item[\textbf{(A)}] $6 = 3 + 2 + 1$ would correspond to only $n_c = 2$ clock qubits; here $n_c = 5$.
\item[\textbf{(C)}] $8$ forgets the ancilla ($3 + 5 = 8$).
\item[\textbf{(D)}] $14$ uses the dimension $d = 8$ instead of $\log_2 d = 3$.
\end{itemize}
\textbf{\color{brand}Key takeaway.} Qubits $= n_c + \log_2 d + 1$; the clock, not the data, usually dominates small examples.
\end{solbox}
\qsep

\needspace{13\baselineskip}
\begin{qitem}{49}
\qmeta{Statement I/II}{Why uncompute?}
\textbf{Statement I:} If the clock register is not uncomputed, the post-selected system register is in general not in the pure state $\ket x$,
because it remains entangled with the clock.\\[1mm]
\textbf{Statement II:} Uncomputation can be replaced by simply measuring the clock register, with the same output $\ket x$.
\begin{itemize}[leftmargin=2.6em, labelsep=0.4em, itemsep=2pt, topsep=3pt]
\item[\optlabel{A}] Both Statement I and Statement II are true
\item[\optlabel{B}] Both Statement I and Statement II are false
\item[\correctlabel{C}] Statement I is true, but Statement II is false
\item[\optlabel{D}] Statement I is false, but Statement II is true
\end{itemize}
\end{qitem}
\begin{solbox}{C}
\textbf{\color{brand}Correct answer.} \textcolor{okgreen}{\textbf{(C)}} Statement I is true, but Statement II is false

\textbf{\color{brand}Concepts required.} Uncomputation coherently erases the clock; measurement would collapse the eigencomponent superposition.

\textbf{\color{brand}Relevant formulae.} $\sum_j\frac{\beta_j}{\lambda_j}\ket{k_j}\ket{u_j} \xrightarrow{\text{QPE}^{-1}} \ket0\sum_j\frac{\beta_j}{\lambda_j}\ket{u_j}$.

\textbf{\color{brand}Step-by-step explanation.} After the rotation and post-selection the state is $\propto\sum_j\frac{\beta_j}{\lambda_j}\ket{k_j}\ket{u_j}$.
\textbf{I is true:} with distinct $\ket{k_j}$ the clock labels the eigencomponents; tracing it out leaves the mixture $\sum_j p_j\ket{u_j}\!\bra{u_j}$,
not the superposition $\ket x$.
\textbf{II is false:} measuring the clock gives one value $k_j$ and collapses the system to the single eigenvector $\ket{u_j}$ --- the
interference between eigencomponents that forms $\ket x$ is destroyed. Uncomputation \emph{coherently} resets $\ket{k_j} \to \ket0$. Answer: (C).

\textbf{\color{brand}Why the other options are incorrect.}
\begin{itemize}[leftmargin=2.4em, itemsep=1pt, topsep=2pt]
\item[\textbf{(A)}] accepts II, but measuring the clock leaves a single eigenvector $\ket{u_j}$, not $\ket x$.
\item[\textbf{(B)}] rejects I, but an un-erased clock labels the components and the system's reduced state is mixed.
\item[\textbf{(D)}] rejects I and accepts II --- both reversed.
\end{itemize}
\textbf{\color{brand}Key takeaway.} Erase, don't measure: the eigenvalue register must be uncomputed for the components to interfere into $\ket x$.
\end{solbox}
\qsep

\needspace{13\baselineskip}
\begin{qitem}{50}
\qmeta{Assertion--Reason}{Choice of $C$}
\textbf{Assertion (A):} In HHL the constant $C$ must satisfy $C \le \min_j|\lambda_j|$ (more precisely, $\le$ the smallest $|\tilde\lambda|$ that can occur).\\[1mm]
\textbf{Reason (R):} The controlled rotation puts the amplitude $C/\tilde\lambda$ on the ancilla state $\ket1$, and an amplitude cannot exceed $1$ in magnitude.
\begin{itemize}[leftmargin=2.6em, labelsep=0.4em, itemsep=2pt, topsep=3pt]
\item[\correctlabel{A}] Both Assertion (A) and Reason (R) are true, and R is the correct explanation of A
\item[\optlabel{B}] Both A and R are true, but R is \emph{not} the correct explanation of A
\item[\optlabel{C}] A is true, but R is false
\item[\optlabel{D}] A is false, but R is true
\end{itemize}
\end{qitem}
\begin{solbox}{A}
\textbf{\color{brand}Correct answer.} \textcolor{okgreen}{\textbf{(A)}} Both Assertion (A) and Reason (R) are true, and R is the correct explanation of A

\textbf{\color{brand}Concepts required.} The controlled rotation needs $|\sin\frac\theta2| = |C/\tilde\lambda| \le 1$; $C$ sets the overall scale of the success amplitude.

\textbf{\color{brand}Relevant formulae.} $\sin\frac{\theta}{2} = \frac{C}{\tilde\lambda}$, $\ |C| \le \min|\tilde\lambda|$; $\ P_{\text{success}} \propto C^2$.

\textbf{\color{brand}Step-by-step explanation.} \textbf{R is true:} $R_Y(\theta)\ket0 = \cos\frac\theta2\ket0 + \sin\frac\theta2\ket1$ with $\sin\frac\theta2 = C/\tilde\lambda$, which requires $|C/\tilde\lambda| \le 1$.
\textbf{A is true} and follows from R: for every eigenvalue that occurs, $|C/\tilde\lambda| \le 1 \iff C \le |\tilde\lambda|$, so $C$ is bounded by the smallest one.
(The notebook sets $C = 2\pi/(2^{n_c}t)$, the smallest non-zero representable value, and clips the argument of $\arcsin$ to $[-1, 1]$.)
Choosing $C$ as large as allowed maximises the success probability $\propto C^2$. Answer: (A).

\textbf{\color{brand}Why the other options are incorrect.}
\begin{itemize}[leftmargin=2.4em, itemsep=1pt, topsep=2pt]
\item[\textbf{(B)}] denies the link, but the bound on $C$ is exactly the requirement that the amplitude $C/\tilde\lambda$ be a valid sine.
\item[\textbf{(C)}] rejects R, but $R_Y$ puts $\sin\frac\theta2$ on $\ket1$, which is bounded by $1$.
\item[\textbf{(D)}] rejects A; a larger $C$ would make $\arcsin(C/\tilde\lambda)$ undefined for the smallest eigenvalue.
\end{itemize}
\textbf{\color{brand}Key takeaway.} Choose $C$ as large as allowed, i.e.\ about $\lambda_{\min}$; then $P_{\text{success}} \gtrsim 1/\kappa^2$.
\end{solbox}
\qsep

\needspace{13\baselineskip}
\begin{qitem}{51}
\qmeta{Single correct}{Negative eigenvalues (two's complement)}
In the default notebook circuit ($n_c = 4$, $t = \pi/8$, $C = 1$, two's-complement clock), the clock value $k = 14$ ($1110$) is treated as the eigenvalue
$\tilde\lambda$ and rotated by $\theta$, where
\par\smallskip\noindent\begin{tabularx}{\linewidth}{@{}*{2}{L}@{}}
\optlabel{A}\,$\tilde\lambda = -2$, $\ \theta = +\pi/3$ & \optlabel{B}\,$\tilde\lambda = 2$, $\ \theta = \pi/3$ \\[3pt]
\correctlabel{C}\,$\tilde\lambda = -2$, $\ \theta = -\pi/3$ & \optlabel{D}\,$\tilde\lambda = 14$, $\ \theta = 2\arcsin\tfrac1{14}$
\end{tabularx}
\end{qitem}
\begin{solbox}{C}
\textbf{\color{brand}Correct answer.} \textcolor{okgreen}{\textbf{(C)}} $\tilde\lambda = -2$, $\ \theta = -\pi/3$

\textbf{\color{brand}Concepts required.} Negative eigenvalues produce phases $e^{i\lambda t}$ with $\lambda t < 0$, which wrap around to clock values near $2^{n_c}$ (two's complement).

\textbf{\color{brand}Relevant formulae.} signed $k = k - 2^{n_c}$ for $k \ge 2^{n_c-1}$; $\ \tilde\lambda = kC$; $\ \theta = 2\arcsin(C/\tilde\lambda)$, an odd function.

\textbf{\color{brand}Step-by-step explanation.} $k = 14 \ge 2^{n_c-1} = 8$, so it represents $14 - 16 = -2$ (two's complement), and $\tilde\lambda = -2\,C = -2$.
(Physically: $e^{-2it} = e^{i(2\pi - 2t)}$, so a negative eigenvalue produces the phase of $k = 16 - 2 = 14$.)
$\theta = 2\arcsin\frac{C}{\tilde\lambda} = 2\arcsin(-\tfrac12) = -\frac\pi3$: the ancilla amplitude is $-\frac12$. The \emph{sign} matters --- it is the
sign of $1/\lambda$ in $A^{-1}$; $+\pi/3$ would solve the system with $|A|$ instead of $A$.

\textbf{\color{brand}Why the other options are incorrect.}
\begin{itemize}[leftmargin=2.4em, itemsep=1pt, topsep=2pt]
\item[\textbf{(A)}] $\theta = +\pi/3$ loses the sign: $\arcsin$ is odd, and the sign of $1/\lambda$ must survive.
\item[\textbf{(B)}] $\tilde\lambda = 2$ drops the sign entirely, i.e.\ inverts $|A|$ instead of $A$.
\item[\textbf{(D)}] $\tilde\lambda = 14$ reads the clock as unsigned; with $t = \pi/8$ an eigenvalue $14$ would alias, and the notebook interprets $k \ge 8$ as negative.
\end{itemize}
\textbf{\color{brand}Key takeaway.} Signs of eigenvalues are kept by two's complement and by the sign of the rotation angle.
\end{solbox}
\qsep

\needspace{13\baselineskip}
\begin{qitem}{52}
\qmeta{Single correct}{Code: $X$ gates of the inversion}
The eigenvalue-inversion loop of \texttt{HHL.ipynb} runs over $k = 1, \dots, 2^{n_c}-1$; for each $k$ it applies \texttt{kernel.x} to every clock qubit whose bit
of $k$ is $0$, then the multi-controlled \texttt{cry}, then the same \texttt{x} gates again. For $n_c = 4$ the total number of \texttt{x} gates is
\par\smallskip\noindent\begin{tabularx}{\linewidth}{@{}*{4}{L}@{}}
\correctlabel{A}\,$56$ & \optlabel{B}\,$28$ & \optlabel{C}\,$60$ & \optlabel{D}\,$32$
\end{tabularx}
\end{qitem}
\begin{solbox}{A}
\textbf{\color{brand}Correct answer.} \textcolor{okgreen}{\textbf{(A)}} $56$

\textbf{\color{brand}Concepts required.} Each clock value $k$ is selected by flipping its zero bits, applying the multi-controlled rotation, and flipping back.

\textbf{\color{brand}Relevant formulae.} $\#X = 2\sum_{k=1}^{2^{n_c}-1}(n_c - \mathrm{popcount}(k))$; $\ \sum_{k=0}^{2^{n}-1}\mathrm{popcount}(k) = n2^{n-1}$.

\textbf{\color{brand}Step-by-step explanation.} For each $k$ the number of zero bits is $4 - \mathrm{popcount}(k)$, and each zero bit costs $2$ gates. Summed over $k = 1, \dots, 15$:
$\sum_k\mathrm{popcount}(k) = 4\cdot2^3 = 32$ (each bit is $1$ in half of the $16$ values; $k = 0$ adds nothing), so
\[ 2\sum_{k=1}^{15}(4 - \mathrm{popcount}(k)) = 2\,(60 - 32) = 56. \]
This matches the recorded circuit statistics (\texttt{'x': 56}, \texttt{'ccccry': 15}). $28$ forgets the second (restoring) layer.

\textbf{\color{brand}Why the other options are incorrect.}
\begin{itemize}[leftmargin=2.4em, itemsep=1pt, topsep=2pt]
\item[\textbf{(B)}] $28$ counts only the first layer of $X$ gates, not the restoring layer.
\item[\textbf{(C)}] $60 = 15\times4$ assumes every clock value has four zero bits.
\item[\textbf{(D)}] $32$ counts the one-bits (popcount sum) instead of the zero bits.
\end{itemize}
\textbf{\color{brand}Key takeaway.} Code-level counts can be checked against \texttt{cudaq.estimate\_resources}: the recorded statistics show \texttt{'x': 56}.
\end{solbox}
\qsep

\needspace{13\baselineskip}
\begin{qitem}{53}
\qmeta{Single correct}{Norm recovery $\vec x = \lVert\vec b\rVert\,\vec a/C$}
Let $\vec a$ be the (unnormalised) amplitudes of the system register in the branch clock $= 0\dots0$, ancilla $= 1$ of an exact HHL run. The notebook recovers
the solution, including its norm, as
\par\smallskip\noindent\begin{tabularx}{\linewidth}{@{}*{2}{L}@{}}
\optlabel{A}\,$\vec x = \dfrac{\vec a}{\sqrt{P_{\text{success}}}}$ & \correctlabel{B}\,$\vec x = \dfrac{\lVert\vec b\rVert}{C}\,\vec a$ \\[3pt]
\optlabel{C}\,$\vec x = \dfrac{C}{\lVert\vec b\rVert}\,\vec a$ & \optlabel{D}\,$\vec x = \vec a$
\end{tabularx}
\end{qitem}
\begin{solbox}{B}
\textbf{\color{brand}Correct answer.} \textcolor{okgreen}{\textbf{(B)}} $\vec x = \dfrac{\lVert\vec b\rVert}{C}\,\vec a$

\textbf{\color{brand}Concepts required.} The post-selected branch is $C\,A^{-1}\ket b$; $\ket b$ is $\vec b$ divided by its norm.

\textbf{\color{brand}Relevant formulae.} $\vec a = \frac{C}{\lVert\vec b\rVert}\vec x$, $\ \lVert\vec a\rVert^2 = P_{\text{success}}$, so $\vec x = \frac{\lVert\vec b\rVert}{C}\vec a$.

\textbf{\color{brand}Step-by-step explanation.} The surviving branch is $\vec a = \sum_j\beta_j\frac{C}{\lambda_j}\ket{u_j} = C\,A^{-1}\ket b = C\,A^{-1}\frac{\vec b}{\lVert\vec b\rVert} = \frac{C}{\lVert\vec b\rVert}\vec x$.
Solving for $\vec x$: $\vec x = \frac{\lVert\vec b\rVert}{C}\vec a$ (code: \texttt{x\_quantum\_full = amplitudes * b\_norm / C\_constant}).
$\vec a/\sqrt{P_{\text{success}}}$ is the \emph{normalised} $\ket x$ (correct direction, norm $1$), not $\vec x$. Note that $\lVert\vec a\rVert^2 = P_{\text{success}}$, so on hardware the norm
of $\vec x$ can be estimated from the measured success rate.

\textbf{\color{brand}Why the other options are incorrect.}
\begin{itemize}[leftmargin=2.4em, itemsep=1pt, topsep=2pt]
\item[\textbf{(A)}] $\vec a/\sqrt{P_{\text{success}}}$ is the normalised $\ket x$ --- correct direction, but norm $1$.
\item[\textbf{(C)}] $C/\lVert\vec b\rVert$ inverts the conversion factor.
\item[\textbf{(D)}] $\vec a$ itself equals $\vec x$ only if $\lVert\vec b\rVert = C$ (true in the default run, but not in general).
\end{itemize}
\textbf{\color{brand}Key takeaway.} Scale information lives in $\lVert\vec b\rVert$, $C$ and $P_{\text{success}}$, not in the normalised state.
\end{solbox}
\qsep

\needspace{13\baselineskip}
\begin{qitem}{54}
\qmeta{Single correct}{Default $4\times4$ solution}
For the notebook's default system ($A$ with eigenvalues $1, 2, 3, 4$ and $\vec b = (1, 0, 0, 0)$, see conventions) the exact solution is
\par\smallskip\noindent\begin{tabularx}{\linewidth}{@{}*{2}{L}@{}}
\correctlabel{A}\,$\vec x = \tfrac1{48}(25,\ 7,\ 11,\ 5)$ & \optlabel{B}\,$\vec x = \tfrac1{48}(25,\ -7,\ -11,\ 5)$ \\[3pt]
\optlabel{C}\,$\vec x = \tfrac{1}{\sqrt{820}}(25,\ 7,\ 11,\ 5)$ & \optlabel{D}\,$\vec x = (1,\ \tfrac12,\ \tfrac13,\ \tfrac14)$
\end{tabularx}
\end{qitem}
\begin{solbox}{A}
\textbf{\color{brand}Correct answer.} \textcolor{okgreen}{\textbf{(A)}} $\vec x = \tfrac1{48}(25,\ 7,\ 11,\ 5)$

\textbf{\color{brand}Concepts required.} The default matrix is $H_4\operatorname{diag}(1,2,3,4)H_4$ with the self-inverse orthogonal $H_4$; invert in the eigenbasis.

\textbf{\color{brand}Relevant formulae.} $A^{-1}\vec b = H_4\Lambda^{-1}H_4\vec b$; $\ H_4\vec b = \frac12(1,1,1,1)$ for $\vec b = (1, 0, 0, 0)$.

\textbf{\color{brand}Step-by-step explanation.} $A = H_4\Lambda H_4$ with $H_4 = \frac12\begin{psmallmatrix}1&1&1&1\\1&-1&1&-1\\1&1&-1&-1\\1&-1&-1&1\end{psmallmatrix}$ ($H_4^2 = I$), so $A^{-1} = H_4\Lambda^{-1}H_4$.
$H_4\vec b = \frac12(1, 1, 1, 1)$ (all $\beta_j = \frac12$); $\Lambda^{-1}$ gives $\frac12(1, \frac12, \frac13, \frac14) = \frac1{24}(12, 6, 4, 3)$; then
\[ H_4\cdot\tfrac{1}{24}(12, 6, 4, 3) = \tfrac1{48}(12{+}6{+}4{+}3,\ 12{-}6{+}4{-}3,\ 12{+}6{-}4{-}3,\ 12{-}6{-}4{+}3) = \tfrac1{48}(25, 7, 11, 5). \]
$= (0.520833, 0.145833, 0.229167, 0.104167)$, exactly the recorded output. Check the first row: $2.5\cdot25 - 0.5\cdot7 - 1\cdot11 = 48$.
$\frac{1}{\sqrt{820}}(\dots)$ is the normalised $\ket x$; $(1, \frac12, \frac13, \frac14)$ inverts the eigenvalues without changing basis.

\textbf{\color{brand}Why the other options are incorrect.}
\begin{itemize}[leftmargin=2.4em, itemsep=1pt, topsep=2pt]
\item[\textbf{(B)}] Negative signs would require negative entries of $A^{-1}\vec b$; all four sums $12\pm6\pm4\pm3$ are positive.
\item[\textbf{(C)}] $\frac{1}{\sqrt{820}}(25, 7, 11, 5)$ is the normalised state $\ket x$, not the solution vector ($\lVert\vec x\rVert \approx 0.597$).
\item[\textbf{(D)}] $(1, \frac12, \frac13, \frac14)$ inverts the eigenvalues but forgets to transform to and from the eigenbasis.
\end{itemize}
\textbf{\color{brand}Key takeaway.} Recorded result $= \frac1{48}(25, 7, 11, 5)$ exactly; fidelity $1$ because every eigenvalue hits an integer clock value.
\end{solbox}
\qsep

\needspace{13\baselineskip}
\begin{qitem}{55}
\qmeta{Matching lists}{Algorithm step $\leftrightarrow$ CUDA-Q code}
Match each step of HHL (List-I) with the CUDA-Q / Python code from \texttt{HHL.ipynb} that implements it (List-II).
\par\smallskip\noindent{\renewcommand{\arraystretch}{1.35}%
\begin{tabularx}{\linewidth}{@{}c>{\raggedright\arraybackslash}p{0.33\linewidth}@{\hspace{1.2em}}cL@{}}
\toprule
\multicolumn{2}{@{}l}{\textbf{List-I}} & \multicolumn{2}{l@{}}{\textbf{List-II}} \\ \midrule
\textbf{(P)} & amplitude encoding of $\ket b$ & \textbf{(1)} & \texttt{kernel.cry(theta, clock, ancilla)} \\
\textbf{(Q)} & controlled $U^{2^j} = e^{iAt2^j}$ & \textbf{(2)} & \texttt{cudaq.register\_operation(name, matrix\_power(U, 2**j))} \\
\textbf{(R)} & eigenvalue inversion & \textbf{(3)} & \texttt{bits[:n\_c] == \textquotesingle0\textquotesingle*n\_c and bits[-1] == \textquotesingle1\textquotesingle} \\
\textbf{(S)} & post-selection of the shots & \textbf{(4)} & \texttt{kernel.qalloc(initial\_state\_vector)} \\
\bottomrule
\end{tabularx}}
\par\smallskip\noindent\begin{tabularx}{\linewidth}{@{}*{2}{L}@{}}
\optlabel{A}\,P$\to$4, Q$\to$1, R$\to$2, S$\to$3 & \optlabel{B}\,P$\to$4, Q$\to$2, R$\to$3, S$\to$1 \\[3pt]
\optlabel{C}\,P$\to$2, Q$\to$4, R$\to$1, S$\to$3 & \correctlabel{D}\,P$\to$4, Q$\to$2, R$\to$1, S$\to$3
\end{tabularx}
\end{qitem}
\begin{solbox}{D}
\textbf{\color{brand}Correct answer.} \textcolor{okgreen}{\textbf{(D)}} P$\to$4, Q$\to$2, R$\to$1, S$\to$3

\textbf{\color{brand}Concepts required.} Mapping the algorithm onto the notebook's kernel-builder code.

\textbf{\color{brand}Relevant formulae.} Bit strings are ordered clock $|$ system $|$ ancilla; success: clock $= 0\dots0$, ancilla $= 1$.

\textbf{\color{brand}Step-by-step explanation.} \begin{itemize}\itemsep0pt
\item P: the initial state $\ket{0\dots0}_{\text{clock}}\ket b\ket0$ is loaded in one step with \texttt{qalloc(vector)} (4).
\item Q: each $U^{2^j}$ is an exact matrix registered as a custom gate and applied controlled by one clock qubit (2).
\item R: the multi-controlled $R_Y(\theta_k)$ on the ancilla, $\theta_k = 2\arcsin(C/\tilde\lambda_k)$ (1).
\item S: in the sampled bit strings (clock $|$ system $|$ ancilla) keep clock $= 0\dots0$ and ancilla $= 1$ (3).
\end{itemize}

\textbf{\color{brand}Why the other options are incorrect.}
\begin{itemize}[leftmargin=2.4em, itemsep=1pt, topsep=2pt]
\item[\textbf{(A)}] Q$\to$1, R$\to$2: the \texttt{cry} is the rotation, the registered matrix powers are the controlled evolutions.
\item[\textbf{(B)}] R$\to$3, S$\to$1: post-selection is classical filtering of bit strings, the rotation is a gate.
\item[\textbf{(C)}] P$\to$2, Q$\to$4 swaps state preparation (\texttt{qalloc}) with the registered matrix powers.
\end{itemize}
\textbf{\color{brand}Key takeaway.} Four lines of code carry the four ideas of HHL: load $\ket b$, controlled $e^{iAt2^j}$, controlled $R_Y$, post-select.
\end{solbox}
\qsep

\needspace{13\baselineskip}
\begin{qitem}{56}
\qmeta{Single correct}{Preparing $\ket b$}
Starting from $\ket{00}$, which circuit prepares $\ket b = \frac12(\ket{00} - \ket{01} + \ket{10} - \ket{11})$, i.e.\ $\vec b \propto (1, -1, 1, -1)$?
\par\smallskip\noindent\begin{tabularx}{\linewidth}{@{}*{2}{L}@{}}
\optlabel{A}\,$H$ on both qubits only & \correctlabel{B}\,$H$ on both qubits, then $Z$ on $q_1$ \\[3pt]
\optlabel{C}\,$H$ on both qubits, then $Z$ on $q_0$ & \optlabel{D}\,$X$ on both qubits, then $H$ on both qubits
\end{tabularx}
\end{qitem}
\begin{solbox}{B}
\textbf{\color{brand}Correct answer.} \textcolor{okgreen}{\textbf{(B)}} $H$ on both qubits, then $Z$ on $q_1$

\textbf{\color{brand}Concepts required.} Product vectors are prepared with single-qubit gates; signs come from $Z$ (or from $H$ acting on $\ket1$).

\textbf{\color{brand}Relevant formulae.} $\ket+\ket- = \frac12(1, -1, 1, -1)$; $\ Z\ket+ = \ket-$; $\ H\ket1 = \ket-$.

\textbf{\color{brand}Step-by-step explanation.} Factorise: $\frac12(\ket0 + \ket1)_{q_0}(\ket0 - \ket1)_{q_1} = \ket+\ket-$. $H\ket0 = \ket+$ on $q_0$; on $q_1$, $ZH\ket0 = Z\ket+ = \ket-$.
\begin{itemize}\itemsep0pt
\item $Z$ on $q_0$ gives $\ket-\ket+ \propto (1, 1, -1, -1)$; $H\otimes H$ alone gives $(1, 1, 1, 1)$; $X\otimes X$ then $H\otimes H$ gives $\ket-\ket- \propto (1, -1, -1, 1)$.
\end{itemize}
Structured vectors like this are cheap to prepare; a \emph{general} $\vec b$ needs $O(d)$ gates (or a QRAM), which can cancel HHL's speed-up --- the notebook
side-steps this with exact amplitude initialisation on a simulator.

\textbf{\color{brand}Why the other options are incorrect.}
\begin{itemize}[leftmargin=2.4em, itemsep=1pt, topsep=2pt]
\item[\textbf{(A)}] $H\otimes H$ alone gives $(1, 1, 1, 1)/2$ --- no minus signs.
\item[\textbf{(C)}] $Z$ on $q_0$ produces $\ket-\ket+ \propto (1, 1, -1, -1)$.
\item[\textbf{(D)}] $X\otimes X$ then $H\otimes H$ gives $\ket-\ket- \propto (1, -1, -1, 1)$.
\end{itemize}
\textbf{\color{brand}Key takeaway.} Efficient preparation of $\ket b$ is an assumption of HHL; structured vectors satisfy it, arbitrary data may not.
\end{solbox}
\qsep

\needspace{13\baselineskip}
\begin{qitem}{57}
\qmeta{Single correct}{Role of the inverse QFT}
Inside the phase-estimation part of HHL, the inverse QFT on the clock register
\begin{itemize}[leftmargin=2.6em, labelsep=0.4em, itemsep=2pt, topsep=3pt]
\item[\optlabel{A}] computes $1/\lambda$ for every eigenvalue
\item[\correctlabel{B}] turns the kicked-back phases $e^{2\pi i\,kj/2^{n_c}}$ into the computational-basis state $\ket{k}$ (the binary eigenvalue estimate)
\item[\optlabel{C}] measures the eigenvalues and sends them to the classical computer
\item[\optlabel{D}] prepares $\ket b$ in the eigenbasis of $A$
\end{itemize}
\end{qitem}
\begin{solbox}{B}
\textbf{\color{brand}Correct answer.} \textcolor{okgreen}{\textbf{(B)}} turns the kicked-back phases $e^{2\pi i\,kj/2^{n_c}}$ into the computational-basis state $\ket{k}$ (the binary eigenvalue estimate)

\textbf{\color{brand}Concepts required.} QPE = Hadamards + controlled powers (phase kickback creates a Fourier state) + inverse QFT (reads it out in binary).

\textbf{\color{brand}Relevant formulae.} $\frac{1}{\sqrt{2^{n_c}}}\sum_je^{2\pi ijk/2^{n_c}}\ket j = \mathrm{QFT}\ket k$; $\ \mathrm{QFT}^{-1}$ gives $\ket k$.

\textbf{\color{brand}Step-by-step explanation.} After the Hadamards and controlled powers the clock holds $\frac{1}{\sqrt{2^{n_c}}}\sum_{j}e^{i\lambda t j}\ket j = \mathrm{QFT}\ket{k}$ with $k = \lambda t2^{n_c}/2\pi$
(phase kickback). The inverse QFT maps this Fourier-basis state back to $\ket k$, so the eigenvalue is written as a binary number that a controlled
rotation can read. The division by $\lambda$ happens later (rotation), nothing is measured, and the eigenbasis expansion of $\ket b$ is never
computed explicitly --- linearity does it.

\textbf{\color{brand}Why the other options are incorrect.}
\begin{itemize}[leftmargin=2.4em, itemsep=1pt, topsep=2pt]
\item[\textbf{(A)}] The inverse QFT computes no reciprocals; $1/\lambda$ comes from the controlled rotation.
\item[\textbf{(C)}] Nothing is measured in the middle of HHL; the clock stays coherent and is later uncomputed.
\item[\textbf{(D)}] $\ket b$ is prepared before QPE, and the eigenbasis expansion is never performed explicitly.
\end{itemize}
\textbf{\color{brand}Key takeaway.} Inverse QFT: Fourier-encoded phase $\to$ binary eigenvalue register.
\end{solbox}
\qsep

\needspace{13\baselineskip}
\begin{qitem}{58}
\qmeta{Statement I/II}{$\ket b$ an eigenvector}
Assume $\ket b$ is itself an eigenvector $\ket{u_j}$ of $A$ and $\lambda_j$ is exactly representable by the clock.\\[1mm]
\textbf{Statement I:} The success probability is $C^2/\lambda_j^2$.\\[1mm]
\textbf{Statement II:} After phase estimation the clock register is in the single basis state $\ket{k_j}$ (not entangled with the system).
\begin{itemize}[leftmargin=2.6em, labelsep=0.4em, itemsep=2pt, topsep=3pt]
\item[\correctlabel{A}] Both Statement I and Statement II are true
\item[\optlabel{B}] Both Statement I and Statement II are false
\item[\optlabel{C}] Statement I is true, but Statement II is false
\item[\optlabel{D}] Statement I is false, but Statement II is true
\end{itemize}
\end{qitem}
\begin{solbox}{A}
\textbf{\color{brand}Correct answer.} \textcolor{okgreen}{\textbf{(A)}} Both Statement I and Statement II are true

\textbf{\color{brand}Concepts required.} With a single eigencomponent there is nothing to interfere: QPE gives a product state and HHL reduces to one rotation.

\textbf{\color{brand}Relevant formulae.} $\beta_j = 1 \Rightarrow P_{\text{success}} = C^2/\lambda_j^2$; $\ \ket0\ket{u_j} \to \ket{k_j}\ket{u_j}$.

\textbf{\color{brand}Step-by-step explanation.} With $\beta_j = 1$ and all other $\beta$ zero:
\textbf{I is true:} $P = \sum|\beta|^2C^2/\lambda^2 = C^2/\lambda_j^2$.
\textbf{II is true:} QPE maps $\ket0\ket{u_j} \to \ket{k_j}\ket{u_j}$, a product state; with an exactly representable phase the clock is a single basis state.
The output is $\ket x = \ket{u_j}$ (since $A^{-1}\ket{u_j} = \lambda_j^{-1}\ket{u_j}$), as expected. Answer: (A).

\textbf{\color{brand}Why the other options are incorrect.}
\begin{itemize}[leftmargin=2.4em, itemsep=1pt, topsep=2pt]
\item[\textbf{(B)}] ``both false'' contradicts the general formulas specialised to one component.
\item[\textbf{(C)}] rejects II, but an exact single phase gives a single clock basis state.
\item[\textbf{(D)}] rejects I, but $\sum_j|\beta_j|^2C^2/\lambda_j^2$ has only one term.
\end{itemize}
\textbf{\color{brand}Key takeaway.} Eigenvector input is the simplest sanity test of an HHL implementation.
\end{solbox}
\qsep

\needspace{20\baselineskip}
\begin{passagebox}{Paragraph for Questions 59 and 60}
HHL is run by hand on $A = \begin{pmatrix}\frac32&\frac12\\[2pt] \frac12&\frac32\end{pmatrix}$, $\vec b = (1, 0)$ ($n_b = 1$). The eigenpairs are
$\lambda_1 = 1$, $\ket{u_1} = \frac{1}{\sqrt2}(\ket0 - \ket1)$ and $\lambda_2 = 2$, $\ket{u_2} = \frac{1}{\sqrt2}(\ket0 + \ket1)$. Phase estimation uses $n_c = 2$
clock qubits and $t = \pi/2$; because both eigenvalues are positive, the clock value $k$ is read as an \emph{unsigned} integer here
($\tilde\lambda_k = kC$). States are written as $\ket{\text{clock}}\ket{\text{system}}$; the ancilla starts in $\ket0$ and is not shown until the rotation.
\end{passagebox}
\needspace{13\baselineskip}
\begin{qitem}{59}
\qmeta{Paragraph based}{$2\times2$ HHL trace: after QPE}
The state after phase estimation is
\par\smallskip\noindent\begin{tabularx}{\linewidth}{@{}*{2}{L}@{}}
\optlabel{A}\,$\frac{1}{\sqrt2}\left(\ket{00}\ket{u_1} + \ket{01}\ket{u_2}\right)$ & \optlabel{B}\,$\frac{1}{\sqrt2}\left(\ket{01} + \ket{10}\right)\otimes\ket0$ \\[3pt]
\optlabel{C}\,$\frac{1}{\sqrt2}\left(\ket{10}\ket{u_1} + \ket{01}\ket{u_2}\right)$ & \correctlabel{D}\,$\frac{1}{\sqrt2}\left(\ket{01}\ket{u_1} + \ket{10}\ket{u_2}\right)$
\end{tabularx}
\end{qitem}
\begin{solbox}{D}
\textbf{\color{brand}Correct answer.} \textcolor{okgreen}{\textbf{(D)}} $\frac{1}{\sqrt2}\left(\ket{01}\ket{u_1} + \ket{10}\ket{u_2}\right)$

\textbf{\color{brand}Concepts required.} Expand $\ket b$ in the eigenbasis and apply QPE to each component separately.

\textbf{\color{brand}Relevant formulae.} $\ket0 = \frac{1}{\sqrt2}(\ket{u_1} + \ket{u_2})$; $\ k = \lambda t2^{n_c}/2\pi = \lambda$ for $t = \pi/2$, $n_c = 2$.

\textbf{\color{brand}Step-by-step explanation.} Expand $\ket b = \ket0 = \frac{1}{\sqrt2}(\ket{u_1} + \ket{u_2})$, so $\beta_1 = \beta_2 = \frac{1}{\sqrt2}$.
Clock values: $k = \lambda t2^{n_c}/2\pi = \lambda\cdot\frac\pi2\cdot\frac{4}{2\pi} = \lambda$, so $\lambda_1 = 1 \to \ket{01}$ and $\lambda_2 = 2 \to \ket{10}$, both exact.
Hence $\frac{1}{\sqrt2}(\ket{01}\ket{u_1} + \ket{10}\ket{u_2})$: clock and system are entangled. Also $C = 2\pi/(2^{n_c}t) = 2\pi/(4\cdot\frac\pi2) = 1$.
(Option with $\ket0$ factored out ignores this entanglement; the swapped option exchanges the eigenvalues.)

\textbf{\color{brand}Why the other options are incorrect.}
\begin{itemize}[leftmargin=2.4em, itemsep=1pt, topsep=2pt]
\item[\textbf{(A)}] $\ket{00}$ and $\ket{01}$ correspond to $k = \lambda - 1$; $k = 0$ would mean $\lambda = 0$.
\item[\textbf{(B)}] Factoring out the system ignores that each clock value belongs to a different eigenvector (entanglement).
\item[\textbf{(C)}] Swapping $\ket{01}$ and $\ket{10}$ assigns $\lambda = 2$ to $\ket{u_1}$ and $\lambda = 1$ to $\ket{u_2}$.
\end{itemize}
\textbf{\color{brand}Key takeaway.} After QPE each eigenvector carries its own eigenvalue label.
\end{solbox}
\qsep

\needspace{13\baselineskip}
\begin{qitem}{60}
\qmeta{Paragraph based}{$2\times2$ HHL trace: output}
After the controlled rotation, uncomputation and a successful ancilla measurement, the success probability and the system state are
\par\smallskip\noindent\begin{tabularx}{\linewidth}{@{}*{2}{L}@{}}
\optlabel{A}\,$P = \tfrac58$; $\ \frac{1}{\sqrt{10}}(3\ket0 + \ket1)$ & \optlabel{B}\,$P = \tfrac34$; $\ \frac{1}{\sqrt{10}}(3\ket0 - \ket1)$ \\[3pt]
\correctlabel{C}\,$P = \tfrac58$; $\ \frac{1}{\sqrt{10}}(3\ket0 - \ket1)$ & \optlabel{D}\,$P = 1$; $\ \ket0$
\end{tabularx}
\end{qitem}
\begin{solbox}{C}
\textbf{\color{brand}Correct answer.} \textcolor{okgreen}{\textbf{(C)}} $P = \tfrac58$; $\ \frac{1}{\sqrt{10}}(3\ket0 - \ket1)$

\textbf{\color{brand}Concepts required.} Rotation multiplies each component by $C/\lambda_j$; after uncomputation the components interfere into $C\,A^{-1}\ket b$.

\textbf{\color{brand}Relevant formulae.} branch $= \sum_j\beta_j\frac{C}{\lambda_j}\ket{u_j}$; $\ P = \lVert\text{branch}\rVert^2$; $\ \ket x = \text{branch}/\sqrt P$.

\textbf{\color{brand}Step-by-step explanation.} Rotation: $\ket{01} \to$ ancilla amplitude $C/1 = 1$; $\ket{10} \to C/2 = \frac12$. After uncomputation the ancilla-$1$ branch is
\[ \tfrac{1}{\sqrt2}\Big(1\cdot\ket{u_1} + \tfrac12\ket{u_2}\Big) = \tfrac12\Big[(\ket0 - \ket1) + \tfrac12(\ket0 + \ket1)\Big] = \tfrac34\ket0 - \tfrac14\ket1 . \]
$P = (\frac34)^2 + (\frac14)^2 = \frac{10}{16} = \frac58$ $\big(= \frac12\cdot1 + \frac12\cdot\frac14\big)$, and the normalised state is $\frac{1}{\sqrt{10}}(3\ket0 - \ket1)$.
Classical check: $A^{-1}\vec b = \frac12\begin{psmallmatrix}3/2&-1/2\\-1/2&3/2\end{psmallmatrix}(1, 0)^T = (\frac34, -\frac14)$ --- with $\lVert\vec b\rVert = C = 1$ the branch \emph{is} $\vec x$.
($P = \frac34$ forgets to square the amplitudes.)

\textbf{\color{brand}Why the other options are incorrect.}
\begin{itemize}[leftmargin=2.4em, itemsep=1pt, topsep=2pt]
\item[\textbf{(A)}] $+\ket1$ has the wrong sign: $\ket{u_1}$ (weight $1$) contributes $-\ket1$, $\ket{u_2}$ (weight $\frac12$) only $+\frac12\ket1$.
\item[\textbf{(B)}] $P = \frac34$ adds the amplitudes $\frac12(1 + \frac12)$ instead of the probabilities $\frac12(1 + \frac14)$.
\item[\textbf{(D)}] $P = 1$ and $\ket0$ would require $A = I$; the run succeeds only with probability $\frac58$.
\end{itemize}
\textbf{\color{brand}Key takeaway.} A full HHL run by hand: $\ket x = \frac{1}{\sqrt{10}}(3\ket0 - \ket1) \propto A^{-1}\vec b = (\frac34, -\frac14)$.
\end{solbox}
\qsep

\newpage

\begin{sectionbanner}SECTION D \enspace$\cdot$\enspace Advanced HHL and Implementation \hfill Q.61--80\end{sectionbanner}

\needspace{13\baselineskip}
\begin{qitem}{61}
\qmeta{Single correct}{Complexity of HHL}
For an $s$-sparse Hermitian $N\times N$ matrix with condition number $\kappa$, target error $\epsilon$, and efficient preparation of $\ket b$, the run time
of the original HHL algorithm (2009) scales as
\par\smallskip\noindent\begin{tabularx}{\linewidth}{@{}*{2}{L}@{}}
\optlabel{A}\,$O(\log N)$, independent of $\kappa$ and $\epsilon$ & \optlabel{B}\,$\tilde O(N\,s\,\kappa)$ \\[3pt]
\correctlabel{C}\,$\tilde O\!\left(\log(N)\,s^2\kappa^2/\epsilon\right)$ & \optlabel{D}\,$O(N^3)$
\end{tabularx}
\end{qitem}
\begin{solbox}{C}
\textbf{\color{brand}Correct answer.} \textcolor{okgreen}{\textbf{(C)}} $\tilde O\!\left(\log(N)\,s^2\kappa^2/\epsilon\right)$

\textbf{\color{brand}Concepts required.} HHL's cost combines sparse Hamiltonian simulation, phase-estimation precision and the number of repetitions.

\textbf{\color{brand}Relevant formulae.} $T_{\text{HHL}} = \tilde O(\log(N)s^2\kappa^2/\epsilon)$; Gaussian elimination $O(N^3)$; conjugate gradient $O(Ns\sqrt\kappa\log(1/\epsilon))$.

\textbf{\color{brand}Step-by-step explanation.} Harrow, Hassidim and Lloyd showed a run time $\tilde O(\log(N)s^2\kappa^2/\epsilon)$: $\log N$ from the number of qubits and sparse Hamiltonian
simulation, $s^2$ from simulating $e^{iAt}$ for an $s$-sparse matrix, one factor $\kappa/\epsilon$ from the phase-estimation precision and one factor $\kappa$
from the $O(\kappa^2)$ repetitions reduced to $O(\kappa)$ by amplitude amplification. $O(N^3)$ is Gaussian elimination; $\tilde O(Ns\kappa)$ is roughly
classical conjugate gradient ($Ns\sqrt\kappa$ for positive-definite $A$). No algorithm is independent of $\kappa$.

\textbf{\color{brand}Why the other options are incorrect.}
\begin{itemize}[leftmargin=2.4em, itemsep=1pt, topsep=2pt]
\item[\textbf{(A)}] No linear-system algorithm is independent of $\kappa$: even the optimal quantum algorithm needs $\Omega(\kappa)$.
\item[\textbf{(B)}] $\tilde O(Ns\kappa)$ is the classical iterative (conjugate-gradient-type) scaling, linear in $N$.
\item[\textbf{(D)}] $O(N^3)$ is classical Gaussian elimination.
\end{itemize}
\textbf{\color{brand}Key takeaway.} Exponential improvement in $N$, polynomial cost in $\kappa$ and $1/\epsilon$ --- the speed-up survives only if $\kappa$ is small.
\end{solbox}
\qsep

\needspace{13\baselineskip}
\begin{qitem}{62}
\qmeta{Single correct}{Repetitions and amplitude amplification}
Suppose $P_{\text{success}} = 1/\kappa^2$ with $\kappa = 10$. The expected number of repetitions needed for one successful run without, and with,
amplitude amplification is about
\begin{itemize}[leftmargin=2.6em, labelsep=0.4em, itemsep=2pt, topsep=3pt]
\item[\optlabel{A}] $1$ without; $\ 1$ with
\item[\correctlabel{B}] $100$ without; $\ 10$ with
\item[\optlabel{C}] $100$ without; $\ 100$ with (no improvement is possible)
\item[\optlabel{D}] $10$ without; $\ \sqrt{10}$ with
\end{itemize}
\end{qitem}
\begin{solbox}{B}
\textbf{\color{brand}Correct answer.} \textcolor{okgreen}{\textbf{(B)}} $100$ without; $\ 10$ with

\textbf{\color{brand}Concepts required.} Repeat-until-success needs $1/P$ trials on average; amplitude amplification needs only $O(1/\sqrt P)$.

\textbf{\color{brand}Relevant formulae.} $\mathbb E[\text{trials}] = 1/P$; amplitude amplification: $\approx\frac{\pi}{4\sqrt P}$ iterations.

\textbf{\color{brand}Step-by-step explanation.} Repeat-until-success: the number of trials is geometric with mean $1/P = \kappa^2 = 100$.
Amplitude amplification (Grover-type rotations around the success subspace) raises the success amplitude $\sqrt P = 1/\kappa$ to $O(1)$ with
$O(1/\sqrt P) = O(\kappa) \approx 10$ applications of the HHL circuit and its inverse --- a quadratic saving. This is how the original analysis
turns $\kappa^2$ repetitions into a factor $\kappa$ in the run time. (It requires coherent repetitions, not just re-running and measuring.)

\textbf{\color{brand}Why the other options are incorrect.}
\begin{itemize}[leftmargin=2.4em, itemsep=1pt, topsep=2pt]
\item[\textbf{(A)}] $P = 0.01$ is far from $1$: most runs fail.
\item[\textbf{(C)}] Amplitude amplification gives a quadratic improvement, so $100 \to 10$ is possible.
\item[\textbf{(D)}] $10$ without amplification would be $1/\sqrt P$; repetitions scale like $1/P$.
\end{itemize}
\textbf{\color{brand}Key takeaway.} Post-selection costs $\kappa^2$ repetitions; amplitude amplification brings it down to $\kappa$.
\end{solbox}
\qsep

\needspace{13\baselineskip}
\begin{qitem}{63}
\qmeta{Matching lists}{Quantum linear-system algorithms}
Match each quantum linear-system algorithm (List-I) with its characteristic improvement or complexity (List-II).
\par\smallskip\noindent{\renewcommand{\arraystretch}{1.35}%
\begin{tabularx}{\linewidth}{@{}c>{\raggedright\arraybackslash}p{0.33\linewidth}@{\hspace{1.2em}}cL@{}}
\toprule
\multicolumn{2}{@{}l}{\textbf{List-I}} & \multicolumn{2}{l@{}}{\textbf{List-II}} \\ \midrule
\textbf{(P)} & HHL (Harrow, Hassidim, Lloyd, 2009) & \textbf{(1)} & precision dependence $\mathrm{polylog}(1/\epsilon)$ via Fourier/Chebyshev fitting of $1/x$ (LCU) \\
\textbf{(Q)} & Variable-time amplitude amplification (Ambainis, 2012) & \textbf{(2)} & almost linear dependence on $\kappa$ \\
\textbf{(R)} & Childs, Kothari and Somma (2017) & \textbf{(3)} & $\tilde O(\kappa^2/\epsilon)$ dependence on $\kappa$ and $\epsilon$ \\
\textbf{(S)} & Discrete adiabatic method (Costa et al., 2022) & \textbf{(4)} & optimal $O(\kappa\log(1/\epsilon))$ query complexity \\
\bottomrule
\end{tabularx}}
\par\smallskip\noindent\begin{tabularx}{\linewidth}{@{}*{2}{L}@{}}
\optlabel{A}\,P$\to$2, Q$\to$3, R$\to$1, S$\to$4 & \optlabel{B}\,P$\to$3, Q$\to$2, R$\to$4, S$\to$1 \\[3pt]
\correctlabel{C}\,P$\to$3, Q$\to$2, R$\to$1, S$\to$4 & \optlabel{D}\,P$\to$3, Q$\to$1, R$\to$2, S$\to$4
\end{tabularx}
\end{qitem}
\begin{solbox}{C}
\textbf{\color{brand}Correct answer.} \textcolor{okgreen}{\textbf{(C)}} P$\to$3, Q$\to$2, R$\to$1, S$\to$4

\textbf{\color{brand}Concepts required.} Quantum linear-system algorithms improved HHL's $\kappa$ and $\epsilon$ dependence step by step.

\textbf{\color{brand}Relevant formulae.} HHL $\tilde O(\kappa^2/\epsilon)$; VTAA $\tilde O(\kappa)$; CKS $\mathrm{polylog}(1/\epsilon)$; discrete adiabatic $O(\kappa\log\frac1\epsilon)$.

\textbf{\color{brand}Step-by-step explanation.} \begin{itemize}\itemsep0pt
\item P: HHL uses phase estimation, which gives $\kappa^2/\epsilon$ (3).
\item Q: Ambainis' variable-time amplitude amplification stops ``branches'' with large eigenvalues early and reaches $\tilde O(\kappa\log^3\kappa)$ (2).
\item R: Childs--Kothari--Somma replace phase estimation by a linear combination of unitaries approximating $1/x$, giving $\mathrm{polylog}(1/\epsilon)$ (1).
\item S: Costa et al.\ use a discrete adiabatic walk and reach the optimal $O(\kappa\log(1/\epsilon))$ (4). Quantum singular value transformation (QSVT) gives
a unified framework for these methods.
\end{itemize}

\textbf{\color{brand}Why the other options are incorrect.}
\begin{itemize}[leftmargin=2.4em, itemsep=1pt, topsep=2pt]
\item[\textbf{(A)}] P$\to$2, Q$\to$3 attributes the original $\kappa^2/\epsilon$ to Ambainis and near-linear $\kappa$ to HHL.
\item[\textbf{(B)}] R$\to$4, S$\to$1: the optimal $O(\kappa\log(1/\epsilon))$ belongs to the discrete adiabatic method, not to CKS.
\item[\textbf{(D)}] Q$\to$1, R$\to$2 swaps Ambainis (improved $\kappa$) and Childs--Kothari--Somma (improved $\epsilon$).
\end{itemize}
\textbf{\color{brand}Key takeaway.} Modern methods reach optimal $\kappa\log(1/\epsilon)$; the read-out and input caveats remain.
\end{solbox}
\qsep

\needspace{13\baselineskip}
\begin{qitem}{64}
\qmeta{Single correct}{Read-out cost of all $x_i$}
Why can HHL \emph{not} efficiently output all $d$ components of $\vec x$ for $d = 2^{20}$?
\begin{itemize}[leftmargin=2.6em, labelsep=0.4em, itemsep=2pt, topsep=3pt]
\item[\optlabel{A}] The no-cloning theorem forbids measuring the system register
\item[\correctlabel{B}] Sampling gives each $|x_i|^2$ only statistically; estimating all $d$ values needs $\Omega(d)$ repetitions (about $d/\epsilon^2$ for precision $\epsilon$), destroying the $\log d$ scaling
\item[\optlabel{C}] Complex amplitudes cannot be measured at all
\item[\optlabel{D}] The ancilla measurement destroys $\ket x$ every time, so the system register never contains $\ket x$
\end{itemize}
\end{qitem}
\begin{solbox}{B}
\textbf{\color{brand}Correct answer.} \textcolor{okgreen}{\textbf{(B)}} Sampling gives each $|x_i|^2$ only statistically; estimating all $d$ values needs $\Omega(d)$ repetitions (about $d/\epsilon^2$ for precision $\epsilon$), destroying the $\log d$ scaling

\textbf{\color{brand}Concepts required.} Sampling a $d$-dimensional distribution to precision $\epsilon$ needs about $d/\epsilon^2$ samples.

\textbf{\color{brand}Relevant formulae.} one run $\to$ one sample $i$ with $P(i) = |x_i|^2/\lVert\vec x\rVert^2$; $\ $ tomography $\sim d/\epsilon^2$ runs.

\textbf{\color{brand}Step-by-step explanation.} One run yields one sample $i$ with probability $|x_i|^2/\lVert\vec x\rVert^2$. To see all $d$ components even once needs at least $\sim d$ samples, and to estimate
each with additive precision $\epsilon$ needs about $d/\epsilon^2$ (tomography); signs and phases need further measurement bases. Each sample costs a full
HHL run, so the total is $\Omega(d)$ --- no better than classical. Measurement itself is allowed (no-cloning only forbids \emph{copying} unknown states);
a successful ancilla outcome \emph{leaves} $\ket x$ in the system register.

\textbf{\color{brand}Why the other options are incorrect.}
\begin{itemize}[leftmargin=2.4em, itemsep=1pt, topsep=2pt]
\item[\textbf{(A)}] No-cloning forbids \emph{copying} an unknown state; measuring it is allowed.
\item[\textbf{(C)}] Complex amplitudes are inferred from measurements in several bases --- possible, just expensive.
\item[\textbf{(D)}] A successful ancilla outcome leaves $\ket x$ in the system register; only failed runs lose it.
\end{itemize}
\textbf{\color{brand}Key takeaway.} Printing $\vec x$ costs $\Omega(d)$ --- use HHL for global properties of $\vec x$.
\end{solbox}
\qsep

\needspace{13\baselineskip}
\begin{qitem}{65}
\qmeta{Assertion--Reason}{``Exponential speed-up'' claims}
\textbf{Assertion (A):} HHL gives an exponential speed-up for solving any linear system $A\vec x = \vec b$ and printing the solution $\vec x$.\\[1mm]
\textbf{Reason (R):} Under suitable assumptions (sparse, well-conditioned $A$, efficient preparation of $\ket b$, only a property of $\ket x$ wanted)
the run time of HHL is polylogarithmic in the dimension.
\begin{itemize}[leftmargin=2.6em, labelsep=0.4em, itemsep=2pt, topsep=3pt]
\item[\optlabel{A}] Both Assertion (A) and Reason (R) are true, and R is the correct explanation of A
\item[\optlabel{B}] Both A and R are true, but R is \emph{not} the correct explanation of A
\item[\optlabel{C}] A is true, but R is false
\item[\correctlabel{D}] A is false, but R is true
\end{itemize}
\end{qitem}
\begin{solbox}{D}
\textbf{\color{brand}Correct answer.} \textcolor{okgreen}{\textbf{(D)}} A is false, but R is true

\textbf{\color{brand}Concepts required.} HHL's polylogarithmic run time holds only under several assumptions; printing $\vec x$ violates one of them.

\textbf{\color{brand}Relevant formulae.} $\tilde O(\log(N)s^2\kappa^2/\epsilon)$ assumes sparse $A$, small $\kappa$, efficient $\ket b$, and output = a property of $\ket x$.

\textbf{\color{brand}Step-by-step explanation.} \textbf{A is false}, for three independent reasons: ``any'' system ignores the sparsity and $\kappa$ requirements (for $\kappa \sim d$ the advantage vanishes),
loading an arbitrary $\vec b$ can cost $O(d)$, and printing $\vec x$ needs $\Omega(d)$ repetitions (Q.64).
\textbf{R is true}: this is precisely the HHL result, $\tilde O(\log(N)s^2\kappa^2/\epsilon)$. Aaronson's ``read the fine print'' lists exactly these caveats. Answer: (D).

\textbf{\color{brand}Why the other options are incorrect.}
\begin{itemize}[leftmargin=2.4em, itemsep=1pt, topsep=2pt]
\item[\textbf{(A)}] needs A true, but A overstates HHL: not ``any'' system (sparsity, $\kappa$) and not printing $\vec x$ ($\Omega(d)$ read-out).
\item[\textbf{(B)}] also needs A true; the fine print makes A false.
\item[\textbf{(C)}] needs A true and R false, but A is false and R is precisely the HHL theorem.
\end{itemize}
\textbf{\color{brand}Key takeaway.} Read the fine print: sparse, well conditioned, cheap $\ket b$, and only an expectation value as output.
\end{solbox}
\qsep

\needspace{13\baselineskip}
\begin{qitem}{66}
\qmeta{Single correct}{Too few clock qubits: rescaled $t$}
The notebook is run with \texttt{NUM\_CLOCK\_QUBITS = 3} on the default matrix (eigenvalues $1, 2, 3, 4$). Since $\kappa = 4$ exceeds the $2^{n_c-1} - 1 = 3$ positive levels,
``auto'' sets $t = 2\pi\cdot3/(2^{n_c}\lambda_{\max})$. The ideal clock values $k_j = \lambda_jt2^{n_c}/2\pi$ are
\begin{itemize}[leftmargin=2.6em, labelsep=0.4em, itemsep=2pt, topsep=3pt]
\item[\correctlabel{A}] $0.75,\ 1.5,\ 2.25,\ 3$: only $\lambda = 4$ is exactly representable
\item[\optlabel{B}] $1,\ 2,\ 3,\ 4$: all exact
\item[\optlabel{C}] $0.5,\ 1,\ 1.5,\ 2$: $\lambda = 2$ and $4$ exact
\item[\optlabel{D}] $1.5,\ 3,\ 4.5,\ 6$: aliasing of $\lambda = 3, 4$
\end{itemize}
\end{qitem}
\begin{solbox}{A}
\textbf{\color{brand}Correct answer.} \textcolor{okgreen}{\textbf{(A)}} $0.75,\ 1.5,\ 2.25,\ 3$: only $\lambda = 4$ is exactly representable

\textbf{\color{brand}Concepts required.} If $\kappa$ exceeds the number of positive clock levels, $t$ must be reduced so the largest eigenvalue fits; small eigenvalues then become non-integer.

\textbf{\color{brand}Relevant formulae.} $t = \frac{2\pi k_{\max}}{2^{n_c}\lambda_{\max}}$, $k_{\max} = 2^{n_c-1} - 1$; $\ k_j = \lambda_jt2^{n_c}/2\pi$; $\ C = 2\pi/(2^{n_c}t)$.

\textbf{\color{brand}Step-by-step explanation.} $t = \frac{2\pi\cdot3}{8\cdot4} = \frac{3\pi}{16}$. Then $k_j = \lambda_j\cdot\frac{3\pi}{16}\cdot\frac{8}{2\pi} = \frac34\lambda_j$: $0.75, 1.5, 2.25, 3$.
Only $3$ is an integer, so phase estimation rounds the other three (with spread-out probability, Q.29) and the eigenvalue inversion uses wrong
values of $1/\lambda$: the fidelity drops below $1$. Also $C = 2\pi/(8t) = \frac43 > \lambda_{\min} = 1$: the eigenvalue $1$ is read either as clock value $1$ ($\tilde\lambda = \frac43$, wrong $1/\lambda$) or as $0$ (no rotation, component lost).
$0.5, 1, 1.5, 2$ would be $t = \pi/8$ unchanged. This is the ``Things to try next'' experiment of the notebook.

\textbf{\color{brand}Why the other options are incorrect.}
\begin{itemize}[leftmargin=2.4em, itemsep=1pt, topsep=2pt]
\item[\textbf{(B)}] $1, 2, 3, 4$ would need $k_{\max} \ge 4$, i.e.\ at least $4$ clock qubits.
\item[\textbf{(C)}] $0.5, 1, 1.5, 2$ keeps $t = \pi/8$, but ``auto'' recomputes $t$ when $\kappa > k_{\max}$.
\item[\textbf{(D)}] $1.5, 3, 4.5, 6$ would exceed the representable range; ``auto'' is designed to prevent aliasing.
\end{itemize}
\textbf{\color{brand}Key takeaway.} Too few clock qubits $\Rightarrow$ rounded eigenvalues $\Rightarrow$ wrong $1/\lambda$ $\Rightarrow$ fidelity $< 1$.
\end{solbox}
\qsep

\needspace{13\baselineskip}
\begin{qitem}{67}
\qmeta{Single correct}{Clock qubits for precision $\epsilon$}
Using the rule of thumb $n_c \approx \lceil\log_2(\kappa/\epsilon)\rceil$ for the number of clock qubits, a system with $\kappa = 100$ and target relative precision
$\epsilon = 0.01$ needs about
\par\smallskip\noindent\begin{tabularx}{\linewidth}{@{}*{2}{L}@{}}
\optlabel{A}\,$7$ clock qubits & \optlabel{B}\,$10$ clock qubits \\[3pt]
\optlabel{C}\,$100$ clock qubits & \correctlabel{D}\,$14$ clock qubits
\end{tabularx}
\end{qitem}
\begin{solbox}{D}
\textbf{\color{brand}Correct answer.} \textcolor{okgreen}{\textbf{(D)}} $14$ clock qubits

\textbf{\color{brand}Concepts required.} The clock must resolve the smallest eigenvalue to relative precision $\epsilon$ while spanning the largest.

\textbf{\color{brand}Relevant formulae.} $2^{n_c} \gtrsim \kappa/\epsilon \Rightarrow n_c \approx \lceil\log_2(\kappa/\epsilon)\rceil$; $\ \log_2 10 \approx 3.32$.

\textbf{\color{brand}Step-by-step explanation.} $\kappa/\epsilon = 100/0.01 = 10^4$, $\log_2 10^4 = 4\log_2 10 \approx 4\times3.32 = 13.29$, rounded up: $14$.
Reason: the smallest eigenvalue $\lambda_{\min} = \lambda_{\max}/\kappa$ must be resolved with relative error $\epsilon$, i.e.\ absolute resolution $\epsilon\lambda_{\min}$, while
the clock must also span $\lambda_{\max}$: $2^{n_c} \gtrsim \lambda_{\max}/(\epsilon\lambda_{\min}) = \kappa/\epsilon$. ($7 = \lceil\log_2100\rceil$ ignores $\epsilon$.)
The phase-estimation cost (total evolution time $\propto 2^{n_c}$) grows like $\kappa/\epsilon$.

\textbf{\color{brand}Why the other options are incorrect.}
\begin{itemize}[leftmargin=2.4em, itemsep=1pt, topsep=2pt]
\item[\textbf{(A)}] $7 = \lceil\log_2 100\rceil$ uses only $\kappa$ and ignores the precision.
\item[\textbf{(B)}] $10$ has no justification from the rule; $2^{10} = 1024 < 10^4$.
\item[\textbf{(C)}] $100$ confuses the number of qubits with $\kappa$; qubits grow logarithmically.
\end{itemize}
\textbf{\color{brand}Key takeaway.} Clock qubits $\sim\log(\kappa/\epsilon)$, but QPE time $\sim\kappa/\epsilon$: precision is cheap in qubits, expensive in depth.
\end{solbox}
\qsep

\needspace{13\baselineskip}
\begin{qitem}{68}
\qmeta{Single correct}{Memory of a large HHL simulation}
The notebook's state-vector approach is used for a Hermitian $2^{20}\times2^{20}$ matrix with $n_c = 10$ (fp64, $16$ bytes per amplitude). The number of qubits
and the memory of the state vector are
\par\smallskip\noindent\begin{tabularx}{\linewidth}{@{}*{2}{L}@{}}
\optlabel{A}\,$31$ qubits; $\ \approx 17.2$ GB & \optlabel{B}\,$31$ qubits; $\ 496$ bytes \\[3pt]
\optlabel{C}\,$2^{20} + 11$ qubits; $\ $ beyond any computer & \correctlabel{D}\,$31$ qubits; $\ \approx 34.4$ GB
\end{tabularx}
\end{qitem}
\begin{solbox}{D}
\textbf{\color{brand}Correct answer.} \textcolor{okgreen}{\textbf{(D)}} $31$ qubits; $\ \approx 34.4$ GB

\textbf{\color{brand}Concepts required.} State-vector memory is $16\times2^{\text{qubits}}$ bytes; the exact-matrix approach also stores dense $d\times d$ gates.

\textbf{\color{brand}Relevant formulae.} qubits $= n_b + n_c + 1$; $\ $memory $= 16\cdot2^{n}$ B; $\ 2^{30}$ B $\approx 1.07$ GB.

\textbf{\color{brand}Step-by-step explanation.} $n_b = 20$, $n_c = 10$, ancilla $1$: $31$ qubits. Memory $16\times2^{31} = 2^{35}$ bytes $= 34.36\times10^9$ bytes $\approx 34.4$ GB --- more than most
single GPUs, i.e.\ a multi-GPU run. ($17.2$ GB is $2^{30}$ amplitudes, one qubit fewer.) Note also that registering $e^{iAt2^j}$ as a dense
$2^{20}\times2^{20}$ custom gate would need $16\times2^{40}$ bytes $\approx 17.6$ TB per matrix: exact matrix exponentials do not scale; real
implementations use Hamiltonian simulation.

\textbf{\color{brand}Why the other options are incorrect.}
\begin{itemize}[leftmargin=2.4em, itemsep=1pt, topsep=2pt]
\item[\textbf{(A)}] $17.2$ GB is $16\times2^{30}$, i.e.\ $30$ qubits; $31$ qubits double it.
\item[\textbf{(B)}] $496$ bytes is $16\times31$: one number per qubit instead of per amplitude.
\item[\textbf{(C)}] $2^{20} + 11$ counts one qubit per dimension; $n_b = \log_2 2^{20} = 20$.
\end{itemize}
\textbf{\color{brand}Key takeaway.} Simulating HHL at scale is limited by exponential memory --- and by the dense $e^{iAt}$ matrices.
\end{solbox}
\qsep

\needspace{13\baselineskip}
\begin{qitem}{69}
\qmeta{Single correct}{Trotter steps for $e^{iAt}$}
$e^{iAt}$ with $A = A_1 + A_2$ ($[A_1, A_2] \neq 0$) is approximated by $r$ first-order Trotter steps $\left(e^{iA_1t/r}e^{iA_2t/r}\right)^r$. To keep the total
error below $\epsilon$, $r$ must grow as
\begin{itemize}[leftmargin=2.6em, labelsep=0.4em, itemsep=2pt, topsep=3pt]
\item[\correctlabel{A}] $r = O\!\left(t^2\lVert[A_1, A_2]\rVert/\epsilon\right)$
\item[\optlabel{B}] $r = O(\log(1/\epsilon))$, independent of $t$
\item[\optlabel{C}] $r = O(t/\epsilon^2)$
\item[\optlabel{D}] $r = 1$ always suffices, because the error cancels
\end{itemize}
\end{qitem}
\begin{solbox}{A}
\textbf{\color{brand}Correct answer.} \textcolor{okgreen}{\textbf{(A)}} $r = O\!\left(t^2\lVert[A_1, A_2]\rVert/\epsilon\right)$

\textbf{\color{brand}Concepts required.} First-order Trotterisation error comes from the commutator and accumulates linearly over the steps.

\textbf{\color{brand}Relevant formulae.} $\lVert e^{i(A_1+A_2)t} - (e^{iA_1t/r}e^{iA_2t/r})^r\rVert \le \frac{t^2}{2r}\lVert[A_1, A_2]\rVert$.

\textbf{\color{brand}Step-by-step explanation.} Baker--Campbell--Hausdorff: $e^{iA_1\tau}e^{iA_2\tau} = e^{i(A_1+A_2)\tau - \frac{\tau^2}{2}[A_1, A_2] + O(\tau^3)}$, so each step with $\tau = t/r$ has error
$\approx\frac{t^2}{2r^2}\lVert[A_1, A_2]\rVert$. Over $r$ steps the errors add: $\approx\frac{t^2}{2r}\lVert[A_1, A_2]\rVert \le \epsilon \Rightarrow r = O(t^2\lVert[A_1,A_2]\rVert/\epsilon)$.
Higher-order product formulas, qubitisation or QSVT reduce this cost. (The notebook avoids the issue entirely with exact $e^{iAt}$ gates.)

\textbf{\color{brand}Why the other options are incorrect.}
\begin{itemize}[leftmargin=2.4em, itemsep=1pt, topsep=2pt]
\item[\textbf{(B)}] $O(\log(1/\epsilon))$ is achievable only with advanced methods (qubitisation/QSVT), not first-order Trotter, and never independent of $t$.
\item[\textbf{(C)}] $O(t/\epsilon^2)$ has the wrong powers; the error per step is quadratic in $t/r$.
\item[\textbf{(D)}] A single step leaves error $\frac{t^2}{2}\lVert[A_1, A_2]\rVert$, which does not cancel for non-commuting terms.
\end{itemize}
\textbf{\color{brand}Key takeaway.} Hamiltonian simulation is a major cost of real HHL circuits; the notebook replaces it by exact matrices.
\end{solbox}
\qsep

\needspace{13\baselineskip}
\begin{qitem}{70}
\qmeta{Statement I/II}{What the notebook implements}
About \texttt{HHL.ipynb}:\\[1mm]
\textbf{Statement I:} The controlled evolutions $e^{iAt2^j}$ are implemented as exact matrix exponentials (\texttt{scipy.linalg.expm}) registered as custom gates, not
by Hamiltonian simulation.\\[1mm]
\textbf{Statement II:} The notebook executes HHL on quantum-hardware targets when \texttt{BACKEND = "hardware"}.
\begin{itemize}[leftmargin=2.6em, labelsep=0.4em, itemsep=2pt, topsep=3pt]
\item[\optlabel{A}] Both Statement I and Statement II are true
\item[\optlabel{B}] Both Statement I and Statement II are false
\item[\correctlabel{C}] Statement I is true, but Statement II is false
\item[\optlabel{D}] Statement I is false, but Statement II is true
\end{itemize}
\end{qitem}
\begin{solbox}{C}
\textbf{\color{brand}Correct answer.} \textcolor{okgreen}{\textbf{(C)}} Statement I is true, but Statement II is false

\textbf{\color{brand}Concepts required.} The notebook is a simulator demonstration: exact $e^{iAt}$ gates, amplitude encoding, CPU/GPU simulators only.

\textbf{\color{brand}Relevant formulae.} \texttt{U\_step = expm(1j * A\_work * t\_evolution)}; allowed methods: \texttt{statevector}, \texttt{density\_matrix}.

\textbf{\color{brand}Step-by-step explanation.} \textbf{I is true:} \texttt{U\_step = expm(1j * A\_work * t\_evolution)}; each power is registered with \texttt{cudaq.register\_operation}.
\textbf{II is false:} the backend cell stops with an error for \texttt{"hardware"}: ``This HHL implementation uses exact matrix exponentials as custom gates and
amplitude encoding, which hardware targets cannot execute. Use a simulator backend.'' The allowed methods are \texttt{statevector} and
\texttt{density\_matrix}; the recorded results come from the CPU simulator \texttt{qpp-cpu}. Answer: (C).

\textbf{\color{brand}Why the other options are incorrect.}
\begin{itemize}[leftmargin=2.4em, itemsep=1pt, topsep=2pt]
\item[\textbf{(A)}] accepts II, but the code raises an error for hardware backends.
\item[\textbf{(B)}] rejects I, but the controlled evolutions are registered exact matrices.
\item[\textbf{(D)}] rejects I and accepts II --- both reversed.
\end{itemize}
\textbf{\color{brand}Key takeaway.} No HHL hardware run was made or is possible with this code; recorded results come from \texttt{qpp-cpu}.
\end{solbox}
\qsep

\needspace{13\baselineskip}
\begin{qitem}{71}
\qmeta{Single correct}{Recorded $P_{\text{success}}$: shot statistics}
In the recorded run the ideal success probability is $0.355903$ and $7086$ of $20\,000$ shots ($0.3543$) were accepted. The binomial standard error of the
measured fraction and the verdict are
\begin{itemize}[leftmargin=2.6em, labelsep=0.4em, itemsep=2pt, topsep=3pt]
\item[\optlabel{A}] $\approx0.0007$; the deviation is $\approx2.3\sigma$
\item[\optlabel{B}] $\approx0.0034$; the deviation is $\approx4.7\sigma$, a circuit error
\item[\correctlabel{C}] $\approx0.0034$; the deviation is $\approx0.5\sigma$, consistent with shot noise
\item[\optlabel{D}] $\approx0.012$; the deviation is $\approx0.1\sigma$
\end{itemize}
\end{qitem}
\begin{solbox}{C}
\textbf{\color{brand}Correct answer.} \textcolor{okgreen}{\textbf{(C)}} $\approx0.0034$; the deviation is $\approx0.5\sigma$, consistent with shot noise

\textbf{\color{brand}Concepts required.} The fraction of accepted shots is binomial; compare deviations with its standard error.

\textbf{\color{brand}Relevant formulae.} $\sigma_{\hat p} = \sqrt{p(1-p)/N}$; $\ z = (\hat p - p)/\sigma_{\hat p}$.

\textbf{\color{brand}Step-by-step explanation.} $\sigma = \sqrt{\frac{p(1-p)}{N}} = \sqrt{\frac{0.355903\times0.644097}{20000}} = \sqrt{1.146\times10^{-5}} \approx 0.00339$.
Deviation: $0.3543 - 0.3559 = -0.0016$, i.e.\ $-0.0016/0.0034 \approx -0.47\sigma$ --- well within statistical fluctuations. ($0.012 \approx 1/\sqrt{7086}$ is the
notebook's rough error scale for the \emph{accepted} shots, not this standard error.) The state-vector run gives exactly $0.355903$.

\textbf{\color{brand}Why the other options are incorrect.}
\begin{itemize}[leftmargin=2.4em, itemsep=1pt, topsep=2pt]
\item[\textbf{(A)}] $0.0007$ is about five times too small: $\sqrt{0.3559\times0.6441/20000} = 0.0034$.
\item[\textbf{(B)}] $4.7\sigma$ uses the wrong deviation (or $\sigma$ ten times too small); the observed $z \approx -0.5$.
\item[\textbf{(D)}] $0.012 = 1/\sqrt{7086}$ is the notebook's error scale for the accepted shots, not for this fraction of $20\,000$.
\end{itemize}
\textbf{\color{brand}Key takeaway.} Before blaming the circuit, compare with $\sqrt{p(1-p)/N}$: here everything is consistent.
\end{solbox}
\qsep

\needspace{13\baselineskip}
\begin{qitem}{72}
\qmeta{Single correct}{Recorded read-out of $|x_0|^2$}
In the recorded run, the post-selected shots ($7086$ accepted) estimate $|x_0|^2/\lVert\vec x\rVert^2$ as $0.7691$; the exact value is $0.7622$. The binomial standard
error of this estimate and the deviation in units of it are
\par\smallskip\noindent\begin{tabularx}{\linewidth}{@{}*{2}{L}@{}}
\correctlabel{A}\,$\approx0.0051$; $\ \approx1.4\sigma$ & \optlabel{B}\,$\approx0.012$; $\ \approx0.6\sigma$, the exact standard error \\[3pt]
\optlabel{C}\,$\approx0.0051$; $\ \approx14\sigma$ & \optlabel{D}\,$\approx0.0030$; $\ \approx2.3\sigma$
\end{tabularx}
\end{qitem}
\begin{solbox}{A}
\textbf{\color{brand}Correct answer.} \textcolor{okgreen}{\textbf{(A)}} $\approx0.0051$; $\ \approx1.4\sigma$

\textbf{\color{brand}Concepts required.} Read-out statistics use only the accepted (post-selected) shots.

\textbf{\color{brand}Relevant formulae.} $\sigma = \sqrt{p(1-p)/N_{\text{acc}}}$, $N_{\text{acc}} = P_{\text{success}}\cdot\text{SHOTS}$; $\ \sqrt{p(1-p)} \le \frac12$.

\textbf{\color{brand}Step-by-step explanation.} $\sigma = \sqrt{\frac{p(1-p)}{N_{\text{acc}}}} = \sqrt{\frac{0.7622\times0.2378}{7086}} = \sqrt{2.558\times10^{-5}} \approx 0.00506$.
Deviation $0.7691 - 0.7622 = 0.0069 \approx 1.4\sigma$: an ordinary fluctuation. The notebook prints $1/\sqrt{7086} \approx 0.012$ as a conservative scale
(since $\sqrt{p(1-p)} \le \frac12$, the true standard error is at most $0.0059$). Only the accepted shots count: the effective sample size is
$P_{\text{success}}\times\text{SHOTS}$, so a low success probability directly increases the read-out error.

\textbf{\color{brand}Why the other options are incorrect.}
\begin{itemize}[leftmargin=2.4em, itemsep=1pt, topsep=2pt]
\item[\textbf{(B)}] $1/\sqrt{N_{\text{acc}}}$ is a conservative scale, not the binomial standard error (which is at most half of it).
\item[\textbf{(C)}] $14\sigma$ divides $0.0069$ by $0.0005$; with $\sigma \approx 0.005$ the deviation is $1.4\sigma$.
\item[\textbf{(D)}] $0.0030$ uses all $20\,000$ shots; failed shots carry no information about $\vec x$.
\end{itemize}
\textbf{\color{brand}Key takeaway.} Low success probability means fewer useful shots and a larger read-out error.
\end{solbox}
\qsep

\needspace{13\baselineskip}
\begin{qitem}{73}
\qmeta{Single correct}{Circuit resources vs.\ $n_c$}
The default circuit ($n_c = 4$, $4\times4$ matrix) has $15$ multi-controlled rotations and $8$ controlled $e^{\pm iAt2^j}$ gates on $7$ qubits. With the
same matrix and $n_c = 5$ the notebook's circuit has
\begin{itemize}[leftmargin=2.6em, labelsep=0.4em, itemsep=2pt, topsep=3pt]
\item[\optlabel{A}] $31$ rotations, $\ 10$ controlled $e^{\pm iAt2^j}$, $\ 7$ qubits
\item[\correctlabel{B}] $31$ rotations, $\ 10$ controlled $e^{\pm iAt2^j}$, $\ 8$ qubits
\item[\optlabel{C}] $16$ rotations, $\ 10$ controlled $e^{\pm iAt2^j}$, $\ 8$ qubits
\item[\optlabel{D}] $31$ rotations, $\ 5$ controlled $e^{\pm iAt2^j}$, $\ 8$ qubits
\end{itemize}
\end{qitem}
\begin{solbox}{B}
\textbf{\color{brand}Correct answer.} \textcolor{okgreen}{\textbf{(B)}} $31$ rotations, $\ 10$ controlled $e^{\pm iAt2^j}$, $\ 8$ qubits

\textbf{\color{brand}Concepts required.} The naive eigenvalue inversion uses one multi-controlled rotation per non-zero clock value.

\textbf{\color{brand}Relevant formulae.} rotations $= 2^{n_c} - 1$; controlled evolutions $= 2n_c$; qubits $= n_c + n_b + 1$.

\textbf{\color{brand}Step-by-step explanation.} Rotations: one per non-zero clock value, $2^{n_c} - 1 = 31$ (exponential in $n_c$ --- this naive inversion does not scale; practical versions compute
$1/\lambda$ arithmetically or use QSVT). Controlled evolutions: $n_c$ in phase estimation plus $n_c$ in uncomputation, $2\times5 = 10$.
Qubits: $5 + 2 + 1 = 8$. (Each rotation is now controlled by $5$ qubits.)

\textbf{\color{brand}Why the other options are incorrect.}
\begin{itemize}[leftmargin=2.4em, itemsep=1pt, topsep=2pt]
\item[\textbf{(A)}] $7$ qubits is the $n_c = 4$ count; one more clock qubit means $8$.
\item[\textbf{(C)}] $16 = 2^4$ is the old count $+1$; it should be $2^5 - 1$.
\item[\textbf{(D)}] $5$ counts only the QPE half; uncomputation needs another $5$ inverse evolutions.
\end{itemize}
\textbf{\color{brand}Key takeaway.} The look-up-table inversion is exponential in $n_c$: fine for demonstrations, not for scaling.
\end{solbox}
\qsep

\needspace{13\baselineskip}
\begin{qitem}{74}
\qmeta{Matching lists}{Error sources and their effects}
Match each error source (List-I) with its main effect on HHL (List-II).
\par\smallskip\noindent{\renewcommand{\arraystretch}{1.35}%
\begin{tabularx}{\linewidth}{@{}c>{\raggedright\arraybackslash}p{0.33\linewidth}@{\hspace{1.2em}}cL@{}}
\toprule
\multicolumn{2}{@{}l}{\textbf{List-I}} & \multicolumn{2}{l@{}}{\textbf{List-II}} \\ \midrule
\textbf{(P)} & bit flip on the ancilla just before measurement & \textbf{(1)} & statistical error $\sim 1/\sqrt{N_{\text{acc}}}$ in the estimated $|x_i|^2$ \\
\textbf{(Q)} & phase error on a clock qubit during phase estimation & \textbf{(2)} & wrong eigenvalue estimate, so a wrong factor $1/\tilde\lambda$ is applied \\
\textbf{(R)} & imperfect uncomputation of the clock & \textbf{(3)} & failed runs are accepted (and good ones rejected): biased $P_{\text{success}}$ and contaminated output \\
\textbf{(S)} & finite number of shots & \textbf{(4)} & residual clock--system entanglement: the output is a mixed state \\
\bottomrule
\end{tabularx}}
\par\smallskip\noindent\begin{tabularx}{\linewidth}{@{}*{2}{L}@{}}
\optlabel{A}\,P$\to$3, Q$\to$2, R$\to$1, S$\to$4 & \optlabel{B}\,P$\to$2, Q$\to$3, R$\to$4, S$\to$1 \\[3pt]
\optlabel{C}\,P$\to$3, Q$\to$4, R$\to$2, S$\to$1 & \correctlabel{D}\,P$\to$3, Q$\to$2, R$\to$4, S$\to$1
\end{tabularx}
\end{qitem}
\begin{solbox}{D}
\textbf{\color{brand}Correct answer.} \textcolor{okgreen}{\textbf{(D)}} P$\to$3, Q$\to$2, R$\to$4, S$\to$1

\textbf{\color{brand}Concepts required.} Each part of HHL has a characteristic failure mode.

\textbf{\color{brand}Relevant formulae.} ancilla flip: wrong branch; clock error: $1/\tilde\lambda \neq 1/\lambda$; residual entanglement: mixed $\rho$; shots: $\sigma \propto 1/\sqrt{N_{\text{acc}}}$.

\textbf{\color{brand}Step-by-step explanation.} \begin{itemize}\itemsep0pt
\item P: the ancilla is the success flag; flipping it mixes the $\sqrt{1 - C^2/\lambda^2}$ branch (not $\propto A^{-1}$) into the accepted data (3).
\item Q: QPE reads phases; a phase error shifts the estimated $k$, and the rotation then uses the wrong $\tilde\lambda$ (2).
\item R: if $\ket{k_j}$ is not reset to $\ket0$, the eigencomponents stay labelled and do not interfere: mixed output (4).
\item S: sampling noise with the \emph{accepted} shots as sample size (1).
\end{itemize}
HHL's deep circuits (controlled $e^{iAt2^j}$ with large $2^j$) make it far beyond today's noisy hardware for non-trivial sizes.

\textbf{\color{brand}Why the other options are incorrect.}
\begin{itemize}[leftmargin=2.4em, itemsep=1pt, topsep=2pt]
\item[\textbf{(A)}] R$\to$1, S$\to$4: finite shots give statistical error, not entanglement.
\item[\textbf{(B)}] P$\to$2, Q$\to$3 swaps the ancilla and clock errors: phase errors in QPE corrupt eigenvalue estimates.
\item[\textbf{(C)}] Q$\to$4, R$\to$2 swaps clock errors and imperfect uncomputation.
\end{itemize}
\textbf{\color{brand}Key takeaway.} Debugging HHL: check eigenvalue registers, the success flag and the uncomputation separately.
\end{solbox}
\qsep

\needspace{13\baselineskip}
\begin{qitem}{75}
\qmeta{Single correct}{Useful outputs of HHL}
Which output can HHL provide efficiently (under its assumptions)?
\begin{itemize}[leftmargin=2.6em, labelsep=0.4em, itemsep=2pt, topsep=3pt]
\item[\optlabel{A}] The full inverse matrix $A^{-1}$
\item[\correctlabel{B}] An estimate of an expectation value $\braket{x|M|x}$ for an efficiently measurable observable $M$
\item[\optlabel{C}] The determinant of $A$
\item[\optlabel{D}] The complete vector $\vec x$ to 16 significant digits
\end{itemize}
\end{qitem}
\begin{solbox}{B}
\textbf{\color{brand}Correct answer.} \textcolor{okgreen}{\textbf{(B)}} An estimate of an expectation value $\braket{x|M|x}$ for an efficiently measurable observable $M$

\textbf{\color{brand}Concepts required.} HHL is efficient for estimating scalar properties of $\ket x$, not for outputting vectors or matrices.

\textbf{\color{brand}Relevant formulae.} $\braket{x|M|x}$ to precision $\epsilon$: $O(1/\epsilon^2)$ samples ($O(1/\epsilon)$ with amplitude estimation).

\textbf{\color{brand}Step-by-step explanation.} Measuring an observable $M$ on $\ket x$ repeatedly estimates $\braket{x|M|x}$ to precision $\epsilon$ with $O(1/\epsilon^2)$ runs (or $O(1/\epsilon)$ with amplitude estimation),
independently of $d$. Examples: the total weight of a region of $\vec x$, an overlap $|\braket{y|x}|^2$ via a swap test, a cost function in machine learning.
Printing $\vec x$ needs $\Omega(d)$ runs (Q.64); $\det A$ and $A^{-1}$ ($d^2$ numbers) are not produced by HHL at all.

\textbf{\color{brand}Why the other options are incorrect.}
\begin{itemize}[leftmargin=2.4em, itemsep=1pt, topsep=2pt]
\item[\textbf{(A)}] $A^{-1}$ has $d^2$ entries; HHL only applies $A^{-1}$ to one state.
\item[\textbf{(C)}] HHL does not compute determinants.
\item[\textbf{(D)}] The complete vector needs $\Omega(d)$ runs (Q.64).
\end{itemize}
\textbf{\color{brand}Key takeaway.} Ask HHL scalar questions about $\vec x$.
\end{solbox}
\qsep

\needspace{13\baselineskip}
\begin{qitem}{76}
\qmeta{Single correct}{Dequantisation}
``Dequantisation'' results (Tang 2019 and successors) give classical algorithms with run time polylogarithmic in $d$ for linear systems when
\begin{itemize}[leftmargin=2.6em, labelsep=0.4em, itemsep=2pt, topsep=3pt]
\item[\optlabel{A}] $A$ is Hermitian
\item[\correctlabel{B}] $A$ has low rank and $\ell^2$-norm sample-and-query access to $A$ and $\vec b$ is available
\item[\optlabel{C}] $\vec b$ is a computational basis vector
\item[\optlabel{D}] $A$ is sparse and well conditioned with $\kappa = \mathrm{polylog}(d)$, and $\vec b$ is given as an efficient circuit
\end{itemize}
\end{qitem}
\begin{solbox}{B}
\textbf{\color{brand}Correct answer.} \textcolor{okgreen}{\textbf{(B)}} $A$ has low rank and $\ell^2$-norm sample-and-query access to $A$ and $\vec b$ is available

\textbf{\color{brand}Concepts required.} Quantum-inspired classical algorithms use sample-and-query access (the classical analogue of QRAM) and exploit low rank.

\textbf{\color{brand}Relevant formulae.} classical time $\mathrm{poly}(\text{rank}, \kappa, 1/\epsilon)\cdot\mathrm{polylog}(d)$ for low-rank $A$ with $\ell^2$ sampling access.

\textbf{\color{brand}Step-by-step explanation.} Tang's quantum-inspired algorithms show that if one can sample rows/entries with probability proportional to their squared norms (the classical analogue
of amplitude encoding / QRAM), low-rank problems such as recommendation systems and low-rank linear systems can be solved classically in
time $\mathrm{poly}(\text{rank}, \kappa, 1/\epsilon)\,\mathrm{polylog}(d)$. For sparse, well-conditioned, full-rank systems with efficiently prepared $\ket b$, HHL-type problems
are BQP-complete and no such classical algorithm is expected. Hermiticity or a basis vector $\vec b$ alone gives no classical shortcut.

\textbf{\color{brand}Why the other options are incorrect.}
\begin{itemize}[leftmargin=2.4em, itemsep=1pt, topsep=2pt]
\item[\textbf{(A)}] Hermiticity alone gives no classical shortcut.
\item[\textbf{(C)}] A basis-vector $\vec b$ does not make $A^{-1}\vec b$ classically easy for a general $A$.
\item[\textbf{(D)}] Sparse, well-conditioned, full-rank systems with efficient $\ket b$ are BQP-complete; dequantisation is not expected there.
\end{itemize}
\textbf{\color{brand}Key takeaway.} Low rank + sampling access = classical polylog; sparse full-rank + small $\kappa$ = where quantum advantage may remain.
\end{solbox}
\qsep

\needspace{13\baselineskip}
\begin{qitem}{77}
\qmeta{Assertion--Reason}{$\kappa^2$ run-time dependence}
\textbf{Assertion (A):} At fixed $N$, $s$ and $\epsilon$, doubling the condition number roughly quadruples the run time of the original HHL algorithm.\\[1mm]
\textbf{Reason (R):} The original HHL run time scales as $\tilde O(\log(N)s^2\kappa^2/\epsilon)$.
\begin{itemize}[leftmargin=2.6em, labelsep=0.4em, itemsep=2pt, topsep=3pt]
\item[\correctlabel{A}] Both Assertion (A) and Reason (R) are true, and R is the correct explanation of A
\item[\optlabel{B}] Both A and R are true, but R is \emph{not} the correct explanation of A
\item[\optlabel{C}] A is true, but R is false
\item[\optlabel{D}] A is false, but R is true
\end{itemize}
\end{qitem}
\begin{solbox}{A}
\textbf{\color{brand}Correct answer.} \textcolor{okgreen}{\textbf{(A)}} Both Assertion (A) and Reason (R) are true, and R is the correct explanation of A

\textbf{\color{brand}Concepts required.} With other parameters fixed, HHL's run time is quadratic in $\kappa$.

\textbf{\color{brand}Relevant formulae.} $T \propto \kappa^2 \Rightarrow T(2\kappa)/T(\kappa) = 4$.

\textbf{\color{brand}Step-by-step explanation.} \textbf{R is true} (Q.61). \textbf{A is true and follows from R}: with everything else fixed, $T \propto \kappa^2$, so $\kappa \to 2\kappa$ gives $T \to 4T$ (up to
logarithmic factors). The two factors of $\kappa$ come from phase-estimation precision ($t \sim \kappa/\epsilon$) and from the repetitions after
amplitude amplification ($\sim\kappa$). Improved algorithms (Q.63) reduce this to linear. Answer: (A).

\textbf{\color{brand}Why the other options are incorrect.}
\begin{itemize}[leftmargin=2.4em, itemsep=1pt, topsep=2pt]
\item[\textbf{(B)}] denies the explanatory link, but A is computed directly from R.
\item[\textbf{(C)}] rejects R, which is the original HHL bound.
\item[\textbf{(D)}] rejects A, but $\kappa^2$ scaling gives exactly a factor $4$.
\end{itemize}
\textbf{\color{brand}Key takeaway.} $\kappa$ is the price tag of HHL; improved algorithms make it linear.
\end{solbox}
\qsep

\needspace{13\baselineskip}
\begin{qitem}{78}
\qmeta{Single correct}{Solution inside the embedded system}
For a non-Hermitian invertible $A$ the notebook solves $\begin{pmatrix}0 & A\\ A^\dagger & 0\end{pmatrix}\begin{pmatrix}\vec y\\ \vec z\end{pmatrix} = \begin{pmatrix}\vec b\\ \vec 0\end{pmatrix}$. The solution satisfies
\par\smallskip\noindent\begin{tabularx}{\linewidth}{@{}*{2}{L}@{}}
\optlabel{A}\,$\vec y = A^{-1}\vec b$ (first half) and $\vec z = \vec 0$ & \optlabel{B}\,$\vec y = \vec z = A^{-1}\vec b/\sqrt2$ \\[3pt]
\optlabel{C}\,$\vec y = A^{-1}\vec b$ and $\vec z = (A^\dagger)^{-1}\vec b$ & \correctlabel{D}\,$\vec z = A^{-1}\vec b$ (second half) and $\vec y = \vec 0$
\end{tabularx}
\end{qitem}
\begin{solbox}{D}
\textbf{\color{brand}Correct answer.} \textcolor{okgreen}{\textbf{(D)}} $\vec z = A^{-1}\vec b$ (second half) and $\vec y = \vec 0$

\textbf{\color{brand}Concepts required.} The Hermitian embedding solves the original system in one block of a doubled system.

\textbf{\color{brand}Relevant formulae.} $\begin{psmallmatrix}0&A\\A^\dagger&0\end{psmallmatrix}\begin{psmallmatrix}\vec y\\ \vec z\end{psmallmatrix} = \begin{psmallmatrix}A\vec z\\ A^\dagger\vec y\end{psmallmatrix}$.

\textbf{\color{brand}Step-by-step explanation.} Multiply out: the first block row gives $A\vec z = \vec b$, so $\vec z = A^{-1}\vec b = \vec x$; the second gives $A^\dagger\vec y = \vec 0$, so $\vec y = \vec 0$ ($A^\dagger$ is invertible).
Hence the solution sits in the \emph{second} half of the doubled system register (code: \texttt{solution\_slice = slice(d0, 2 * d0)}). The extra
qubit distinguishes the halves; post-selecting it on $1$ isolates $\ket x$ (here it is $1$ with certainty, since $\vec y = \vec 0$).

\textbf{\color{brand}Why the other options are incorrect.}
\begin{itemize}[leftmargin=2.4em, itemsep=1pt, topsep=2pt]
\item[\textbf{(A)}] The first block row involves $A\vec z$, so the solution is $\vec z$ (second half), not $\vec y$.
\item[\textbf{(B)}] The solution is not split over both halves; $\vec y$ vanishes exactly.
\item[\textbf{(C)}] $A^\dagger\vec y = \vec 0$ forces $\vec y = \vec 0$; neither half contains $(A^\dagger)^{-1}\vec b$.
\end{itemize}
\textbf{\color{brand}Key takeaway.} Embedding costs one qubit; the answer is in the half selected by that qubit.
\end{solbox}
\qsep

\needspace{20\baselineskip}
\begin{passagebox}{Paragraph for Questions 79 and 80}
The previously recorded default run of \texttt{HHL.ipynb} (7 qubits, \texttt{qpp-cpu}) printed:\\[2pt]
{\small\ttfamily Success probability: state 0.355903, shots 0.3543 (7086 of 20000), ideal 0.355903\\
Fidelity |<x\_classical|x\_quantum>|\^{}2 : 1.0000000000\\
Relative error of the recovered vector: 1.655e-15}\\[2pt]
In the eigenbasis of the default $A$ (eigenvalues $1, 2, 3, 4$) the vector $\ket b = \ket{00}$ has all four coefficients $\beta_j = \tfrac12$; $C = 1$ and $\lVert\vec b\rVert = 1$.
\end{passagebox}
\needspace{13\baselineskip}
\begin{qitem}{79}
\qmeta{Paragraph based}{Recorded run: $P_{\text{success}}$ exactly}
The exact value of the ideal success probability $0.355903$ is
\par\smallskip\noindent\begin{tabularx}{\linewidth}{@{}*{4}{L}@{}}
\correctlabel{A}\,$\dfrac{205}{576}$ & \optlabel{B}\,$\dfrac{25}{48}$ & \optlabel{C}\,$\dfrac{205}{144}$ & \optlabel{D}\,$\dfrac14$
\end{tabularx}
\end{qitem}
\begin{solbox}{A}
\textbf{\color{brand}Correct answer.} \textcolor{okgreen}{\textbf{(A)}} $\dfrac{205}{576}$

\textbf{\color{brand}Concepts required.} Exact success probability from the eigenbasis coefficients and the eigenvalues.

\textbf{\color{brand}Relevant formulae.} $P = \sum_j|\beta_j|^2C^2/\lambda_j^2$; $\ 1 + \frac14 + \frac19 + \frac1{16} = \frac{205}{144}$.

\textbf{\color{brand}Step-by-step explanation.} $P_{\text{success}} = \sum_j|\beta_j|^2\frac{C^2}{\lambda_j^2} = \frac14\left(1 + \frac14 + \frac19 + \frac1{16}\right) = \frac14\cdot\frac{144 + 36 + 16 + 9}{144} = \frac{205}{576} = 0.3559028$.
It also equals $C^2\lVert\vec x\rVert^2/\lVert\vec b\rVert^2 = \frac{25^2 + 7^2 + 11^2 + 5^2}{48^2} = \frac{820}{2304} = \frac{205}{576}$. ($\frac{205}{144}$ forgets $|\beta_j|^2 = \frac14$;
$\frac{25}{48}$ is $x_0$.) Since every $\lambda_j$ is exactly representable, the state-vector run reproduces it to all printed digits.

\textbf{\color{brand}Why the other options are incorrect.}
\begin{itemize}[leftmargin=2.4em, itemsep=1pt, topsep=2pt]
\item[\textbf{(B)}] $\frac{25}{48}$ is the first solution component $x_0$.
\item[\textbf{(C)}] $\frac{205}{144} > 1$ forgets $|\beta_j|^2 = \frac14$ --- impossible for a probability.
\item[\textbf{(D)}] $\frac14 = |\beta_j|^2$ for one component; it ignores the $1/\lambda_j^2$ factors.
\end{itemize}
\textbf{\color{brand}Key takeaway.} $P_{\text{success}} = C^2\lVert\vec x\rVert^2/\lVert\vec b\rVert^2$: here $\frac{205}{576}$ exactly.
\end{solbox}
\qsep

\needspace{13\baselineskip}
\begin{qitem}{80}
\qmeta{Paragraph based}{Recorded run: $\lVert\vec x\rVert$ from $P_{\text{success}}$}
Without knowing $\vec x$ in advance, the norm $\lVert\vec x\rVert$ can be inferred from the success probability. Its value is
\par\smallskip\noindent\begin{tabularx}{\linewidth}{@{}*{2}{L}@{}}
\optlabel{A}\,$\dfrac{24}{\sqrt{205}} \approx 1.676$ & \optlabel{B}\,$\dfrac{205}{576} \approx 0.356$ \\[3pt]
\correctlabel{C}\,$\dfrac{\sqrt{205}}{24} \approx 0.597$ & \optlabel{D}\,$1$, because quantum states are normalised
\end{tabularx}
\end{qitem}
\begin{solbox}{C}
\textbf{\color{brand}Correct answer.} \textcolor{okgreen}{\textbf{(C)}} $\dfrac{\sqrt{205}}{24} \approx 0.597$

\textbf{\color{brand}Concepts required.} The success probability measures the norm of the post-selected branch, from which $\lVert\vec x\rVert$ follows.

\textbf{\color{brand}Relevant formulae.} $\lVert\vec x\rVert = \frac{\lVert\vec b\rVert}{C}\sqrt{P_{\text{success}}}$.

\textbf{\color{brand}Step-by-step explanation.} The accepted branch is $\vec a = \frac{C}{\lVert\vec b\rVert}\vec x$ (Q.53) and $\lVert\vec a\rVert^2 = P_{\text{success}}$, hence
\[ \lVert\vec x\rVert = \frac{\lVert\vec b\rVert}{C}\sqrt{P_{\text{success}}} = \sqrt{\frac{205}{576}} = \frac{\sqrt{205}}{24} \approx 0.5966 . \]
Check: $\frac1{48}\sqrt{820} = \frac{2\sqrt{205}}{48} = \frac{\sqrt{205}}{24}$. On hardware the measured success rate ($0.3543 \pm 0.0034$) gives $\lVert\vec x\rVert \approx 0.595 \pm 0.003$.
The normalised state $\ket x$ has norm $1$, but the solution vector does not.

\textbf{\color{brand}Why the other options are incorrect.}
\begin{itemize}[leftmargin=2.4em, itemsep=1pt, topsep=2pt]
\item[\textbf{(A)}] $\frac{24}{\sqrt{205}}$ inverts the relation.
\item[\textbf{(B)}] $\frac{205}{576}$ is $\lVert\vec x\rVert^2$; the square root is missing.
\item[\textbf{(D)}] $1$ is the norm of the normalised state $\ket x$, not of $\vec x$.
\end{itemize}
\textbf{\color{brand}Key takeaway.} Measure the success rate to recover the scale of $\vec x$ that the normalised state forgets.
\end{solbox}
\qsep

\end{document}